\documentclass[11pt]{amsart}
\pdfinfoomitdate=1
\pdftrailerid{}
\pdfsuppressptexinfo=15
\usepackage[T1]{fontenc}
\usepackage{lmodern,amsmath,amssymb,amsthm,mathtools,microtype}
\usepackage[margin=1in]{geometry}
\usepackage{array,longtable,booktabs}
\usepackage[hidelinks]{hyperref}
\emergencystretch=1em
\newtheorem{theorem}{Theorem}[section]
\newtheorem{proposition}[theorem]{Proposition}
\newtheorem{lemma}[theorem]{Lemma}
\newtheorem{corollary}[theorem]{Corollary}
\theoremstyle{definition}
\newtheorem{definition}[theorem]{Definition}
\newtheorem{remark}[theorem]{Remark}
\newcommand{\E}{\mathbb E}
\newcommand{\Pp}{\mathbb P}
\newcommand{\R}{\mathbb R}
\newcommand{\N}{\mathbb N}
\newcommand{\TV}{\mathrm{TV}}
\newcommand{\cE}{\mathfrak E}
\newcommand{\cF}{\mathfrak F}
\newcommand{\cH}{\mathcal H}
\newcommand{\diam}{\operatorname{diam}}
\newcommand{\Var}{\operatorname{Var}}
\newcommand{\supp}{\operatorname{supp}}
\newcommand{\law}{\operatorname{law}}
\newcommand{\cp}{\mathrm{cp}}
\newcommand{\on}{\mathrm{on}}
\newcommand{\ip}[2]{\langle #1,#2\rangle}
\newcommand{\norm}[1]{\lVert#1\rVert}
\title[Acquired geometry and causal resource transfer]{General Theta Foundations I:\\Acquired Geometry and Causal Resource Transfer}
\author{Qian Qi}
\date{October 7, 2026}
\subjclass[2020]{Primary 62B05, 60G35; Secondary 94A29, 28A80, 93E11}
\keywords{causal experiment, predictive quotient, acquired law, finite memory, singular quantization, experiment comparison}
\begin{document}
\begin{abstract}
A prepared causal experiment induces a predictive quotient and an actual acquired law. We prove a two-sided resource--resolution theorem in which old approximation error may pass between rare predictive strata without a reset. A forward squared-gain drift, coupled with raw occupation bounds, yields an implementable online quantizer with the same multiscale order as the checkpoint optimum. A finite-matrix criterion and invariant-reference comparison verify the hypotheses without inverse rare-state probabilities; exact suspension introduces no inverse activity factor. The geometric lower permits finitely overlapping, countable score neighborhoods. A recurrent non-erasing symbolic experiment verifies the theory on singular acquired sets with no limiting quantization exponent, while a nonlinear hidden-state filter supplies a different realization. A finite-acquisition variant has an explicit evolving predictive law and a coupled acquired-depth--memory--calibration minimax law for one terminal task. Resource-certified causal morphisms compose common-task errors while retaining acquisition, simulator, controller, workspace, program, precision and clock costs. Complementary finite-chart, reset-preserved, information-converse and minimax results are retained in the appendices.
\end{abstract}
\maketitle
\section{The experimental problem}\label{sec:experiment}
The mother problem is to pass from a physically executable experiment to a quantitative law for the predictions that can actually be retained. A formally reachable set need not carry appreciable acquired probability. A good codebook at a checkpoint need not admit a good recursive implementation. The first obstruction is probabilistic and the second causal; neither is measured by ambient dimension alone.

The difficulty becomes sharper when an experiment changes predictive strata. Old error need not disappear at the transition, and a bound normalized by a rare destination mass can become meaningless even when the physical update is stable. Our central result, Theorem~\ref{thm:nonreset}, controls the marked old-error kernel before such a normalization. A positive forward drift propagates every allocation-dependent approximation at once. Actual occupation bounds compare the forcing created by an update with the terminal acquisition. Exact holds are handled inside the energy recursion, so an arbitrarily long wait retains its physical cost but introduces no fictitious new rounding error.

Combined with global covering and arbitrary-center small-ball estimates, this mechanism gives matching checkpoint and online laws on finitely or countably many acquired components. A bounded overlap multiplicity replaces disjoint score neighborhoods. Proposition~\ref{prop:nr-matrix} derives the drift from raw cross-stratum kernels; Proposition~\ref{prop:nr-reference} derives occupation comparison from a reference invariant measure while keeping the actual prepared law nonstationary. Neither proposition assumes an optimum, a covariance, a Fisher matrix or reset at every component change.

Theorems~\ref{thm:nr-moran} and \ref{thm:hmm} verify the general mechanism for different raw experiments. The first carries old symbolic tails through recurrent expanding and contracting transitions between arbitrarily rare strata, without recurrent state-revealing resets. Its singular resolution profile can lack a limiting exponent. The second begins with a hidden-state transition and observation density, derives the nonlinear predictive filter and its actual acquired density, and verifies the same transfer hypotheses. In Theorem~\ref{thm:recurrent-acquisition}, only finitely many initial bits are acquired: the actual predictive law is computed from a recurrent known-depth chain. A single task then has a coupled acquisition--memory--calibration minimax law, with every initial reading and waiting call charged. This is a consequence of the common-oracle transfer principle, Proposition~\ref{prop:nr-acquisition}, rather than a replacement of the general theorem by a selected realization.

Resource comparison is made in a common task metric. Theorem~\ref{thm:morphism} supplies parameter-independent causal simulation with its strategy lift, internal state, clock and readout. Metric approximation errors add before squaring; complete interactive total-variation defects add in bounded loss. Persistent cardinalities multiply only for a stated product implementation. Executable continuation challenges give separate information lower bounds. Fixed-protocol matching laws are not unrestricted optimal-control or unknown-kernel learning theorems.

Predictive-state minimality \cite{shalizi,psr}, probability quantization \cite{graf}, singular Moran quantization \cite{zhu-moran}, quantized filtering \cite{pages,sellami,feuer}, belief-state control \cite{saldi}, average contraction and positive moment stability \cite{diaconis,costa,ogura}, and statistical experiment comparison \cite{blackwell,lecam,torgersen} supply established neighboring theories. We claim no novelty for conditional projection, a positive-matrix resolvent, or a cylinder coding. The contribution developed here is their experiment-level coupling through nonreset error transport, actual acquisition and an executable two-sided resource law. Section~\ref{sec:scope} distinguishes that assertion from the classical tools and records the remaining mathematical questions. Complementary finite anisotropic, reset-preserved, stationary-moment and converse arguments are retained in the appendices with their original hypotheses.

\subsection{Raw data, strategies, and preparation}
A discrete-time experiment consists of standard Borel spaces $X_t,\mathcal A_t,Y_{t+1}$, a probability preparation $\mu$ on $X_0$, and nonnegative normalized kernels
\begin{equation}\label{eq:raw}
 K_t(dx',dy,dc\mid x,a).
\end{equation}
The visible history $h_t$ records actions, reports, and the published nonnegative resource increments $c$. These increments include failed preparations and physical elapsed time. Positivity means nonnegativity and normalization, not a positive lower bound for every report. Deterministic transitions are permitted. An unknown physical parameter, when present, is included in the hidden state with a specified prior for the Bayesian assertions below. Parameter-uniform experiment comparison will be stated separately.

A legal strategy is a stochastic kernel $\alpha_t(da\mid h_t)$ supported on the declared available actions. All randomness used by the apparatus or algorithm is independent of the hidden preparation unless generated by \eqref{eq:raw}. A stopping rule is a stopping time for the visible filtration. The main risk theorem fixes a deterministic raw-acquisition horizon $N$ and an exploration strategy $\beta$ independently of the encoder. It does not optimize $\beta$. At most $N$ actual calls, including unsuccessful calls, are made. A stopped comparison is always explicitly bounded by $N$.

Successive kernel integration gives the path law $P^\alpha$. The complete conditional belief $\pi_t(h_t)$ on the hidden state has a measurable version. We specify versions on a declared admissible history domain; identities on impossible reports are not inferred by division by zero.

\subsection{Executable future tests and the quotient}
A future test is a legal finite continuation followed by a bounded measurable function of its actual reports. Hidden-state functions qualify only when a terminal measurement producing them is part of the experiment. Tests include determining report events and the continuation interface: legal actions, remaining budget, and every clock value affecting future experiments. They are closed under legal one-step continuation, finite randomization, and the bounded monotone limits needed to determine the designated laws.

For an admissible prepared belief $b$, let $\mathsf e_t(b)$ be the evaluations of a countable determining family of these tests, together with the continuation interface. Equality of $\mathsf e_t$ is \emph{predictive equivalence}. The family determines the entire designated continuation law, not merely the finite scoring menu introduced below. Here and throughout, a statement about an unavailable action is not a continuation assertion.

We impose the following realization certificate, verified from the instruments. The admissible belief domain $\mathcal B_t$ and the action domain are compact metric. The determining evaluations are continuous on $\mathcal B_t$. For a continuous report component there is a compact metric report domain $Y$, a finite dominating measure $\lambda$ with full support on $Y$, and an instrument density $T_{t,a,y}$. For every determining next-time test $f$, the functions
\begin{equation}\label{eq:density-pullback}
 n_f(b,a,y)=bT_{t,a,y}f,\qquad g(b,a,y)=bT_{t,a,y}1
\end{equation}
are jointly continuous. The functions $b\mapsto\int_E n_f(b,a,y)\lambda(dy)$, for report events $E$ in a determining algebra, are evaluations of executable one-step continuations, hence are constant on predictive fibres. The same is required for $f=1$. This is an \emph{event-instrument closure}, not the assertion that a report singleton has positive probability. The normalized belief update $U_{t,a,y}$ preserves the declared domain and is jointly continuous where $g>0$. If it is used at $g=0$, a continuous, predictive-fibre-compatible extension is supplied separately. Finite report components use the corresponding atomic instruments. Deterministic reveal components may instead be handled by an explicitly verified continuous quotient update, as in Section~\ref{sec:sensor}.

\begin{lemma}[Measurable factorization with completed histories]\label{lem:factor}
Let $R:\Omega\to E$ and $S:\Omega\to S_0$ be random variables, where $E$ is standard Borel and $S_0$ is a nonempty Borel subset of a countable product of compact intervals. If every coordinate of $S$ is measurable with respect to the completion of $\sigma(R)$ under $P$, there is a Borel map $f:E\to S_0$ such that $S=f(R)$ almost surely. It is unique $P_R$-almost everywhere among factors of $S$.
\end{lemma}
\begin{proof}
A completed-measurable real random variable has an uncompleted-measurable version: replace the finitely many measurable sets in each dyadic simple approximation by $\sigma(R)$-sets differing only by null sets, and take a pointwise limit on the common full-measure set. Every set in $\sigma(R)$ is $R^{-1}(B)$ for some Borel $B\subset E$, since these inverse images already form a sigma algebra. The same simple approximations therefore give Borel coordinate factors on $E$; assign zero wherever their limits fail to exist. Combine the countably many factors into a product-valued Borel map $f_0$. Since $S_0$ is Borel, $\{e:f_0(e)\in S_0\}$ is Borel and has full $P_R$-measure. Replace its complement by a fixed point of $S_0$. This proves existence. Equality of two factors after composition with $R$ is exactly their equality $P_R$-almost everywhere.
\end{proof}

\begin{proposition}[Predictive realization and pointwise update descent]\label{prop:quotient}
Under the preceding certificate, $S_t=\mathsf e_t(\mathcal B_t)$ is compact metrizable and realizes predictive equivalence. Report laws descend continuously. The normalized update descends continuously on the positive-density domain; a declared continuous fibre-compatible extension descends on its declared extended domain. In particular, equivalent beliefs have equivalent updated beliefs for \emph{every} positive-density report, not just almost every report. Any standard-Borel-valued history statistic through which the determining evaluations and continuation interface are completed-measurable admits a measurable predictive factor almost surely. Setwise sufficiency gives a unique setwise factor on the image, but no additional Borel conclusion for an arbitrary measurable quotient.
\end{proposition}
\begin{proof}
Rescale the determining evaluations to $[-1,1]$. Their continuous map into a countable product has compact metrizable image. The determining property identifies its fibres with equality of the designated continuation laws.

Fix equivalent $b,b'$ and an available action $a$. Event-instrument closure gives equal integrals of $n_f(b,a,\cdot)$ and $n_f(b',a,\cdot)$ over a determining algebra. Uniqueness of finite signed measures extends this equality to all Borel events, and equality of densities follows $\lambda$-almost everywhere. A continuous function vanishing almost everywhere under a full-support measure vanishes everywhere: any nonzero value would have a same-sign neighborhood of positive measure. Hence the two numerator functions, and the two denominators, agree for every $y$. Where $g>0$, dividing proves equality of every determining next-time evaluation. The continuation interface is among these evaluations. At a zero denominator only the separately verified extension is used.

A continuous surjection from a compact space onto a Hausdorff space is a quotient map. Its product with a compact action or report domain has the same property. Thus the continuous fibre-constant report evaluations and extended updates descend. Restriction to the saturated open positive-density domain also descends, since that domain is the inverse image of an open set under the descended continuous $g$. Apply Lemma~\ref{lem:factor} to the countable predictive coordinates for measurable minimality. Setwise minimality is the elementary inclusion of sufficient-statistic fibres in predictive fibres.
\end{proof}

This proposition fixes conditional versions rather than inferring them from null events. Finite exact interface components may be included as a disjoint compact union. A noncompact report or belief domain needs a separate Borel realization or a verified localization. No universal smooth realization of an arbitrary predictive equivalence relation is asserted.

\subsection{The common prediction task}
At time $N$, after the label has been retained, an independent index $i\in\{1,\ldots,q\}$ is selected with fixed positive weights $v_i$, $\sum_i v_i=1$. The corresponding legal terminal probe returns $T_i\in[0,1]$. The algorithm reports $a_i\in[0,1]$ and incurs $(T_i-a_i)^2$. Each expectation refers to its own executable probe; a joint counterfactual vector of all probe outcomes is unnecessary. For each possible terminal time $t$, write $T_{t,i}$ for its corresponding executable probe; at the selected horizon abbreviate $T_i=T_{N,i}$. Put
\begin{equation}\label{eq:score}
 p_t(s)=\big(\sqrt{v_i}\,\E[T_{t,i}\mid s]\big)_{i=1}^q,
 \qquad B_N=\sum_i v_i\E\Var(T_i\mid H_N),
 \qquad \nu_N=\law(p_N(s_N)).
\end{equation}
A varying terminal time uses the corresponding declared probe. The finite scoring menu need not determine the entire predictive quotient. Our assumptions below require only its Lipschitz upper bound globally and a separating lower bound on each acquired chart. Unmeasured interface coordinates cannot be silently charged to this task's lower bound.

A checkpoint encoder reads $H_N$ and retains at most $m$ labels. An online encoder has at most $m$ persistent labels at every time, updates from its previous label and the current input/report, and has no readable history tape. A terminal decoder may depend on the public horizon and on the freshly selected probe index, but not on discarded observations. For these fixed-exploration quantization statements, the encoder/decoder shared seed is independent of the entire prepared path and task law, including the exploration randomness. Shared randomness does not carry observation-dependent state. The risks $R_\cp(N,m)$ and $R_\on(N,m)$ are infima of the full expected score over these respective classes, under the fixed exploration $\beta$. Temporary workspace, numerical input access, and program description are recorded separately in Section~\ref{sec:transfer}.

The mathematical label class allows Borel transitions on exact current reports. Its risk is not the risk of a finite-bit processor. For an implementation with finite input precision and arithmetic, we write $R_{\rm alg}$ for its actual risk and retain the errors $u_t$ and $v_N^{\rm num}$ explicitly. The latter is a deterministic upper bound on the terminal score-norm error; a stated $L^2$ bound could be used instead. Appendix~\ref{sec:algorithms} supplies both types of algorithms. A fixed public clock may give the predetermined horizon or protocol phase, but cannot encode observation-dependent past information through adaptive timing. Any such timing channel is retained state and is charged.

\section{Acquired geometry under non-erasing causal transport}\label{sec:nonreset}
A transition between predictive strata need not destroy the information, or the approximation error, inherited from its source. The relevant stability condition must therefore act on the whole marked transition kernel. We give such a condition. It is a forward drift inequality for squared gains, combined with an occupation comparison for the \emph{actual} acquisition. Neither condition refers to an encoder optimum or to an already computed forced-error covariance.

We use $t=0,\ldots,T$ for the initialized recursion. Its paid preparation, subsequent calls and terminal test have a separately recorded raw count $N_{\rm raw}$. In particular, $T$ is not a convention under which an initialization or terminal probe is free. The notation $R_\cp(T,M)$ and $R_\on(T,M)$ denotes the risks of Section~\ref{sec:experiment} for this declared protocol.

\subsection{Geometry and the raw marked kernel}
Let $Z_t$ be a standard Borel factor of the fixed prepared experiment, with a measurable predictive coordinate $s(Z_t)$ in a bounded metric domain $(D,d)$. Any controller coordinates needed to obtain the following conditional law belong to $Z_t$. Conditional on the preceding history and independent algorithm randomness, the next step is an exact hold with probability $1-\vartheta_t$, or has a marked kernel $P_t(dz',d\omega\mid Z_{t-1})$ with probability $\vartheta_t$. The numbers $\vartheta_t\in[0,1]$ are deterministic known protocol data. On a hold both $Z_t$ and the retained representative are copied exactly. This is a statement about predictive and declared controller coordinates: the physical clock still advances and is charged. The time-dependent legal interface is known at the indexed time and belongs to the separate controller/clock ledger. This suspension assumption is not an assertion about state-dependent waiting or stopping. Write
\begin{equation}\label{eq:nr-law}
 \mu_t=(1-\vartheta_t)\mu_{t-1}+\vartheta_t\mu_{t-1}P_t,
 \qquad \widehat\mu_t=\mu_{t-1}P_t.
\end{equation}
Here and below a marked kernel acts on the state by integrating its marks. The conditional-kernel requirement includes independence of the next mark from the past approximation error given $Z_{t-1}$. Analysis-only coordinates are not made available to the algorithm. Encoder--decoder shared randomness is independent of the entire fixed experiment; any exploration randomness is part of the declared experiment and is not an additional decoder oracle.

Let $(D_j)_{j\in\mathcal J}$ be a finite or countable Borel partition of the factor states of full acquired mass, and put $w_{t,j}=\mu_t(D_j)$. The notation $s(D_j)$ means the predictive image, not the set of factor states itself. Fix an anchor $o\in D$ and a terminal score map $p_T:D\to\R^q$, with Lipschitz constant $L$. At $w_{T,j}>0$ let $\nu_{T,j}$ be the conditional acquired score law. The following geometric hypothesis permits overlaps of predictive images.

\smallskip\noindent\textbf{(G) Covering, actual small balls and finite multiplicity.}
There are constants $C_g,c_g,c_s>0$, $D_2\ge1$, an integer $\Lambda\ge1$, and positive nonincreasing scales $r_j(k)$, $k\in\N_0$, such that
\begin{equation}\label{eq:nr-regular}
 r_j(0)=r_j(1)=a_j,\quad \sup_j a_j<\infty,\quad
 r_j(k)^2\le D_2r_j(2k)^2\quad(k\ge1).
\end{equation}
For every $k\ge1$, at most $k$ valid representatives cover $s(D_j)$ at radius $C_g r_j(k)$; also $d(s(z),o)\le C_g a_j$ for $z\in D_j$. The sets $S_{T,j}=p_T(s(D_j))$ are compact. Every point of $\R^q$ belongs to at most $\Lambda$ of the open neighborhoods
\[
 U_{T,j}=\{x:\operatorname{dist}(x,S_{T,j})<c_sa_j\}.
\]
For all $k\ge1$, all $x\in\R^q$, and all positive-mass components,
\begin{equation}\label{eq:nr-smallball}
 \nu_{T,j}(B(x,c_g r_j(k)))\le\frac1{2k}.
\end{equation}
These are uniform certificates for the stated protocol class. No ambient volume is substituted for $\nu_{T,j}$. In particular, atomic components require a separate treatment, such as Theorem~\ref{thm:main}; they do not satisfy \eqref{eq:nr-smallball} at every $k$.

Define the acquired resolution profile
\begin{equation}\label{eq:nr-profile}
 \mathcal Q_T(m)=\inf_{\substack{k_j\in\N_0\\\sum_jk_j\le m}}
                       \sum_jw_{T,j}r_j(k_j)^2.
\end{equation}
Only finitely many $k_j$ are nonzero. An unallocated component is represented by the anchor. Effective selection of an allocation is a separate algorithmic requirement; the infimum in \eqref{eq:nr-profile} need not be attained.

\smallskip\noindent\textbf{(L) Executable local approximation.}
For every finite allocation there is a Borel recursive machine with at most one anchor and $\sum_jk_j$ other labels. It uses only its retained label, the current report and declared program data. If $e_t=d(s(Z_t),\widehat s_t)$, its active-step error satisfies
\begin{equation}\label{eq:nr-local}
 e_t\le g_t(Z_{t-1},Z_t,\omega)e_{t-1}+b_*R_{J_t},
 \qquad R_j=C_g r_j(k_j)+\xi_j,
\end{equation}
where $\xi_j\ge0$, $b_*<\infty$, and $J_t$ is the actual destination component. The exact-label case has $\xi_j=0$. The initialized error is at most $R_{J_0}$. Holds copy the label without new error. The gain $g_t\ge0$ is fixed by the raw update and is independent of the allocation. The rule must still be executable when several unallocated components have the same anchor; the analytic index $J_t$ is not a free retained register. No restriction requires $g_t=0$ on a component-changing edge.

\smallskip\noindent\textbf{(D) Forward squared-gain drift.}
A measurable function $v$ and constants $V<\infty$, $0\le\kappa<1$ satisfy
\begin{equation}\label{eq:nr-drift}
 1\le v\le V,\qquad
 \int g_t(z,z',\omega)^2v(z')P_t(dz',d\omega\mid z)
       \le\kappa v(z),\quad t=1,\ldots,T.
\end{equation}
The inequality is pointwise on the declared factor domain. It controls incoming old error through all destinations, including nonreset cross-stratum edges. Individual gains can exceed one.

\smallskip\noindent\textbf{(O) Actual occupation comparison.}
For a constant $A\ge1$, independent of the allocation,
\begin{equation}\label{eq:nr-occupation}
 \mu_0(D_j)\le A w_{T,j},\qquad
 \widehat\mu_t(D_j)\le A w_{T,j}\quad(t=1,\ldots,T;\ j\in\mathcal J).
\end{equation}
A zero terminal weight forces the corresponding seed and active-injection masses to vanish. This is a restriction on the actual preparation and kernels, not a conclusion inferred from a formal reachable set. Proposition~\ref{prop:nr-reference} verifies it without dividing by any component probability. The single-component case has $A=1$ with no invariant-law assumption.

\subsection{The central transfer theorem}
\begin{theorem}[Nonreset acquired-geometry and causal resource transfer]\label{thm:nonreset}
Fix a prepared positive causal experiment, a legal exploration independent of the encoder, a deterministic initialized horizon $T$, and a common finite squared-prediction task. Realize the predictive quotient by Proposition~\ref{prop:quotient}, or by a separately proved Borel realization with its legal interface and pointwise update versions. Suppose \textup{(G)}, \textup{(L)}, \textup{(D)} and \textup{(O)} hold for the marked factor above. Set
\begin{equation}\label{eq:nr-H}
 H_*=VA\max\left\{1,\frac{b_*^2}{(1-\sqrt\kappa)^2}\right\}.
\end{equation}
There are $c,C>0$, depending only on $C_g,c_g,c_s,D_2,\Lambda,L$, such that for every $M\ge2$,
\begin{equation}\label{eq:nr-main}
 B_T+c\mathcal Q_T(M)\le R_\cp(T,M)\le R_\on(T,M)
                 \le B_T+C H_*\mathcal Q_T(M).
\end{equation}
The upper is an infimum over explicitly admissible recursive label machines, not over checkpoint maps. For any allocation with $\sum k_j\le M-1$, finite local errors $\xi_j$, and score readout error at most $\zeta_T$, the same construction has
\begin{equation}\label{eq:nr-finite}
 R_{\rm alg}-B_T\le
 \left[L\sqrt{H_*}\left(C_g\sqrt{\sum_jw_{T,j}r_j(k_j)^2}
                  +\sqrt{\sum_jw_{T,j}\xi_j^2}\right)+\zeta_T\right]^2.
\end{equation}
Approximate minimizing allocations give \eqref{eq:nr-main}, or its finite-error counterpart, with arbitrarily small additive allocation error.

The constants have no dependence on $\min_j w_{T,j}$, the number of components, or $\min_t\vartheta_t$. The dependence on $A,V,\kappa$ is precisely that displayed in $H_*$. Exact Borel labels do not imply an effective finite-bit implementation. When the local routines and a causal morphism are effective, their acquisition, simulator-state, controller, clock, workspace, program and precision costs are the full ledger of Theorem~\ref{thm:morphism}. Complete causal total-variation defects enter the common bounded-loss risk additively, whereas the metric errors in \eqref{eq:nr-finite} add before squaring.
\end{theorem}

We prove the geometric and dynamic parts separately. This also isolates which conditions are responsible for each side of the matching law.

\begin{lemma}[Finite-overlap small-ball lower]\label{lem:nr-overlap}
Under \textup{(G)}, the quadratic quantization error of the actual score law is bounded below by $c\mathcal Q_T(M)$ for every $M\ge1$. The lower allows arbitrary decoder centers and independent shared randomness. Moreover,
\begin{equation}\label{eq:nr-coarsen}
 \mathcal Q_T(M-1)\le D_2\mathcal Q_T(M)\quad(M\ge1),\qquad
 \mathcal Q_T(\Lambda M)\ge D_2^{-\lceil\log_2(2\Lambda)\rceil}\mathcal Q_T(M).
\end{equation}
\end{lemma}
\begin{proof}
For centers $x_1,\ldots,x_M$, assign center $x_l$ to every $j$ with $x_l\in U_{T,j}$, and let $k_j$ be the assigned counts. Finite multiplicity gives $\sum_jk_j\le\Lambda M$. If $k_j=0$, every center is at distance at least $c_sa_j$ from every score in $S_{T,j}$. If $k_j>0$, put $r=r_j(k_j)$ and $c_0=\min\{c_g,c_s/2\}$. Centers not assigned to $j$ have distance at least $c_sa_j\ge2c_0r$. The union of the $k_j$ assigned balls of radius $c_0r$ has conditional mass at most $1/2$. Consequently the conditional distortion is at least $c_0^2r_j(k_j)^2/2$, with a possibly smaller uniform constant also covering $k_j=0$. Sum against the actual weights to obtain $c\mathcal Q_T(\Lambda M)$.

For the second inequality in \eqref{eq:nr-coarsen}, take any allocation $k$ at budget $\Lambda M$ and put $l_j=\lfloor k_j/\Lambda\rfloor$. Then $\sum l_j\le M$. If $l_j=0<k_j$, use $r_j(0)=r_j(1)$ and $k_j<\Lambda$; if $l_j\ge1$, use $k_j<\Lambda(l_j+1)\le2\Lambda l_j$. Iterating \eqref{eq:nr-regular} and monotonicity in either case gives $r_j(l_j)^2\le D_2^{\lceil\log_2(2\Lambda)\rceil}r_j(k_j)^2$. The case $k_j=0$ is immediate. Taking approximate infima proves the assertion. For the first inequality, remove a label from an occupied entry. A singleton entry has no cost change. For $k\ge2$, $k\le2(k-1)$ implies $r_j(k-1)^2\le D_2r_j(k)^2$. If no label needs removing, the same inequality is automatic.

For each fixed independent random seed a decoder has at most $M$ centers; random output can be replaced by its conditional mean. Pointwise nearest-center selection is an admissible improvement of a checkpoint encoder. The deterministic lower therefore survives averaging over the seed. Conditional projection, Lemma~\ref{lem:projection}, identifies this quantization error with excess prediction risk.
\end{proof}

\begin{lemma}[Cross-stratum error transport with exact suspension]\label{lem:nr-energy}
Under \textup{(L)}, \textup{(D)} and \textup{(O)}, for every fixed allocation and every nonnegative allowance vector with $\sum_jw_{T,j}R_j^2<\infty$,
\begin{equation}\label{eq:nr-energy}
 \E e_t^2\le\E[v(Z_t)e_t^2]\le H_*\sum_jw_{T,j}R_j^2
                    \quad(0\le t\le T).
\end{equation}
\end{lemma}
\begin{proof}
Put $q^2=VA\sum_jw_{T,j}R_j^2$ and $a_t^2=\E[v(Z_t)e_t^2]$. The seed gives $a_0\le q$. Conditional on an active step, the old-error term in \eqref{eq:nr-local} has weighted $L^2$ norm at most $\sqrt\kappa a_{t-1}$. Indeed the conditional-kernel property makes its squared integral
\[
 \E\left[e_{t-1}^2\int g_t(Z_{t-1},z',\omega)^2v(z')
                              P_t(dz',d\omega\mid Z_{t-1})\right],
\]
which is at most $\kappa a_{t-1}^2$. The forcing term has squared weighted norm at most
\[
 b_*^2V\sum_j\widehat\mu_t(D_j)R_j^2\le b_*^2q^2.
\]
Minkowski's inequality is used \emph{inside the active branch}. On the exact hold branch the energy is unchanged. The deterministic suspension coin therefore gives
\begin{equation}\label{eq:nr-mixture}
 a_t^2\le(1-\vartheta_t)a_{t-1}^2+
                  \vartheta_t(\sqrt\kappa a_{t-1}+b_*q)^2.
\end{equation}
The interval $0\le a\le q\max\{1,b_* /(1-\sqrt\kappa)\}$ is invariant under this inequality, including $\vartheta_t=0$. Induction proves \eqref{eq:nr-energy}. When $q=0$, the seed and forcing vanish and the same induction gives zero energy. For a countable partition the forcing integral is the sum of nonnegative component integrals by monotone convergence. No independence between current error and the component is used, and no error arriving from another component is discarded. In particular the proof does not divide an unconditional bound by a rare terminal mass. It bounds the allocation-weighted global risk directly.
\end{proof}

\begin{proof}[Proof of Theorem~\ref{thm:nonreset}]
The lower follows from Lemma~\ref{lem:nr-overlap} and conditional projection. Every online encoder is a checkpoint encoder, giving the middle inequality. For any allocation at budget $M-1$, use the machine in \textup{(L)}. Lemma~\ref{lem:nr-energy} and the Lipschitz task map bound its excess by $L^2H_*C_g^2\sum_jw_{T,j}r_j(k_j)^2$. Taking approximate infima and using \eqref{eq:nr-coarsen} proves the exact upper. This argument remains valid when an infimum is zero, by additive approximation rather than an undefined relative approximation.

For finite precision apply Minkowski first to the weighted sequence $(C_gr_j(k_j)+\xi_j)_j$, and then to the score error and terminal readout error. This proves \eqref{eq:nr-finite}. A restricted finite-input class is contained in the exact-input checkpoint class, so the geometric lower remains a lower for it. Effectivity and total resource feasibility are separately required by \textup{(L)} and the stated ledger. The common-loss defect assertion is the independently proved composition in Theorem~\ref{thm:morphism}; it neither changes the oracle baseline nor supplies an unproved converse for a precision or simulator register.
\end{proof}

\subsection{Raw criteria, without rare-mass normalization}
\begin{proposition}[Invariant reference and nonstationary acquisition]\label{prop:nr-reference}
Suppose every active raw factor kernel preserves a probability $\lambda$, and the actual prepared law satisfies $c\lambda\le\mu_0\le C\lambda$ as measures, for $0<c\le C<\infty$. Then \textup{(O)} holds with $A=C/c$ for every Borel partition, horizon and deterministic suspension schedule. The actual measures need not be stationary and are still those in \eqref{eq:nr-law}.
\end{proposition}
\begin{proof}
Positivity and $\lambda P_t=\lambda$ imply $c\lambda\le\mu P_t\le C\lambda$ whenever $c\lambda\le\mu\le C\lambda$. Convex combination with the identity preserves these inequalities. Induction gives them for every $\mu_t$ and $\widehat\mu_t$. Hence $\widehat\mu_t(D_j)\le C\lambda(D_j)\le(C/c)\mu_T(D_j)$; the seed is identical. Null components need no division. The measure $\lambda$ is only a comparison device: actual terminal weights and conditional small-ball laws in \textup{(G)} must still be calculated or bounded from $\mu_T$.
\end{proof}

\begin{proposition}[A raw matrix criterion allowing nonreset edges]\label{prop:nr-matrix}
Let the factor be a finite disjoint union of fibers, with reference laws $\sigma_i$. At an active step, a raw transition selects a destination with probabilities $p_{ij}$, and an edge kernel $T_{ij}$ sends $\sigma_i$ to $\sigma_j$. Marks and predictive transformations are part of $T_{ij}$. Suppose the local gains obey $g\le L_{ij}$ on this edge. More generally it is enough that $\int g(z,z',\omega)^2T_{ij}(dz',d\omega\mid z)\le L_{ij}^2$ for every source state $z$; in this form the same edge may carry both contracting and expanding marks. Put
\[
 C_2=(p_{ij}L_{ij}^2)_{i,j}.
\]
If $\rho(C_2)<1$, there are $v_i\ge1$ and $\kappa<1$ such that \textup{(D)} holds with $v(z)=v_i$ on fiber $i$. For a stationary probability vector $\pi$ of $(p_{ij})$, $\lambda=\sum_i\pi_i\sigma_i$ is invariant. A preparation bounded above and below relative to $\lambda$ then satisfies \textup{(O)}, with constants not involving $\min\{\pi_i:\pi_i>0\}$. The geometric partition in \textup{(G)} may refine the fibers and may be countable.

For a time-dependent family of edge kernels the same assertions hold when the kernels have a common invariant $\lambda$ and a common positive vector with $C_{2,t}v\le\kappa v$. For a single finite nonnegative matrix, existence of such a positive vector is equivalent to $\rho(C_2)<1$.
\end{proposition}
\begin{proof}
When $\rho(C_2)<1$, the finite matrix series $v=\sum_{n\ge0}C_2^n\mathbf1$ converges and satisfies $v\ge\mathbf1$ and $C_2v=v-\mathbf1\le(1-1/\max_i v_i)v$. Integrating the pointwise bounds, or the uniform conditional second-moment bounds, over all destinations gives \eqref{eq:nr-drift}. Conversely a positive $v$ with $C_2v\le\kappa v$ makes the induced weighted maximum norm of $C_2$ at most $\kappa$, so its spectral radius is less than one. The invariance calculation is
\[
 \sum_i\pi_i\sum_jp_{ij}\sigma_iT_{ij}
        =\sum_j\left(\sum_i\pi_ip_{ij}\right)\sigma_j
        =\sum_j\pi_j\sigma_j.
\]
Proposition~\ref{prop:nr-reference} gives the occupation comparison for any geometric partition. The time-dependent assertion uses the same calculations at each time; no joint-spectral-radius necessity is asserted. This criterion is a characterization of the stated positive \emph{gain comparison matrix}, not a necessity theorem for optimal finite-state prediction.
\end{proof}

The forward matrix in Proposition~\ref{prop:nr-matrix} is generally not the reverse conditional-moment operator of Theorem~\ref{thm:kernel}. Here the same positive drift controls old error for \emph{every allocation}; occupation comparison then controls all destination-dependent forcing radii at once. There is no normalization of a previously integrated covariance by $\operatorname{diag}(w_T)^{-1/2}$. Exact suspension is dealt with by \eqref{eq:nr-mixture}, rather than by applying Minkowski across hold and active branches, which would introduce a spurious small-activity loss.

\begin{corollary}[A common-oracle law and finite-resource composition]\label{cor:nr-resources}
Under the exact hypotheses of Theorem~\ref{thm:nonreset}, let $b_\circ\le B_T$ be a fixed common ideal baseline. For a class on which the displayed constants are uniform,
\begin{equation}\label{eq:nr-common}
 R_\on(T,M)-b_\circ\asymp
             (B_T-b_\circ)+\mathcal Q_T(M).
\end{equation}
If a certified product implementation uses $R$ simulator, $D_c$ controller and $Q$ phase states, a total persistent budget $M_{\rm tot}$ permits the upper with $M=\lfloor M_{\rm tot}/(RD_cQ)\rfloor\ge2$. Calibration and arithmetic allowances $\xi_j\le u_j\delta+z_j\eta$ give the right side of \eqref{eq:nr-finite} with
\[
 \sqrt{\sum_jw_{T,j}\xi_j^2}
 \le\delta\sqrt{\sum_jw_{T,j}u_j^2}+\eta\sqrt{\sum_jw_{T,j}z_j^2}.
\]
A parameter-independent causal simulator of complete common-task defect $\varepsilon$ adds at most $\varepsilon$ to a loss bounded by one, after the explicit strategy and task lift.
\end{corollary}
\begin{proof}
Add $B_T-b_\circ\ge0$ to both sides of the two-sided excess law and absorb uniform positive constants. The product construction is the finite-state composition in Theorem~\ref{thm:morphism}; feasibility of its other resource coordinates remains required. The precision estimate is the triangle inequality in weighted $\ell^2$. The final assertion is total-variation contraction followed by the bounded-loss inequality on the complete lifted task law. None of these operations gives a product-state lower or a universal $\delta^2$ or $\varepsilon$ lower. Such lower bounds require indistinguishable alternatives or executable information challenges, as in Theorems~\ref{thm:robust}, \ref{thm:soft-state} and \ref{thm:single-probe}.
\end{proof}

\subsection{Acquisition error and an atomic predictive law}
An actual acquired law can be atomic before sufficient observations have been made, so its small balls need not obey \textup{(G)} at arbitrarily fine scales. The next consequence handles this situation without replacing the actual predictive law by a latent-state law.

\begin{proposition}[Common-oracle acquisition transfer]\label{prop:nr-acquisition}
Let $Y\in\R^q$ be the common bounded squared-task vector, with task weights included in its norm, and let $\mathcal G$ be a declared ideal information sigma-field containing the actual history sigma-field $\sigma(H)$. Set $X=\E[Y\mid\mathcal G]$ and $b_\circ=\E\|Y-X\|^2$. The actual conditional prediction is $P=\E[X\mid H]$, with law $\nu$, and its acquisition error is $\mathcal A=\E\|X-P\|^2$. Orthogonal projection gives actual Bayes risk $b_\circ+\mathcal A$. Suppose, at every budget $M$ considered, an executable online machine has excess above that actual Bayes risk at most $C_ur(M)^2$, and the \emph{ideal-score} law has arbitrary-center $M$-point distortion at least $c_lr(M)^2$, with $c_l>0$. The online upper may in particular be supplied by Lemma~\ref{lem:nr-energy}. Then
\begin{equation}\label{eq:nr-acquisition-transfer}
 R_\cp(M)-b_\circ=\mathcal A+q_M(\nu),\qquad
 R_\on(M)-b_\circ\asymp\mathcal A+q_M(\nu)
                      \asymp\mathcal A+r(M)^2,
\end{equation}
where $q_M$ is quadratic probability quantization and the constants depend only on $C_u,c_l$. No small-ball assumption on the possibly atomic actual law $\nu$ is required. The ideal law is only a lower witness; the middle expression uses the actual acquired law and cannot in general omit $\mathcal A$.
\end{proposition}
\begin{proof}
Orthogonal conditional projection gives
\[
 \E\|X-d(S)\|^2=\mathcal A+\E\|P-d(S)\|^2
\]
for every randomized encoder $S$ based on $H$, with independent randomization. Checkpoint nearest-center selection and conditional-mean decoding therefore give the equality of infima in \eqref{eq:nr-acquisition-transfer}. Every such decoder is also an $M$-center approximation to $X$ when all of $X$ is available, so its error is at least $q_M(\law(X))\ge c_lr(M)^2$. It is also at least $\mathcal A$. Thus the checkpoint excess is at least $\tfrac12\min\{1,c_l\}(\mathcal A+r(M)^2)$. The assumed online construction gives the reverse bound $\mathcal A+C_ur(M)^2$. Nesting checkpoint and online classes proves all comparisons. This uses neither an identification of $X$ with the actual prediction $P$ nor a claim that the ideal oracle is implementable at the actual acquisition budget.
\end{proof}

\section{Causal morphisms and the full resource account}\label{sec:transfer}
We now make explicit the implementation structure through which a geometric risk theorem is transferred. In particular, an internal register is not made free by omitting it from a risk argument, and a register appearing in one construction is not automatically mandatory in every construction.

\subsection{Algorithms, resources and admissible interfaces}
Fix a parameter set $\Lambda$, raw experiments $\mathfrak E_\lambda$, a common task with loss in $[0,1]$, and a class of legal controllers. A controller has no access to $\lambda$. A resource-bounded algorithm consists of Borel state-transition kernels, a terminal readout, and, in the finite-bit model, a fixed finite program and a specified numerical input interface. There is no readable discarded-report tape. A random tape is independent of the experiment and immutable; an observation-dependent pointer into it is retained state.

We use the implementation ledger
\begin{equation}\label{eq:full-budget}
 b=(N_{\rm raw},m,R,D_c,Q,W,P,b_{\rm cal},b_{\rm num},\tau).
\end{equation}
Here $N_{\rm raw}$ counts all raw calls, including pilots, failures and every actually executed terminal probe; $m,R,D_c,Q$ are, respectively, representative, simulator, controller and internal-clock state capacities; $W$ is transient workspace in bits; $P$ is program and readonly-data length in bits; the two precision capacities specify the calibration and numerical input interfaces; arithmetic guard bits are charged in workspace; and $\tau$ is physical elapsed time. Actual numerical error is a separate property of those interfaces and routines, not automatically $2^{-b_{\rm num}}$. Workspace is reset between calls. Any observation-dependent information surviving a reset belongs to the persistent account. A table and a generator program have different $P$ and time accounts.

The total persistent capacity of a direct product implementation is $M=mRD_cQ$. To discuss arbitrary fused implementations, we also use
\begin{equation}\label{eq:total-budget}
 \mathbf b=(N_{\rm raw},M,W,P,b_{\rm cal},b_{\rm num},\tau).
\end{equation}
A fused machine is constrained by its total number of observation-dependent states, not by a preferred factorization. Both vectors remain explicit in the relevant risk notation. In Sections~\ref{sec:experiment}--\ref{sec:kernel}, $N$ denotes calls before the forecast; embedding a protocol with one further physical terminal probe requires $N_{\rm raw}\ge N+1$. A menu of counterfactual probes is never charged as though they had all been executed. The resource preorder is coordinatewise comparison of capacities; known interface precision requirements are compared in the same direction. Different input interfaces are different algorithm models, even at the same integer capacity.

The superscript $\chi$ in $\mathcal A^\chi_{\mathfrak E}(b)$ and $\mathcal R^\chi_{\mathfrak E}(b)$ specifies the cost semantics. Unless indicated otherwise, $\chi={\rm pw}$ means pathwise call and time bounds. For $\chi={\rm mean}$ these two bounds hold in expectation uniformly over the declared parameter/controller class. A tail specification includes its failure probabilities explicitly. Persistent capacities, workspace and readonly-description bounds are pathwise in all three semantics. For example, the parameter-uniform risk is
\[
 \mathcal R^\chi_{\mathfrak E}(b)
   =\inf_{A\in\mathcal A^\chi_{\mathfrak E}(b)}
       \sup_{\lambda\in\Lambda}\E^{A}_{\mathfrak E_\lambda}\ell.
\]
The fixed-preparation version replaces the supremum by the specified prepared law. The empty-class value is $+\infty$. For the fixed-exploration prediction theorem the algorithm chooses the encoder, not the exploration.

A public clock is free only when it is the fixed exogenous protocol clock, independent of observed values. An adaptively delayed emission, a data-dependent stopping phase that will be read later, or a remembered failure count is not such a clock. Its observation-dependent states enter $Q$, or enter the total capacity $M$ directly.

\subsection{Complete causal-morphism data}
A morphism $\Gamma:\mathfrak E\to\mathfrak F$ has the following data. Source and target have specified visible filtrations and legal action correspondences. The simulator has an initial state kernel, independent of the unknown parameter, and a microstep kernel. Given its current state, the current target command and the source prefix, this kernel either chooses a legal source action or emits a target report and a new simulator state. Successive target emissions occur at nondecreasing source stopping indices $\sigma_k$; the declared finite call bound prevents infinite internal execution. A terminal command is mapped to a legal source continuation producing the common loss-relevant terminal readout. A target stopping decision is mapped to source stopping, and the published target clocks and costs have specified report transformations. The emitted reports include failures when the target interface includes them. Prefix dependence is a mathematical description, not permission to reread a discarded tape: the resource certificate must implement it through the declared retained state and current inputs.

The resulting target-to-source lift $\mathsf L_\Gamma$ is part of the data. It maps every admitted target controller into an admitted source controller, preserving action availability, task semantics and the stopping convention. This is a statement to be checked on the declared classes, not merely a type name. The state update and report transformation are measurable on standard Borel domains. All auxiliary randomness is parameter independent. In finite-bit assertions the same data have the executable routines and numerical interfaces charged in \eqref{eq:full-budget}.

The complete interactive defect is
\begin{equation}\label{eq:defect}
 \epsilon_\Gamma(b)=\sup_{\lambda,\,\alpha\in\mathcal A_{\mathfrak F}(b)}
 \left\|\law^{\mathfrak E_\lambda}(\Gamma;\mathsf L_\Gamma\alpha)
             -\law^{\mathfrak F_\lambda}(\alpha)\right\|_{\TV}.
\end{equation}
The compared laws include the target-visible transcript, costs and clocks, the common terminal readout and the controller's terminal decision. They also include any independent controller seed shared with a later readout. Equivalently, it suffices to have a uniform comparison conditional on each such seed. This prevents the unjustified extension of a marginal-transcript bound to a correlated controller decision. Our convention is $\|P-Q\|_{\TV}=\sup_A|P(A)-Q(A)|$, one half of the density $L^1$ distance. Taking the infimum over a declared class of these simulators defines causal deficiency. A particular simulator supplies only an upper certificate.

A resource map $\rho_\Gamma$ is a monotone map of the full ledgers with a proof that the lifted implementation lies within its image budget. One useful explicit pathwise map is
\begin{align}\label{eq:resource-map}
 \rho_\Gamma(b)=\big(&n_0+n_\Gamma N_{\rm raw},\ m,\ r_\Gamma R,\ d_\Gamma D_c,
 q_\Gamma Q,\ W+w_\Gamma(b),\ P+p_\Gamma(b),\notag\\
 &\beta_{{\rm cal},\Gamma}(b),\ \beta_{{\rm num},\Gamma}(b),\
 \tau_0+t_\Gamma N_{\rm raw}+t'_{\Gamma}\tau\big).
\end{align}
The nonnegative coefficients and monotone functions are known protocol bounds. Pilots are included in $n_0,\tau_0$. Precision maps specify enough source precision for the target accuracy, not a change of the unknown physical parameter. The workspace sum is a safe simultaneous-storage bound; a proof of reuse may improve it. A direct product machine at the image budget has at most $m(r_\Gamma R)(d_\Gamma D_c)(q_\Gamma Q)$ persistent states. A different verified monotone cost map is permitted; merely writing $\rho$ does not certify one.

\begin{theorem}[Typed resource and common-risk composition]\label{thm:morphism}
Let $\Gamma_1:\mathfrak E\to\mathfrak F$ and $\Gamma_2:\mathfrak F\to\mathfrak G$ have the preceding data, compatible action/readout/clock interfaces, and resource maps $\rho_1,\rho_2$. For each admitted outer budget $b$, assume
\[
 \mathsf L_{\Gamma_2}\mathcal A_{\mathfrak G}(b)
       \subseteq\mathcal A_{\mathfrak F}(\rho_2(b)),
 \quad
 \mathsf L_{\Gamma_1}\mathcal A_{\mathfrak F}(\rho_2(b))
       \subseteq\mathcal A_{\mathfrak E}(\rho_1(\rho_2(b))).
\]
Their serial composition is a causal, parameter-independent morphism with resource map $\rho_1\circ\rho_2$ and defect at most
\[
 \min\{1,\epsilon_1(\rho_2(b))+\epsilon_2(b)\}.
\]
For a common bounded loss and the same preparation or parameter supremum,
\begin{equation}\label{eq:risktransfer}
 \mathcal R^{\rm pw}_{\mathfrak E}(\rho_\Gamma(b))
       \le\mathcal R^{\rm pw}_{\mathfrak F}(b)+\epsilon_\Gamma(b).
\end{equation}
For the prediction implementation of Theorem~\ref{thm:main}, with $N_{\rm raw}\ge N+1$, keeping every coordinate of $b$ and its image,
\begin{equation}\label{eq:joint}
 \mathcal R^{\rm pw}_{\mathfrak E}(\rho_\Gamma(b))
 \le B_N+\left(C\sqrt{\Gamma_N\Psi_N(m)}
       +L\|U_N\|_2+v_N^{\rm num}\right)^2
       +\epsilon_\Gamma(b)+\Delta_{\rm rep}+p_{\rm bad}.
\end{equation}
Here the fixed-preparation risk is used. With parameter-uniform certificates the parameter-uniform version takes the supremum of the entire right side over the parameter; in particular the Bayes baseline is not assumed parameter independent. The term $\Delta_{\rm rep}$ is the bounded-horizon allowance \eqref{eq:pathallowance} only for a separate report-simulation stage not already counted in $\epsilon_\Gamma$. Deterministic approximation certificates may fail on an event of probability at most $p_{\rm bad}$.

Expected-cost versions hold when the inner per-call and pilot bounds hold conditionally at every calling history, uniformly over the lifted class. Tail versions hold with the corresponding joint good-resource event and with the union bound on its declared failure probabilities. Unconditional expected-cost bounds alone do not establish the asserted affine composition.
\end{theorem}
\begin{proof}
Execute the two microstep kernels in their causal order, carrying their joint internal state. Their stopping indices compose because every decision uses only the currently available prefix; the finite resource maps exclude infinite internal execution. The two inclusions in the statement establish legality of the composed lift. State updates, report readouts, clocks and terminal commands compose as measurable kernels. This constructs the morphism and its full resource map, including precision, workspace and readonly data. In particular independent implementation registers multiply their state capacities, while their ideal logarithmic bit capacities add.

Fix a parameter and an outer controller, including its shared independent seed. Insert between the two path laws the law with an exact intermediate experiment. The first discrepancy is at most $\epsilon_1(\rho_2(b))$ by the lifted-class inclusion and the complete interactive defect. The second nonanticipating transformation is a Markov kernel and cannot increase total variation. The remaining discrepancy is at most $\epsilon_2(b)$. The triangle inequality proves the stated defect. The inclusion of the common terminal readout ensures that this is a comparison of the actual task, not of two surrogate outputs.

For $0\le\ell\le1$, the layer-cake formula bounds the difference of its expectations by total variation. Apply this to each target algorithm, then take the same parameter supremum and the algorithm infimum. This gives \eqref{eq:risktransfer}, using arbitrarily near-optimal target algorithms when an infimum is not attained. Substitute Theorem~\ref{thm:main} to obtain \eqref{eq:joint}. On a bad-certificate event the loss is at most one; on its complement the pointwise approximation estimate is integrated against the original law. This argument does not condition the geometric law on the good event.

For mean budgets, write the total inner cost as the pilot cost plus the sum of per-call costs times the indicator that the call is reached. Conditional bounds and the tower property give the affine expectation bound, including random numbers of calls. For tail budgets, intersect the certified good events, compose the pathwise bounds there, and use the union bound for their complement. These arguments also explain the stated restrictions on expected and tail semantics.
\end{proof}

Calibration and arithmetic errors add in the one-step $u_t$ before multiplication by the same transport factors. Minkowski's inequality then adds their propagated root-mean-square contributions before squaring. A simulation defect acts on the common bounded loss and is added linearly. An allowed calibration error or defect need not occur on any particular instance; a converse requires an indistinguishability experiment. Nor does a forward simulator transfer a lower bound in the reverse direction. The next section supplies lower experiments and a mandatory-state theorem rather than making either inference.

\section{A singular recurrent experiment without state-revealing resets}\label{sec:nr-moran}
The next realization transports an unknown tail between predictive strata. Expanding transitions recur, and neither cross-stratum edge overwrites the old state. Its invariant measure is singular, but its actual prepared acquisition is allowed to be nonstationary. The mechanism is thus different from both uniform filter contraction and refresh-dominated error erasure.

\subsection{The raw experiment}
Fix a sequence $\rho_n\in[1/5,1/3]$, put $\ell_0=1$, $\ell_n=\prod_{h=1}^n\rho_h$, and define
\begin{equation}\label{eq:nr-code}
 x(b)=\sum_{n\ge1}(1-\rho_n)\ell_{n-1}b_n,
 \qquad b\in\{0,1\}^{\N}.
\end{equation}
Let $K$ be the image of this coding and let $\sigma$ be the image of the independent fair-bit law. If two sequences first differ at position $n$, summing the tails gives
\begin{equation}\label{eq:nr-gap}
 \frac{\ell_{n-1}}3\le |x(b)-x(c)|\le\ell_{n-1}.
\end{equation}
Consequently the coding is continuous and one-to-one, and has a continuous inverse on $K$. A cylinder of length $n$ has mass $2^{-n}$ and diameter at most $\ell_n$.

The physical state is $(I,b)\in\{0,1\}\times\{0,1\}^{\N}$. Take $\alpha\in(0,1]$, $\beta=1/2$, and let an active command use the mode transition matrix
\begin{equation}\label{eq:nr-mode}
 P=\begin{pmatrix}1-\alpha&\alpha\\[1mm]1/2&1/2\end{pmatrix},
 \qquad \pi=\left(\frac{1/2}{\alpha+1/2},\frac{\alpha}{\alpha+1/2}\right).
\end{equation}
On $0\to0$, independently choose deletion of the first old bit with probability $\gamma=1/4096$, and otherwise prepend two newly sampled fair bits. On $0\to1$ delete the first old bit. On $1\to0$ prepend four newly sampled fair bits. On $1\to1$ copy the sequence. The report gives the source and destination modes, the operation, and the newly sampled bits, if any. It never reveals the remaining old tail. All these rules are positive normalized Borel kernels. An independent suspension at time $t$ copies the state and reports a hold with probability $1-\vartheta_t$; otherwise the active kernel is used.

Let $\lambda$ be the law with modes $\pi$ and independent fair tails. The actual preparation is any known law $\mu_0$ with
\begin{equation}\label{eq:nr-preparation}
 c\lambda\le\mu_0\le C\lambda,\qquad 0<c\le C<\infty.
\end{equation}
A paid initialization reports the prepared mode and sequence. This specifies the exact Borel experiment. A finite-access encoder reads only a stated finite prefix during initialization, at a separately charged input and time cost; it cannot subsequently reread the initialization report. This distinction is essential. Replacing the full initialization by a truncated physical report defines a different experiment and is not asserted to preserve a complete-transcript deficiency or its full-history baseline.

There is one terminal command: a paid Bernoulli probe of mean
\begin{equation}\label{eq:nr-score}
 p(0,b)=\frac18+\frac{x(b)}8,\qquad
 p(1,b)=\frac58+\frac{x(b)}8.
\end{equation}
The probe ends the experiment. Active calls, holds, initialization and this terminal call are all counted. For $T$ post-initialization opportunities the unit-call implementation uses $N_{\rm raw}=T+2$ raw calls. Input-prefix access and arithmetic time are additional ledger entries, not extra information retained for free. The legal strategy for the theorem is this fixed sequence of $T$ opportunities followed by the terminal command. All finite continuations using these commands are the executable tests defining the quotient.

Define
\begin{equation}\label{eq:nr-moran-profile}
 r(k)=\ell_{\lfloor\log_2 k\rfloor}\quad(k\ge1),\qquad r(0)=1.
\end{equation}
Equip the predictive domain $\{0,1\}\times K$ with distance $|x-y|$ within a mode and distance one between modes. This is a bounded metric; take $(0,x(000\cdots))$ as anchor.

\begin{theorem}[Non-erasing recurrent singular realization]\label{thm:nr-moran}
For the raw experiment above the predictive quotient is $(I,x(b))$, with its declared terminal interface. For every $T\ge0$ and $M\ge4$,
\begin{equation}\label{eq:nr-moran-law}
 R_\cp(T,M)-B_T\asymp R_\on(T,M)-B_T
          \asymp \ell_{\lfloor\log_2 M\rfloor}^{\,2},
 \qquad B_T=\E[p(I_T,b_T)(1-p(I_T,b_T))].
\end{equation}
The constants depend only on $c,C$ and the displayed universal coding bounds. They are uniform in $T$, $\alpha\in(0,1]$, every deterministic suspension schedule, and every sequence $(\rho_n)$ in $[1/5,1/3]$. In particular they remain bounded when mode one becomes arbitrarily rare and when all activity probabilities tend to zero. The gain drift can be taken with
\begin{equation}\label{eq:nr-explicit-drift}
 v_0=450,\quad v_1=1,\quad V=450,\quad\kappa=17/20,
 \qquad A=C/c.
\end{equation}
Both cross-mode transitions preserve a nonconstant dependence on the old infinite tail. For $0<\alpha<1$, arbitrarily long runs of the expanding within-mode deletion have positive probability. The theorem therefore does not reduce to a uniform deterministic block-contraction assumption and does not use recurrent state-revealing refreshes.

An exact Borel implementation uses at most $M$ labels and no additional retained mode flag. With effective finite-prefix input, computable ratios and certified terminal arithmetic, it is a finite-input implementation with the same label count and with explicitly nonzero readout error as in \eqref{eq:nr-finite}. For ratios in $\{1/3,1/5\}$, initialization need read at most $\lceil\log_2M\rceil$ tail bits; the subsequent report has at most four new bits and finite mode data. The input-access, program, temporary-space and time costs are given below the proof.
\end{theorem}
\begin{proof}
After initialization the full visible history determines the mode and tail by the reported transformations. Future kernel laws are consequently functions of $(I,b)$. Conversely the terminal mean in \eqref{eq:nr-score} identifies the mode, since its two intervals are disjoint, and then identifies $x$ and hence $b$ by \eqref{eq:nr-gap}. Thus two different such states are separated by an executable terminal test. This proves sufficiency and minimality of the stated Borel quotient. The update maps are defined on every valid state for its compatible reports. Incompatible source-mode reports have probability zero; a declared extension applies the reported sequence operation to the supplied predictive coordinate and sets the reported destination mode. It is never justified by null-event Bayes division.

Each active edge sends the fair product law to itself: deletion of the first fair bit leaves a fair independent tail, and prepending independent fair bits has the same property. Thus the edge laws send $\sigma$ to $\sigma$, and \eqref{eq:nr-mode} gives $\lambda P=\lambda$. Proposition~\ref{prop:nr-reference} proves \textup{(O)} and $c\lambda\le\mu_t\le C\lambda$, including the active candidate laws. In particular the actual conditional tail law in each mode is bounded above by $(C/c)\sigma$ and below by $(c/C)\sigma$. Its actual weight is $w_{t,i}=\mu_t(\{i\}\times K)$, not an assigned replacement $\pi_i$.

Deleting the first bit has Lipschitz constant at most $15$ on $K$. Indeed, for a first difference at $n\ge2$, use the upper bound $\ell_{n-2}$ for the shifted difference and the lower bound $\ell_{n-1}/3$ for the original difference. Their ratio is at most $3/\rho_{n-1}\le15$. For $n=1$, the original difference is at least $1/3$ and the shifted difference at most one. Prepending $m$ common bits has Lipschitz constant at most $3^{1-m}$: the ratio of the two bounds in \eqref{eq:nr-gap} is at most $3\ell_{n+m-1}/\ell_{n-1}\le3^{1-m}$. The prepend-two, deletion, prepend-four and copy gains may be chosen as $1/3,15,1/27,1$. On the $0\to0$ edge their conditional squared-gain bound is $(1-\gamma)/9+225\gamma=85/512$. This is an average over its raw reported marks, uniformly in the source state, not an average with respect to the old error. No such bound says deletion is a contraction.

The forward gain matrix is
\begin{equation}\label{eq:nr-C2}
 C_2=\begin{pmatrix}(1-\alpha)85/512&225\alpha\\[1mm]1/1458&1/2\end{pmatrix}.
\end{equation}
With $v=(450,1)$ its first component is $(19125/256)(1-\alpha)+225\alpha\le225<(17/20)450$, while its second is $450/1458+1/2=131/162<17/20$. This proves \textup{(D)} uniformly even though an expanding nonreset edge has squared gain $225$. The argument uses no factor $1/\pi_1$. If $\alpha>0$, cross-mode passages recur almost surely along infinitely many active opportunities, since the two-state active chain is irreducible; an infinite run of physical holds is not counted as activity. The finite-horizon theorem does not require an infinite-activity assumption. For $0<\alpha<1$, a run of $n$ within-mode deletions has conditional probability $((1-\alpha)\gamma)^n>0$. The $n$-fold deletion separates the zero tail and a tail with its first one at position $n+1$ by a distance ratio at least $(5/6)3^n$. Hence arbitrary expanding stretches genuinely occur in the raw map family, even though its forced second moment is stable.

For the cover, use one zero-tail representative for each cylinder of length $n=\lfloor\log_2 k\rfloor$. There are $2^n\le k$ such representatives and their covering radius is at most $\ell_n$. Also $r(k)^2/r(2k)^2=\rho_{n+1}^{-2}\le25$. To verify actual small balls uniformly in $k$, put $K_0=C/c$ and choose a fixed integer $h\ge1$ with $2^h\ge4K_0$. Distinct cylinders of length $n+h$ have mutual gap at least $\ell_{n+h-1}/3\ge\ell_n/(3\cdot5^{h-1})$. Hence a Euclidean ball of radius $\ell_n/(12\cdot5^{h-1})$ meets at most one such cylinder. Its actual conditional mass is at most $K_0 2^{-(n+h)}\le2^{-n-2}\le1/(2k)$, since $k<2^{n+1}$. The score shrinks within-mode distances by $1/8$, so reduce the ball constant by another factor eight. This establishes \eqref{eq:nr-smallball} for arbitrary centers, including centers outside the Cantor set. The two score sets have a fixed positive gap; neighborhoods of radius $1/16$ are disjoint. Condition \textup{(G)} holds with $\Lambda=1$, $a_0=a_1=1$ and constants depending only on $K_0$.

We give the local machine, since a geometric cover alone is insufficient. For an allocation $(k_0,k_1)$, an allocated mode $i$ uses all binary words of length $n_i=\lfloor\log_2 k_i\rfloor$ as labels, each decoded as its zero-tail extension in that mode. Unallocated modes use the common anchor. On initialization select the corresponding prefix label or the anchor. On a hold copy the label. On an active report, apply the reported deletion, prepend or copy operation to the finite word with its implicit zero tail. If the destination is allocated, truncate or pad to its prefix depth and retain that label; otherwise retain the anchor. The report supplies the destination during the update, not as a surviving side channel. At the forecast cut the decoder sees only the retained label.

When source and retained modes agree, the preceding Lipschitz estimates and final cylinder truncation give \eqref{eq:nr-local} with $b_*=1$, $C_g=1$. If an unallocated source has been encoded by the other-mode anchor, its metric error is one. The coordinate discrepancy after deletion is at most one, after prepending $m$ common bits at most $\ell_m\le3^{-m}\le3^{1-m}$, and after copying at most one. These are bounded by the declared gain times that metric error. Final rounding adds $r(k_i)$. If the destination is unallocated, its metric error is at most one, already bounded by $R_i=r(0)=1$. This checks the local rule even when the analytic component was discarded. Initialization error is at most the corresponding radius. The terminal score is globally one-Lipschitz in the declared metric. Active report laws depend only on the mode, so their total variation is zero within a mode and at most the cross-mode distance one; the terminal Bernoulli kernel has the same continuity bound. Every hypothesis of Theorem~\ref{thm:nonreset} is now derived from the raw experiment.

Finally the mixture profile has the stated uniform scale. Since $k_i\le M$, every term $r(k_i)^2$ is at least $r(M)^2$, giving $\mathcal Q_T(M)\ge r(M)^2$. Conversely allocate $\lfloor M/2\rfloor$ to each mode for $M\ge4$. The depth loss is at most two, whence $\mathcal Q_T(M)\le25^2r(M)^2$, independently of the two actual weights. Apply \eqref{eq:nr-main}. This also covers $T=0$, when the initialized quantizer is the online machine and there are no updates.
\end{proof}

\paragraph{Finite-access and bit costs.}
The initial encoder need inspect at most $n_{\max}=\max_i n_i\le\lceil\log_2M\rceil$ bits, during that one paid initialization, plus a mode bit. It may erase each inspected bit after accumulating the finite prefix label. Later updates operate only on the retained word and at most four fresh reported bits. Padding means appending zeros from the program, not consulting the discarded input. A straightforward implementation uses $O(n_{\max}+b_{\rm out})$ temporary bits, including the current report and arithmetic, and at most $\lceil\log_2M\rceil$ persistent bits for its label. Temporary bits are erased before the next cut. For rational ratios $1/3$ or $1/5$, a zero-tail representative is a rational whose numerator and denominator have $O(n_{\max})$ bits. Schoolbook integer arithmetic gives a conservative $O((n_{\max}+b_{\rm out})^3)$ bit-operation bound for its terminal evaluation to absolute error $2^{-b_{\rm out}}$; word updates cost $O(n_{\max})$. The program includes the effective ratio schedule, depths, transition rules and calibrated mode probabilities. The schedule's description length is not included in the label count. Sampling and transmitting bits, all suspended calls, and the terminal probe retain their physical-time charges. These are constructive costs, not optimal time--space converses. No full-transcript TV comparison between a continuous raw report and a finite packet is used.

\begin{corollary}[A nonreset singular law with no limiting exponent]\label{cor:nr-oscillation}
There is a computable choice $\rho_n\in\{1/3,1/5\}$ in Theorem~\ref{thm:nr-moran} for which the logarithmic resolution slope
\[
 \frac{-\log(R_\on(T,M)-B_T)}{\log M}
\]
has subsequences with limits $2\log 3/\log2$ and $2\log5/\log2$, uniformly along any choices of $T$, $\alpha$ and suspension schedule with fixed preparation constants. In particular no single power-law dimension describes this acquired resolution profile.
\end{corollary}
\begin{proof}
Take integers $n_j=2^{2^j}$, and make $\rho_n$ alternate between $1/3$ and $1/5$ on the successive blocks $(n_{j-1},n_j]$, with any finite initial block. Then $n_{j-1}/n_j\to0$. At $M_j=2^{n_j}$,
\[
 \frac{-\log\ell_{n_j}^2}{\log M_j}
  =\frac{2\sum_{h\le n_j}\log(1/\rho_h)}{n_j\log2}
\]
converges to the indicated two values on alternating block endpoints. Uniform multiplicative constants in \eqref{eq:nr-moran-law} contribute $O(1)/\log M_j$, which vanishes. These are positive decay exponents, not the signed slope of $\log$ risk. The argument proves a continuum statement through exact cylinder geometry, independently of any finite regression.
\end{proof}

\subsection{A second, structurally different realization of the same theorem}
\begin{corollary}[The nonlinear filtering class under the general transfer]\label{cor:nr-filter}
The controlled binary hidden-state experiments in Theorem~\ref{thm:hmm}, with flip probability $q_0\in(0,1/2)$ and actions in $[a_-,a_+]\subset(0,1)$, satisfy the dynamic hypotheses of Theorem~\ref{thm:nonreset} with a single geometric component, $v=1$, $A=1$ and $\kappa=(1-2q_0)^2$. Their actual acquired small-ball and covering profiles are comparable to $r(k)=k^{-1}$, with $r(0)=r(1)$ after fixed scaling. Consequently their checkpoint and online excess risks have the same $M^{-2}$ order uniformly over positive horizons and the specified encoder-independent control rules. This conclusion does not require stationary preparation.
\end{corollary}
\begin{proof}
From the raw density \eqref{eq:hmm-kernel}, normalization gives \eqref{eq:hmm-update}; the terminal hidden-bit probe identifies the scalar posterior quotient. The log-odds derivative computation in the proof of Theorem~\ref{thm:hmm} gives a pointwise same-report gain at most $1-2q_0$, so \eqref{eq:nr-drift} holds for the factor enlarged by any needed controller coordinates. A one-component partition makes \eqref{eq:nr-occupation} an equality. The positive lower bound on the raw report-to-posterior derivative \eqref{eq:hmm-jacobian}, together with the report-density upper bound, gives a uniform upper bound on the \emph{actual} posterior density at every positive horizon. Since the invariant interval has a fixed finite length, an equally spaced valid grid supplies the covering radius, and the density upper gives the arbitrary-center small-ball inequality after choosing $c_g$ small enough. Log-odds and Euclidean distance are equivalent on that interval. For the single score set $\Lambda=1$, and the dyadic profile bound is $D_2=4$. The recursive posterior update followed by grid rounding verifies \textup{(L)}, including initialization. Report-kernel continuity and the effective arithmetic interface are the explicit computations in Theorem~\ref{thm:hmm} and Appendix~\ref{sec:algorithms}. These computations prove the hypotheses from the raw kernel; the filtering realization is not a premise of the general transfer proof.
\end{proof}

\section{Joint acquisition, calibration and memory in one recurrent experiment}\label{sec:recurrent-acquisition}
We now replace the infinite initialization report by genuinely finite paid observations. This produces an atomic predictive law whose acquisition depth evolves through nonreset transitions. Its risk law depends on acquisition and memory through the same random depth, rather than through an independently selected terminal menu.

Use the physical coding \eqref{eq:nr-code}, modes \eqref{eq:nr-mode}, and stationary physical preparation $\lambda$. A preparation call reports the mode but no tail bits. The next $B\ge0$ calls reveal, one at a time and without altering the state, the first $B$ bits of that prepared tail. The observer may not revisit any discarded bit. There follow $T$ opportunities with the same independent suspension schedule as before. The active operations are now: on $0\to0$, delete one old bit with probability $\gamma=1/4096$ and otherwise prepend two fair reported bits; delete one old bit on $0\to1$; prepend \emph{six} fair reported bits on $1\to0$; and copy on $1\to1$. This six-bit choice supplies a uniform gain bound on the larger domain of partially known tails. No transition reveals the old unknown tail. Each call, including a hold, costs one acquisition. The terminal call ends the experiment, so
\begin{equation}\label{eq:ra-budget}
 N_{\rm raw}=B+T+2.
\end{equation}
Current-report packet bits and their reading time are also charged; a six-bit report is not a free infinite input.

Let $W_t$ be the prefix known to the \emph{full-history} observer after the $t$th opportunity, and $H_t=|W_t|$. Given that history, all remaining bits are independent and fair. The actual known depth has $H_0=B$ and evolves by
\begin{equation}\label{eq:ra-height}
 H_t=\begin{cases}
 H_{t-1}+2,&0\to0\text{ active prepend},\\
 \max\{H_{t-1}-1,0\},&\text{active deletion on }0\to0\text{ or }0\to1,\\
 H_{t-1}+6,&1\to0\text{ active},\\
 H_{t-1},&1\to1\text{ active or a suspension}.
 \end{cases}
\end{equation}
All weights below are obtained from this finite-time raw mode--height chain, including failed activity opportunities. Write $\omega_{T,i,h}=\Pp(I_T=i,H_T=h)$, and $\omega_{T,h}=\sum_i\omega_{T,i,h}$.

For a finite word $w$ of length $h$, let
\[
 m(w)=x(w000\cdots)+\tfrac12\ell_h.
\]
It is the mean of the unknown fair tail in its cylinder, not a revealed point of the original Cantor set. Put $\pi_1=\alpha/(\alpha+1/2)$ and define a nominal terminal Bernoulli mean
\begin{equation}\label{eq:ra-score}
 p(i,x)=\tfrac12+\tfrac14(i-\pi_1)+\tfrac18(x-\tfrac12).
\end{equation}
Its physical expectation under $\lambda$ is $1/2$. An unknown calibration parameter $\theta\in\{-\delta,\delta\}$, $0\le\delta\le1/16$, changes the terminal mean to $p(i,x)+\theta$ and has no effect on any preparation or acquisition report. The displayed ranges lie strictly inside $[0,1]$. Programs know the nominal kernel and the bound $\delta$, not the sign of $\theta$.

The nominal acquired predictive law is explicitly
\begin{equation}\label{eq:ra-law}
 \nu_{B,T}=\sum_{i=0}^1\sum_{h=0}^{B+6T}\omega_{T,i,h}\,2^{-h}
        \sum_{w\in\{0,1\}^h}\delta_{p(i,m(w))}.
\end{equation}
Thus this law is actually acquired and is finite at every finite raw budget. No non-atomic reference law is inserted into its definition. Put
\begin{equation}\label{eq:ra-profile}
 \Phi_{B,T}(M)=\sum_{h=0}^{B+6T}\omega_{T,h}
                \ell_{\min\{h,\lfloor\log_2M\rfloor\}}^2.
\end{equation}
The same acquired depth weights govern the unresolved variance and the memory resolution. Define the common ideal baseline
\begin{equation}\label{eq:ra-baseline}
 b_\circ=\E_\lambda[p(I,x)(1-p(I,x))]-\delta^2.
\end{equation}
It equals the oracle Bernoulli risk for both calibration signs, since $\E_\lambda p=1/2$. Let $\mathcal R_\cp$ and $\mathcal R_\on$ be the infimum over parameter-independent checkpoint and online $M$-label encoders of the worst risk over these two signs, minus this same $b_\circ$.

\begin{theorem}[A recurrent acquired-depth resource law]\label{thm:recurrent-acquisition}
For every $B,T\ge0$, $M\ge6$, $\alpha\in(0,1]$, every deterministic suspension schedule, every $\rho_n\in[1/5,1/3]$, and $0\le\delta\le1/16$,
\begin{equation}\label{eq:ra-main}
 \mathcal R_\cp(B,T,M,\delta)\asymp
 \mathcal R_\on(B,T,M,\delta)\asymp
                     \delta^2+\Phi_{B,T}(M),
\end{equation}
with universal positive constants. More precisely,
\begin{equation}\label{eq:ra-exact-cp}
 \mathcal R_\cp=\delta^2+\mathcal A_{B,T}+q_M(\nu_{B,T}),\qquad
 \frac1{576}\E\ell_{H_T}^2\le\mathcal A_{B,T}
                              \le\frac1{256}\E\ell_{H_T}^2,
\end{equation}
where $\mathcal A_{B,T}=\E\Var(p(I_T,x_T)\mid\text{actual history})$.

An online implementation retains a mode and a prefix of variable length at most $m$, encoded together in $2^{m+2}-2\le M$ labels. The length and mode are included in this count. It reads every paid input at most once. Its report and task continuity are explicit below; rational schedules in $\{1/3,1/5\}$ and rational or certified computable $\alpha$ admit finite-input arithmetic. Unbiased dyadic terminal rounding of mesh $2^{-b}$ adds at most $2^{-2b}/4$ to the risk. A further certified complete causal defect $\varepsilon$ adds at most $\varepsilon$ for this bounded loss. These last two are upper allowances, not universal error floors.
\end{theorem}
\begin{proof}
We first derive the predictive state from the raw observations. At the end of the $B$ read calls the observed word is uniform among words of length $B$ and its remaining physical tail is independent fair. A reported prepend concatenates independent known bits. A deletion removes the first known bit if the known word is nonempty; otherwise it removes a fair unobserved bit and leaves an independent fair tail. Copying preserves the description. Induction proves \eqref{eq:ra-height} and the claimed conditional-tail law. Since mode choices are independent of bit values, conditionally on $(I_T,H_T)=(i,h)$ the known word is uniform, giving \eqref{eq:ra-law}.

Let $K^*$ be $K$ together with the centers $m(w)$ of all finite cylinders. This is a compact set. Different centers are different: a cylinder center lies in the open gap between its two children, whereas a proper descendant center lies in a child. No center belongs to $K$. The nominal prediction coordinate $(i,m(w))$ therefore identifies both mode and finite known word; the nominal terminal mean in \eqref{eq:ra-score} identifies this coordinate because the two mode intervals have a gap $1/8$. All future reports and terminal laws are functions of it, so it is the predictive quotient for the nominal experiment. With the symmetric preparation on the unobserved calibration sign the same quotient applies; for the two-point minimax problem the whole terminal family is determined by the same coordinate. Neither construction allows the encoder to read the sign.

Conditional on a word of length $h$, independence of remaining fair bits gives
\[
 \Var(x_T\mid W_T)=\frac14\sum_{n>h}(1-\rho_n)^2\ell_{n-1}^2.
\]
Its first term is at least $\ell_h^2/9$, and its upper bound is $\ell_h^2/4$, since the tail has range $\ell_h$. The factor $1/8$ in the score yields the acquisition bounds in \eqref{eq:ra-exact-cp}.

Next consider any parameter-independent encoder and forecast $a$. Its excess over \eqref{eq:ra-baseline} at sign $\theta$ is $\E(p+\theta-a)^2$. The average of the two signs is $\E(p-a)^2+\delta^2$, giving a lower bound on their maximum. For any fixed encoder the nominal conditional-mean decoder $a=\E[p\mid S]$ is unbiased, and attains that average for both signs. Conditional projection of $p$ onto the actual full history and then checkpoint quantization proves the exact infimum \eqref{eq:ra-exact-cp}, including randomized encoders with independent seeds.

We construct the online upper rather than transferring a checkpoint map. Choose the largest nonnegative $m$ with $2^{m+2}-2\le M$. Store the current mode and a word of length at most $m$. During the $B$ initial reads retain the first $\min\{B,m\}$ bits and discard later bits immediately. A prepend adds the reported bits and truncates at depth $m$; a deletion pops one bit if present; a copy preserves the word. All these operations are functions of the finite label and current report. A missing old tail is never reread or replaced by new information. The label count is
\[
 2\sum_{h=0}^m2^h=2^{m+2}-2.
\]
Decode the stored word by its center, not by an alleged exact old state.

For completeness we verify the dynamic comparison on $K^*$. If two coordinates share an initial prefix and first differ at the next level, their distance is at least one sixth of the parent-cylinder scale and at most that scale. The factor $1/6$ covers the case where one coordinate is the parent center; two opposite children have the stronger gap $1/3$. Consequently deletion, extended to fix the empty-word center, is $30$-Lipschitz on $K^*$; prepending $l$ common bits is at most $6\cdot3^{-l}$-Lipschitz. The proof is the same first-difference comparison as \eqref{eq:nr-gap}, now with the center-to-child gap. Center coordinates with at least one common bit use the ratio of adjacent scales, at most five; the case of no common bit is bounded directly by the diameter. Thus the prepend-two, deletion, prepend-six and copy gains are $2/3,30,2/243,1$. The $0\to0$ edge has uniform conditional squared-gain bound $(1-\gamma)4/9+900\gamma=85/128$; arbitrarily long deletion runs are not excluded.

Take $v=(1800,1)$. The forward gain matrix satisfies
\[
 C_2v=\left(\frac{19125}{16}(1-\alpha)+900\alpha,\ \frac{3600}{59049}+\frac12\right)
                  \le\frac{17}{20}(1800,1).
\]
Truncation of any longer word to its depth-$m$ center changes its coordinate by at most $\ell_m$, and it leaves a shorter word exact. Initial compression has the same bound. Apply the energy proof of Lemma~\ref{lem:nr-energy} to the error between the full-history predictive coordinate and the stored center, with one forcing radius $\ell_m$, $A=1$, $V=1800$ and $\kappa=17/20$. Its argument uses no small-ball condition; the full-history factor here includes the actual known word, and the next marked report is independent of old approximation error given its mode. It follows that the nominal online excess above the actual Bayes risk is at most a universal constant times $\ell_m^2$. The copied suspension branch creates no new error.

The stored word is the actual physical prefix at its retained depth. Its length and its bit origins depend only on the initial read count and mode/operation history, not on the bit values. Conditional on that history and the retained prefix, all bit origins beyond the prefix remain independent fair. Averaging over histories therefore shows that the stated center decoder is exactly $\E[p\mid S]$. It is unbiased and hence adds precisely $\delta^2$ for the calibration ambiguity, with no unproved cancellation in a biased arithmetic rule. The mode-aware gain vector is an analytic proof weight, not another retained register.

For the lower, the physical tail law remains fair on every edge even though only a prefix is observed. The physical mode--tail law is therefore $\lambda$ at time $T$. The arbitrary-center cylinder lower in the proof of Theorem~\ref{thm:nr-moran}, with the separated affine score intervals of \eqref{eq:ra-score}, gives
\[
 q_M(\law_\lambda(p))\ge c_0\ell_{\lfloor\log_2M\rfloor}^2
\]
uniformly in the rare mode probability. This is a lower witness from the joint raw physical law, not the acquired law in \eqref{eq:ra-law}. Proposition~\ref{prop:nr-acquisition} combines it with the proved online upper and the actual acquisition variance; it shows why the witness cannot remove the latter. For $M\ge6$ the difference between $m$ and $\lfloor\log_2M\rfloor$ is at most two, so $\ell_m^2\le25^2\ell_{\lfloor\log_2M\rfloor}^2$. Finally, for $n=\lfloor\log_2M\rfloor$, monotonicity gives
\[
 \max\{\E\ell_{H_T}^2,\ell_n^2\}
 \le\E\ell_{\min(H_T,n)}^2
 \le\E\ell_{H_T}^2+\ell_n^2.
\]
These comparisons prove \eqref{eq:ra-main} from \eqref{eq:ra-exact-cp}. They also show that $\Phi_{B,T}$ is equivalent to the actual-law quantity $\mathcal A_{B,T}+q_M(\nu_{B,T})$, not merely to a formal latent geometric dimension.

Current operation reports depend on the mode and not on the unknown tail. With the metric of within-mode coordinate distance and unit cross-mode distance, their TV modulus is at most one times the distance. The terminal Bernoulli modulus within a mode is $|x-y|/8$, and is at most the cross-mode distance otherwise. The coordinate updates just proved provide the finite-input stability. To round a computed exact nominal center to adjacent dyadic numbers, choose the upper endpoint with probability equal to its relative fractional position. Conditional rounding mean is the center and variance is at most $2^{-2b}/4$. It uses fresh independent randomness, erased after the readout; cross terms with both calibration and nominal prediction error vanish by conditional expectation. A computable implementation of these probabilities with additional arithmetic error uses the explicit allowance in \eqref{eq:nr-finite}. The bounded complete-law TV allowance follows from Theorem~\ref{thm:morphism}. Neither upper bound claims a matching numerical or deficiency lower for every parameter choice.
\end{proof}

For rational ratio schedules and rational or certified computable mode probabilities, the center arithmetic and finite-word routines have the temporary-space and conservative bit-time costs stated in Section~\ref{sec:nr-moran}, with six current input bits rather than four. The encoder's persistent alphabet already includes its word length and mode. An external controller that counts the $B$ reading calls and $T$ opportunities must be charged a phase/clock implementation if such scheduling is not supplied by the declared apparatus. The displayed law is for that fixed paid protocol; it does not assert that this division of $N_{\rm raw}$ between reading and recurrence is an optimal adaptive exploration policy, or that an unknown transition kernel can be learned without pilots.

An exact rational Bernoulli randomization for the unbiased readout can be generated by rejection from independent fair bits; its expected random-bit cost is finite, but its worst-case running time is not bounded by that fact. A truncated sampler must charge its tail event as an additional task defect. The deterministic arithmetic bounds above apply to evaluating the center and rounding probabilities, not to an uncharged worst-case exact random sampler.

\section{Raw-kernel realization theorems}\label{sec:realizations}
\subsection{A controlled binary hidden-state filter}\label{sec:filter}
Let $X_t\in\{0,1\}$ have transition probability $q_0\in(0,1/2)$ of flipping at each raw call. The preparation is $P(X_0=1)=1/2$. A command selects $a$ from a fixed finite subset of $[a_-,a_+]\subset(0,1)$. Conditional on the new hidden state, the report $y\in[-1,1]$ has density
\begin{equation}\label{eq:hmm-kernel}
 g(y\mid X_{t+1}=x,a)=\tfrac12(1+a(2x-1)y).
\end{equation}
Each call costs one acquisition and one unit of physical time. The exploration may be any specified causal command rule, fixed independently of the encoder. A terminal probe reads $X_N$ and ends the experiment. Further filtering reports and this terminal probe are the designated future tests.

\begin{theorem}[Filtering realization]\label{thm:hmm}
For this raw-kernel class, the predictive quotient is the scalar conditional probability $p_t=P(X_t=1\mid H_t)$, with the terminal/continuation interface. It has a common compact invariant interval after initialization. For all $N\ge1$ and $m\ge1$,
\begin{equation}\label{eq:hmm-law}
 R_\cp(N,m)-B_N\asymp R_\on(N,m)-B_N\asymp m^{-2},
 \qquad B_N=\E[p_N(1-p_N)].
\end{equation}
The constants depend only on $q_0,a_-,a_+$, not on $N$ or on the specified command rule. The zero-error upper is an exact Borel finite-label implementation. With rational or certified finite-precision calibration, the separate program in Appendix~\ref{sec:algorithms} has the following nonzero-error guarantee. A uniform update/numerical error $u$ in log-odds distance contributes at most $u/(2q_0)$ to the propagated state error. The constants in this statement are not asserted uniform as $q_0\downarrow0$.
\end{theorem}
\begin{proof}
Direct normalization of \eqref{eq:hmm-kernel} gives, with $r=q_0+(1-2q_0)p$,
\begin{equation}\label{eq:hmm-update}
 f_a(y\mid p)=\tfrac12[1+ay(2r-1)],\qquad
 F_{a,y}(p)=\frac{r(1+ay)}{1+ay(2r-1)}.
\end{equation}
All denominators are at least $1-a_+$. Conditional laws of all future report words and the terminal probe are functions of $p$ by induction. Conversely the terminal probe has mean $p$, so different $p$ are predictively different. This proves quotient identification and minimality without assuming a covariance or information matrix.

Let
\[
 c_- =\frac{q_0(1-a_+)}{1+a_+(1-2q_0)},\qquad c_+=1-c_-.
\]
The interval $D=[c_-,c_+]$ contains $1/2$ and is invariant, because the prediction $r$ always lies in $[q_0,1-q_0]$ and the update is increasing in $r$ and $ay$. Use the metric
$d(p,p')=|\operatorname{logit}(p)-\operatorname{logit}(p')|$.
On $D$ it is equivalent to Euclidean distance, with constants depending only on the displayed parameters. The map $p\mapsto p$ is the task map, and its inverse-coordinate Lipschitz constants follow from the bounded derivative $1/[p(1-p)]$. This provides a one-dimensional global chart and an effective valid domain.

For update stability, put $v=p(1-p)$ and $\rho=1-2q_0$. The derivative of the predicted log-odds with respect to previous log-odds equals
\[
 \frac{\rho v}{q_0(1-q_0)+\rho^2 v}\le\rho.
\]
Indeed $v\le1/4$ and $1-\rho^2=4q_0(1-q_0)$, so the denominator is at least $v$. The observation then adds $\log[(1+ay)/(1-ay)]$, independently of the previous state. Thus every same-report update is $\rho$-Lipschitz in $d$. Proposition~\ref{prop:mechanisms}(a) gives $\Gamma_N\le(2q_0)^{-2}$, including initialization rounding.

It remains to verify actual acquired mass. For each preceding history and command,
\begin{equation}\label{eq:hmm-jacobian}
 \frac{\partial F_{a,y}(p)}{\partial y}
 =\frac{2ar(1-r)}{[1+ay(2r-1)]^2}
 \ge c_y:=\frac{2a_-q_0(1-q_0)}{(1+a_+)^2}>0.
\end{equation}
The conditional density of $y$ is at most $(1+a_+)/2$. Change of variables therefore gives a conditional density of $p_{t+1}$ at most $(1+a_+)/(2c_y)$ on its conditional interval and zero elsewhere. Mixtures over the actual histories and commands preserve this upper bound. Its integral is one. Consequently \eqref{eq:mass} holds on the fixed interval $D$ at every $N\ge1$, with $C_g$ this bound times $|D|$. No uniform posterior law has been inserted. Equations~\eqref{eq:chart} and \eqref{eq:mass} give $\Psi_N(m)=|D|^2/m^2$.

For report continuity, direct integration of the density difference gives
\[
 \norm{f_a(\cdot\mid p)-f_a(\cdot\mid p')}_\TV
 =\frac{a\rho}{2}|p-p'|\le\frac{a_+\rho}{8}d(p,p').
\]
The terminal Bernoulli kernel has total variation $|p-p'|\le d(p,p')/4$. Hence all loss-relevant reports satisfy \textup{(C)}. An equally spaced grid in $p$, the update \eqref{eq:hmm-update}, and exact Borel threshold comparisons define the mathematical label encoder. This is not yet a finite-bit implementation. For rational calibrated constants, Appendix~\ref{sec:algorithms} gives the separate finite-bit routine, including current-report quantization, interval arithmetic and absolute-error conversion. Approximate input and arithmetic with certified absolute errors are absorbed into $u_t$; denominator and log-odds conversion bounds above translate them to the common metric. For rational calibration, fixed-precision interval arithmetic produces any prescribed positive update accuracy; an exact tie need not be resolved, since approximate nearest-grid selection suffices. The persistent datum is the grid index; a bounded-precision register for the current report and arithmetic is temporary, not a retained real state. Summing the geometric error series gives the last error assertion. Apply Theorem~\ref{thm:main} to conclude.
\end{proof}

This is a repeatedly updated nonlinear filter, not a once-prepared finite measurement ensemble. The acquired law is generated by the actual observation density at every positive time. Formula~\eqref{eq:hmm-law} concerns prediction along the selected exploration, not the unrestricted optimal information-gathering controller.

\subsection{Singular geometric acquisition with an absorbing report}\label{sec:sensor}
We next construct a class with no strict contraction during waiting and with genuinely intersecting, different-dimensional acquired supports. Fix finitely many compact boxes $Q_j$ and maps
\[
 z_j:Q_j\longrightarrow[1/4,3/4]^d
\]
that are bi-Lipschitz with uniform constants in physical Euclidean coordinates. They may be nonlinear and their images may intersect. A dimension-zero box is a point. Prepare a latent mark $J_0$ with probabilities $\pi_j$ and, conditionally on $J_0=j$, a point $V$ uniformly on $Q_j$. The hidden target is $Z=z_{J_0}(V)$. After preparation the raw hidden state is $(Z,A_t)$; the latent mark is retained only on the preparation probability space for analyzing its pushforward law. The mark and any zero-width source coordinates are not supplied to the observer.

The apparatus has an unresolved/absorbed flag $A_t\in\{0,1\}$, initially zero. At each raw attempt, when unresolved, an independent coin with success probability $s\in[0,1]$ is tossed. Failure reports $\bot$ and changes nothing. Success reports $Z$ and sets the flag to one. Thereafter every raw attempt reports $\bot$ and preserves $Z$ and the flag. All attempts, including these holds, cost one acquisition and one unit of physical time. The action ``attempt'' is available in both modes, so no unrecorded action-legality bit is assumed. The terminal menu consists of a flag probe returning $A_N$, and $d$ binary probes with success probabilities $Z_i$, each executed only once as a terminal experiment. Give these $d+1$ probes equal weights.

After failure while unresolved, the hidden target distribution is unchanged: failure probability is independent of $Z$. After a successful report the future target probabilities are exactly the reported $Z$, whether or not the latent mark is identifiable. Thus the predictive quotient is the unresolved point $o$, together with the \emph{glued} set of absorbed predictions $z\in\bigcup_jz_j(Q_j)$. At an intersection the mark is quotiented out, not secretly retained. The flag probe distinguishes $o$ from an absorbed state even when their coordinate means coincide.

Set
\begin{equation}\label{eq:sensor-baseline}
 b_* =\frac1{d+1}\sum_{i=1}^d\E[Z_i(1-Z_i)],\quad
 V_* =\frac1{d+1}\sum_{i=1}^d\Var(Z_i),\quad
 h_N=(1-s)^N.
\end{equation}
For $N=0$ take $h_0=1$, including $s=1$. Let $\Psi^{\rm acq}(m)$ be \eqref{eq:psi} for the absorbed component charts with weights $\pi_j$. Widths there are the actual physical chart widths with the fixed task-weight scaling absorbed into the chart constants.

\begin{theorem}[Stratified acquisition law through rank collapse]\label{thm:sensor}
For every $s\in[0,1]$, $N\ge0$, and $m\ge2(J+1)$,
\begin{equation}\label{eq:sensor-law}
 R_\cp(N,m)-b_*\asymp R_\on(N,m)-b_*
 \asymp h_NV_*+(1-h_N)\Psi^{\rm acq}(m).
\end{equation}
The constants depend only on the finite chart and task constants, not on $s,N$, positive widths, or the distance between intersecting strata. When positive widths vanish, zero coordinates are removed from the predictive quotient and the same assertion holds, including point components. If all component score images (including the unresolved point) are separated, Theorem~\ref{thm:main} gives the all-budget version using the full mixture allocation \eqref{eq:psi}, for every $m\ge1$.

The online construction has $\Gamma_N=1$ in the full-mixture transport certificate. It retains no success count or latent component mark. Finite-precision successful reports with metric error at most $\eta$ give a root-mean-square prediction allowance at most $C\sqrt{1-h_N}\eta$, when all hold operations are exact.
\end{theorem}
\begin{proof}
The preparation is the pushforward of the specified box mixture, and the transition is a normalized nonnegative Borel kernel on the compact target space with its two flags. For the success report component one may take the prepared target law as dominating measure; its support is the compact union of the chart images. The success instrument at the unresolved belief has density $s$, and its determining-test numerator has the continuous version $s f(z)$. At absorbed beliefs both densities are zero. The $\bot$ component is atomic. Thus the explicit overwrite below provides exactly the fibre-compatible zero-density extension required in Proposition~\ref{prop:quotient}. The admissible belief domain is the unresolved prepared law together with the point beliefs at absorbed targets. It is a disjoint compact union. All determining future-test evaluations are continuous on it; the explicit overwrite/hold versions below also verify the zero-denominator extension clause of Proposition~\ref{prop:quotient}. The probability of $N$ consecutive failures is exactly $h_N$. Conditional on success by $N$, the target law is still the prepared mixture, because all acquisition coins are independent of it. The acquired predictive law is therefore
\begin{equation}\label{eq:sensor-acquired}
 h_N\delta_o+(1-h_N)\sum_j\pi_j\,(z_j)_\#\operatorname{Unif}(Q_j),
\end{equation}
with the absorbed flag coordinate included in the score. In particular the formal union of all prepared values does not have weight one before it is acquired.

The conditional terminal variance of the flag is zero. On an absorbed history the coordinate variance is $Z_i(1-Z_i)$. On an unresolved history it is
$\E[Z_i(1-Z_i)]+\Var(Z_i)$, by total variance. Averaging proves
\begin{equation}\label{eq:sensor-B}
 B_N=b_*+h_NV_*.
\end{equation}
Conditioned on the raw component mark, the acquired charts have exactly uniform physical box densities. Their bi-Lipschitz constants give \textup{(G)}. Their union can have intersections, different local dimensions, and atomic mass. All these properties follow from the raw preparation and report map.

Equip the quotient with Euclidean distance on its weighted score vector. Report continuity holds as follows. Between two absorbed states the next-attempt kernels are identical, namely $\delta_\bot$. Between unresolved and absorbed modes, the variation distance is at most one, while the flag-score distance is $1/\sqrt{d+1}$; a finite global Lipschitz constant follows. Terminal coordinate kernels are Bernoulli and their variation distance is $|Z_i-Z'_i|$; the flag kernel is bounded by the same mode-separation argument. There is no assertion that distinct Dirac reveal channels are TV-continuous: the reveal action is used only from the single unresolved quotient point. On a successful report define the predictive update to overwrite with that reported absorbed state from every previous representative, including the conditional version at a null report in the absorbed mode. On $\bot$, it is the identity.

The machine initially encodes $o$. At a success it chooses a nearest representative in the union codebook, using the reported physical coordinates; on every failure or later hold it leaves the label unchanged. It never reads $J_0$. Conditional on any raw mark $j$, the success-rounding error is bounded by that chart's cover. Therefore the coefficients are exactly those of Proposition~\ref{prop:mechanisms}(c), with the unresolved point as component zero. This proves $\Gamma_N=1$, irrespective of how many holds occur. It also proves the finite-precision allowance: a single error of size $C\eta$ is introduced on an event of probability $1-h_N$, and none on the complementary event.

For completeness one can see the separated oracle and acquisition terms directly. Lemma~\ref{lem:projection}, conditioned on the success event and its latent mark, and the proof of Corollary~\ref{cor:crossing} give an excess of at least $c(1-h_N)\Psi^{\rm acq}(m)$, even if the decoder also uses centers for the unresolved point. For the upper use one exact unresolved label and $m-1$ absorbed grid labels. Allocation scaling gives
$\Psi^{\rm acq}(m-1)\le C_J\Psi^{\rm acq}(m)$ for $m\ge2(J+1)$: allocating $\lfloor(m-1)/J\rfloor$ to every chart and using the estimate in Corollary~\ref{cor:crossing} proves it. The finite machine just constructed then incurs at most $C(1-h_N)\Psi^{\rm acq}(m)$. Add \eqref{eq:sensor-B}. At $s=0$ or $N=0$ the success term is zero and a single known forecast attains the acquisition term. With zero surviving widths all component charts are points and sufficiently many labels give zero acquired distortion, as the formula states.

For effective implementation, take algebraic calibrated widths and weights, and chart maps with certified finite-precision evaluation. The finite allocation minimum is an algebraic comparison problem. A nearest representative can be selected to within absolute distance $\eta$ by comparing approximated distances; this adds at most $2\eta$ to the best covering radius. No exact comparison of arbitrary input reals is required. A finite readonly table of approximated representatives or a grid-generation program, with separately charged description and temporary workspace, suffices. Thresholds are compared in physical coordinates; no uniform error estimate divides by a collapsing width. If the unresolved target mean also requires numerical integration, its certified score error $\eta_0$ is charged separately; the total readout allowance is then at most $\sqrt{h_N}\eta_0+C\sqrt{1-h_N}\eta$. A table or algorithm for this mean is part of the supplied effective calibration, not a free integration oracle. Appendix~\ref{sec:algorithms} distinguishes these costs. This proves the claimed implementation statements.
\end{proof}

\paragraph{A concrete singular family.}
In two coordinate dimensions take equal mixture weights on
\[
 z_1(u)=(1/2+u,1/2),\quad |u|\le a/2,
 \qquad
 z_2(v,w)=(1/2+v,1/2+w),\quad |v|\le b/2,\ |w|\le c/2,
\]
where $0\le a,b,c\le1/4$. The actual acquired law has a one-dimensional singular component on a line and a two-dimensional component on a rectangle. The former lies inside the latter when their horizontal ranges overlap. If $c=0$, the rectangle becomes a line; if also $b=0$, it becomes an atom. Before acquisition the support consists only of the unresolved predictive point. The physical scales are
\[
 \psi_1(k)=a^2/k^2,\qquad
 \psi_2(k)=\max\{\max(b,c)^2/k^2,bc/k\}\quad(bc>0),
\]
with zero widths removed at the boundary. There is no single smooth acquired manifold at the crossing or a single ambient-dimension exponent. Theorem~\ref{thm:sensor} applies with constants uniform through all these changes. Smooth nonalgebraic bi-Lipschitz distortions of the two charts yield the same conclusion, with their uniform metric constants. Thus the extension is not a statement about a preassigned Hilbert-space rank cut.

\subsection{Recurrent expansion with intermittent refresh}\label{sec:recurrent}
We give a third realization to test the structural kernel theorem, rather than only its contractive and absorbing special mechanisms. Fix $l=1/4$, $0\le a\le1/2$, and $u=l+a$. Prepare $X_0$ uniformly on $[l,u]$; for $a=0$ this is a point. One initial paid call reveals $X_0$. At each subsequent paid tick, independently of the past, a hold occurs with probability $1-\theta$, $0\le\theta\le1$. On a hold, $X$ is unchanged and the report is the symbol $H$. Otherwise an active type is chosen independently: with probability $p=1/10$ it is the fold $E$, and with probability $1-p$ it is the refresh $C$. The raw transition/report kernels are
\begin{align}\label{eq:recurrent-kernel}
 E:&\quad X'=l+2\min\{X-l,u-X\},\quad Y=E,\\
 C:&\quad X'=cX+(1-c)(l+aU),\quad Y=(C,U),
           \quad c=1/4,\quad U\sim\operatorname{Unif}[0,1].\notag
\end{align}
These are normalized nonnegative Borel kernels; the fold is deterministic. There is no further exact reveal of $X'$. A terminal probe is Bernoulli with mean the current $X$. All ticks, including the initial reveal and holds, cost one raw call. The exploration has no choice of informative actions and is fixed. Its protocol clock is exogenous.

\begin{theorem}[Acquired geometry under recurrent expanding updates]\label{thm:recurrent}
For every $N\ge1$, $m\ge1$, $a\in[0,1/2]$ and $\theta\in[0,1]$, the raw experiment \eqref{eq:recurrent-kernel} has predictive quotient $[l,u]$ with its protocol interface, and
\begin{equation}\label{eq:recurrent-risk}
 R_{\cp}(N,m)-B_N\asymp R_{\on}(N,m)-B_N\asymp a^2m^{-2},
 \qquad B_N=\E[X_N(1-X_N)].
\end{equation}
Here $X_N$ denotes the state after the $N$ paid calls, including the initial reveal. The constants are absolute and uniform in all displayed parameters. At $a=0$ both exact excesses are zero. The online bound has the kernel-derived transport constant
\begin{equation}\label{eq:recurrent-gamma}
 \Gamma\le\frac{3040}{667},
\end{equation}
including initial rounding, independently of the hold probability.

If current inputs, physical endpoints and arithmetic have absolute errors controlled by $\eta$ per initial or active update, and hold operations are exact, the finite-bit implementation has
\[
 R_{\rm alg}(N,m)-B_N\le C\bigl(a/m+\eta\bigr)^2.
\]
The upper is independent of $N$ and $\theta$. The endpoint/error guarantee includes a projection back to the calibrated valid interval with its stated absolute accuracy. Its program and workspace costs are those of Appendix~\ref{sec:algorithms}; no stored exact real is used.
\end{theorem}
\begin{proof}
After the initial report the exact state is known. The reported type and, at refresh, the reported $U$, determine each subsequent exact update. All future laws are therefore functions of the current $x$ and the protocol interface. The terminal Bernoulli mean is $x$, so distinct states are predictively distinct. Continuation laws and these explicit updates are continuous on the interval within each report component. The report law before the terminal probe is independent of $x$, hence has zero report-kernel continuity constant. The terminal report has total variation $|x-x'|$. This establishes the quotient and report certificates directly from the raw kernel, without assuming a covariance geometry.

For $a>0$, let $f$ be the actual current density. Initially $f=1/a$. A hold preserves it. The tent fold has two affine inverse branches, each with Jacobian factor $1/2$, so its pushforward density is the average of the two branch densities and has sup norm at most $\|f\|_\infty$. For refresh, condition on the preceding state. The new state is uniform on an interval of length $(1-c)a$; its conditional density is at most $1/((1-c)a)=4/(3a)$. Mixing over the actual preceding law preserves this bound. Induction, including mixing over the reported type and holds, gives
\[
 0\le f_N(x)\le4/(3a),\qquad \int_l^u f_N(x)\,dx=1
\]
for every $N\ge1$. Thus the actual-mass certificate has $C_g=4/3$ on one chart of physical width $a$. It is a bound on the law generated by the experiment, not an imposed invariant uniform law. The task map is $x\mapsto x$ and $\psi(m)=a^2/m^2$.

Use a midpoint grid with $m$ representatives and physical Euclidean distance. Execute holds without arithmetic or rounding. Execute an active update from the previous representative and then round. The fold has Lipschitz constant two; refresh has constant $c=1/4$. Their rounding error is at most $a/(2m)$. Therefore the error recurrence has active gain $G=2$ with probability $1/10$ and $G=1/4$ with probability $9/10$, with forcing one after scaling the grid radius. The coefficients are independent marks, independent of the grid. In the one-state environmental kernel,
\[
 \E G=17/40,\qquad \E G^2=73/160<1.
\]
Theorem~\ref{thm:kernel} and \eqref{eq:scalar-resolvent} yield
\[
 H=\frac{1+17/40}{(1-17/40)(1-73/160)}=\frac{3040}{667}.
\]
The normalized initial rounding seed is at most one; the scalar fixed-point first and second moments are both at least one, so it is an admitted seed. Proposition~\ref{prop:suspension} proves that insertion of holds leaves this bound unchanged for $\theta>0$. At $\theta=0$ only the initial rounding persists, so the same upper holds directly. Apply Lemma~\ref{lem:initialization} to the one paid initial reveal and the $N-1$ subsequent calls. Theorem~\ref{thm:main} then supplies \eqref{eq:recurrent-risk}, and its lower uses the actual density just proved. No realization calculation has entered the proof of the kernel theorem.

On an active or initial update, interval projection and approximate arithmetic add at most $C\eta$ to the same error recurrence. On holds the representative is copied exactly, giving zero new error. The same scalar forcing bound, applied to radius $a/(2m)+C\eta$, proves the finite-bit assertion. The physical formulas in \eqref{eq:recurrent-kernel} use $a$ only multiplicatively; their absolute-error bounds never divide by it. At $a=0$ there is a single exact state, and the exact program can bypass all rounding. This also proves rank-uniformity.
\end{proof}

For every $\theta>0$ the active folds occur repeatedly, and blocks of consecutive folds have positive probability. Thus this example is not uniformly contractive along active paths, not an absorbing reveal after the initial preparation, and not a finite measurement ensemble. Its stability is a forced second-moment conclusion derived from the raw update probabilities. Uniformity in arbitrarily long waiting periods does not remove those periods from the acquisition or physical-time budget.

\section{Consequences and comparison with existing theory}\label{sec:scope}
The central theorem separates actual acquisition, geometry and causal implementation, then proves a quantitative connection. Its forward gain matrix need not be diagonal in the predictive strata. Its occupation comparison is derived from raw invariant-reference bounds in Proposition~\ref{prop:nr-reference}, without identifying the reference measure with the actual law. Lemma~\ref{lem:nr-overlap} permits a finite multiplicity of overlapping score neighborhoods. Proposition~\ref{prop:nr-acquisition} treats atomic acquired laws by retaining the unresolved-variance term rather than imposing a false fine-scale small-ball bound.

The recurrent prefix experiment gives an explicit common-task law $\delta^2+\E\ell_{\min(H_T,\lfloor\log_2M\rfloor)}^2$. Its known depth is generated by the very same nonreset acquisition process that the encoder must compress. The older one-probe and selected-menu minimax theorems remain independent sharpness witnesses, not proofs of adaptive acquisition or of a universal full-resource region. The exact and noisy continuation converses likewise constrain information that truly survives a cut, rather than every register in one preferred implementation.

\subsection{Theorem-level comparison}
\paragraph{Predictive states and sufficient statistics.}
Shalizi and Crutchfield's minimality and uniqueness theorems for causal states \cite[Theorems 2--3]{shalizi} identify minimal prescient representations of a stochastic process. Predictive state representations \cite{psr} use future action--observation tests. Our quotient proposition does not claim priority for either construction. It explicitly includes controlled continuation legality, paid resources and conditional-version domains. Its extra measurable work is the density-instrument descent at every declared continuous report. The categorical treatment of kernels and sufficiency in \cite{fritz} supplies another general language; an abstract sufficient statistic alone does not provide the acquired-mass or finite-machine estimates used here.

\paragraph{Bisimulation and belief-state control.}
Behavioral metrics for Markov decision processes \cite{ferns} compare states through rewards and transition laws and control value differences. Our metric is fixed by executable prediction tasks and valid update geometry; it is not asserted to be a universal bisimulation metric. Saldi, Y\"uksel and Linder \cite[Assumption 1 and Theorem 3]{saldi} obtain near-optimal finite belief-model policies for discounted partially observed control under their continuity, compactness and moment hypotheses. That theorem optimizes a controller, whereas our matching lower is for a fixed exploration and the law it actually acquires. Our finite-horizon two-sided resolution estimate does not imply their unrestricted control conclusion, nor does asymptotic policy near-optimality supply our actual-mass lower bound.

\paragraph{Quantization and nonlinear filtering.}
Conditional-mean projection and optimal probability quantization are classical \cite{graf}. Nonlinear filter approximation by quantization is developed in \cite{pages,sellami,feuer}. The additional obligation here is a matching lower for the very same acquired task, together with the forced-moment condition needed by a persistent finite machine. Zhu's local estimates for Ahlfors--David measures \cite{zhu}, and for Moran measures \cite{zhu-moran}, concern fine uniformity of contributions of optimal cells and successive quantization errors. The cylinder covers in Lemma~\ref{lem:moran-profile} and Theorem~\ref{thm:nr-moran} are not offered as replacements for that theory. Its role is to verify a multiscale profile, then propagate that profile through an actual nonstationary recurrent experiment with countably many weighted components. The finite-budget anisotropic results additionally retain width collapse and finite intersecting mixtures. We do not classify arbitrary singular measures.

\paragraph{Positive moment operators.}
Average contraction of iterated random functions \cite{diaconis} and moment stability of Markov jump systems \cite{costa} are established subjects. Ogura and Martin \cite{ogura} characterize mean stability for positive semi-Markovian jump systems with resets by a spectral-radius condition. That homogeneous stability question is not the joint domination of forced error measures by changing acquired measures required here. We claim no novelty for a Neumann resolvent or the criterion $\rho(K_2)<1$. The role of Theorem~\ref{thm:kernel} is more specific: reverse-edge operators propagate the conditional first and mixed second moments of the \emph{rounding forcing}, and the resulting matrix is normalized by the actual acquisition weights. Exact suspension cancels in the forced profiles, not in the physical budget. This distinction is why a recurrent expanding kernel can still yield a matching finite-memory risk curve. The new forward drift in Theorem~\ref{thm:nonreset} controls all cross-stratum old-error edges and all allocation-dependent radii, after occupation comparison, without that covariance normalization. Its matrix condition is standard positive stability, but its coupling to the actual-law lower and executable online construction is the claim studied here. In the complementary reset-preserved class, Theorem~\ref{thm:bridge} propagates forward error measures and obtains the terminal mass factor directly; Proposition~\ref{prop:recharge-flow} identifies the necessary refresh-versus-expansion balance used by that sufficient criterion. We do not claim that this sufficient condition is a necessary spectral classification of all raw kernels.

\paragraph{Causal rate distortion and sequential reconstruction.}
Nonanticipative rate-distortion theory imposes realizability constraints on reconstruction kernels and relates information minimization to filtering on abstract spaces \cite{nonanticipative}. A communication or directed-information budget is not, without a separate coding theorem, the cardinality of the complete state retained across each computational cut. Our theorem gives explicit recursive labels and a checkpoint converse in the same task, not an identification of those resource models. The finite-prefix acquisition theorem counts word length and mode in a fused finite alphabet, and separately records input, arithmetic and controller costs.

\paragraph{Comparison of experiments.}
Blackwell comparison \cite{blackwell}, Le Cam's deficiency theory \cite{lecam}, and Torgersen's systematic account \cite{torgersen} provide the classical common-decision-loss principle. The triangle inequality and total-variation contraction in Theorem~\ref{thm:morphism} are instances of that principle. The added structure is the finite-resource causal implementation, with explicit lifted controller classes, stopping indices, clocks and numerical interfaces. A forward simulator yields an upper risk transfer, not a lower bound relative to a different oracle. The mandatory-state and robust converses are proved directly in Theorems~\ref{thm:tag} and \ref{thm:robust}, rather than attributed to a forward deficiency bound.

\subsection{Necessary distinctions illustrated by exact examples}
First, the valid envelope $[0,1]$ equipped with a point mass and the same envelope equipped with a uniform law have respectively zero one-point distortion and strictly positive finite-label distortion. Hence a reachability envelope is not an acquired law. More generally a component of probability $w$ contributes with weight $w$, even if its geometric diameter is large.

Second, an identity update has no need for repeated quantization. If the initial state has already been replaced by a representative, the same representative can be retained indefinitely. The bound $Nr$ obtained by blindly injecting a covering radius at every hold is not an inherent resource requirement. The acquisition-weighted certificate corrects precisely this failure of accounting.

Third, a statistical garbling can increase memory needed to match its \emph{own} oracle. Observing a fair hidden bit exactly produces only posterior probabilities zero and one. Adding independent Gaussian noise to this observation produces a continuously varying posterior probability. Two labels reproduce the first oracle exactly, but no finite number reproduces the latter continuously distributed oracle exactly under square score. This is compatible with data processing because the baselines differ.

Fourth, let $P_N=\delta_0$ and $Q_N=(1-e^{-Nc})\delta_0+e^{-Nc}\delta_1$, $c>0$. Their total variation distance tends to zero, but the exponential rates at the point one are infinite and $c$, respectively. The common-risk defect of this paper is therefore not a certificate of large-deviation equivalence.

Finally, a hidden mode or an internal simulator register is not free because its name is omitted from a state vector. Independent pairs of retained states have the product cardinality in Theorem~\ref{thm:morphism}. Conversely that product is a constructive upper bound, not a theorem that every simulator requires it. Theorem~\ref{thm:tag} establishes exact support necessity through zero-error continuation challenges; Theorem~\ref{thm:soft-state} replaces it by an information lower bound when errors are allowed. An exact label machine and a finite-bit numerical implementation are also different objects: the latter retains the nonzero calibration, workspace, and readout terms of \eqref{eq:joint}. Exact-arithmetic zero distortion is not a promise of a finite-precision physical output with zero error.


\subsection{Uniformity and the remaining mathematical questions}
Theorem~\ref{thm:nonreset} allows nonreset transfer between strata, expanding edges, nonstationary acquired laws comparable to an invariant reference, and finite or countable geometric partitions with bounded overlap. The uniform constants are explicit in the drift and occupation bounds. They are not an assertion that every raw kernel admits such bounds. In particular, rapid creation or extinction of strata can violate occupation comparison; the separate reset-preserved forward-measure theorem still handles its own nonstationary refresh-flow class. No necessity characterization of optimal finite-memory prediction is inferred from the spectral-radius criterion for a comparison matrix.

The countable overlap result requires finite multiplicity at the declared relative scales and dyadic regularity of the profiles. Arbitrary unbounded-multiplicity intersections or all singular measures remain outside it. The finite intersecting-stratum and rank-collapse results remain independently available. The singular nonreset examples prove a nonuniform recurrent extension, not a classification of all nonlinear filters, diffusions or hidden-state systems.

The recurrent acquired-depth law treats a known raw kernel and a fixed, fully charged acquisition protocol, with a two-point calibration ambiguity that the reports cannot learn. It does not optimize the division of the raw budget between initial readings and recurrence, estimate unknown transition kernels, or supply an optimal pilot--calibration design. These remain mathematical responsibilities within the same mother problem. A universal matching region for program length, transient workspace, physical runtime and simulator state is also not proved. The information lower bounds and explicit finite-bit implementations give narrower, stated claims rather than an equality for that entire vector.

Finally, a deterministic-horizon energy estimate is not automatically a stopping-time estimate: a stopping rule can select large errors. The bounded stopped common-law comparisons of Theorem~\ref{thm:morphism} retain their separate uniform hypotheses. Infinite-horizon transcript comparison, large-deviation equivalence and unbounded-generator closure require different estimates. Neither successful compilation nor finite regression checks establish any of the continuum theorems proved here.

\appendix
\section{Terminal acquired mass and causal resolution}\label{sec:bridge}
A stationary moment bound need not be commensurate with a rare terminal component. We therefore propagate an error \emph{measure}, rather than a single unconditional moment. The reference measure in the propagation is the actual acquired law at the same time. The resulting theorem applies to nonstationary kernels and to countably many accumulating components. No stationary law or inverse terminal mass occurs.

\subsection{A multiscale geometric profile}
Fix the experiment, exploration, prediction task and horizon of Section~\ref{sec:experiment}. Let $(D,d)$ be a bounded valid predictive domain with a valid anchor $o$. There are disjoint Borel acquired components $D_j$, $j\in\mathcal J$, where $\mathcal J$ is finite or countable. Their union has full acquired mass at every time considered. The anchor may have zero acquired mass. The task map $p_N:D\to\R^q$ is $L$-Lipschitz. Write $\mu_t$ for the actual law of the predictive Markov factor specified below, $\mu_t^j$ for its restriction to component $j$, and $w_{t,j}=\mu_t^j(D_j)$. Extra factor coordinates are suppressed in this notation; the geometric maps act on their predictive coordinate. At positive terminal weights let $\nu_{N,j}$ be the conditional law of the terminal score. The following are uniform certificates, not definitions of an optimum.

\smallskip\noindent\textbf{(MG) Mass, covering and scale separation.}
For constants $C_g,c_g,c_s,D_r>0$ there are scales $0<a_j\le a_*$ and nonincreasing profiles $r_j:\N_0\to(0,\infty)$ such that
\begin{equation}\label{eq:profile-regular}
 r_j(0)=r_j(1)=a_j,\qquad r_j(k-1)^2\le D_r r_j(k)^2\quad(k\ge2).
\end{equation}
For every $k\ge1$ there are at most $k$ valid representatives covering $D_j$ in the predictive metric with radius $C_g r_j(k)$; also $d(s,o)\le C_g a_j$ on $D_j$. All such representatives are in the valid update domain. Put $S_{N,j}=p_N(D_j)$, assumed compact. The open sets
\[
 \{z:\operatorname{dist}(z,S_{N,j})<c_s a_j\}
\]
are pairwise disjoint, and, for every $k\ge1$ and every $z\in\R^q$,
\begin{equation}\label{eq:profile-smallball}
 \nu_{N,j}\big(B(z,c_g r_j(k))\big)\le(2k)^{-1}.
\end{equation}
The latter bound is required only at $w_{N,j}>0$. These laws are obtained from the preparation and raw kernels, not from a substitute geometric measure. The cover is an intermediate-time valid envelope; it is not asserted to be the support of every intermediate acquired law. Constants are uniform over the declared horizons and component indices. For countably many components, $a_j\to0$ along an enumeration is a convenient compactness certificate, but the estimates require only the displayed conditions. Atoms are handled by the separate finite-stratum theorem in Section~\ref{sec:main}; they do not satisfy \eqref{eq:profile-smallball} for all $k$.

For an integer $m\ge0$ define the explicit acquired profile
\begin{equation}\label{eq:bridge-profile}
 \mathcal Q_N(m)=\inf_{\substack{k_j\in\N_0,\ \sum_jk_j\le m}}
             \sum_j w_{N,j}r_j(k_j)^2.
\end{equation}
An allocation has finite support. Its zero entries mean that the corresponding entire components are discarded to the anchor. The infimum need not be attained. This profile is specified by actual masses and separately verified cover/small-ball scales; it is not defined in terms of an optimal encoder. Condition~\eqref{eq:profile-regular} implies
\begin{equation}\label{eq:one-label}
 \mathcal Q_N(m-1)\le D_r\mathcal Q_N(m),\qquad m\ge1,
\end{equation}
with $D_r$ enlarged to at least one: remove one label from any nonzero entry of an approximate allocation; removing the sole label of a component changes nothing because $r_j(0)=r_j(1)$. Then take infima. If the allocation has fewer than $m$ labels, no removal is necessary.

\subsection{An experiment-level balance of error measures}
Let $Z_t$ be a standard-Borel factor of the actual experiment, carrying its predictive state and any controller coordinates needed for the following property. Conditional on the entire preceding history and algorithm randomness, the law of the next factor and analysis mark is
\[
 P_t(dz',d\xi\mid Z_{t-1}).
\]
In particular, a new mark is conditionally independent of the retained approximation error given $Z_{t-1}$. This is a Markov-factor requirement on the chosen fixed exploration, not a claim that every predictive coordinate is Markov under every controller. The finite programs do not read analysis-only coordinates.

A mark specifies either an exact hold, a non-hold continuation, or an overwrite. Holds preserve the exact state and its representative. \textbf{On every non-overwrite edge the actual component is preserved; component changes are overwrites.} This structural restriction belongs to the present complementary theorem, not to Theorem~\ref{thm:nonreset}. An overwrite may enter any component and removes dependence of the exact update on the previous state. Fix $\eta>0$. Nonnegative measurable $g_t(z,\xi,z')$ and $b_t(z,\xi,z')$ are checked from the raw update and the rounding procedure, independently of the allocation. On overwrites set $g_t=0$. For every finite allocation there is a stated Borel finite-label machine, using its label and current report only, for which
\begin{equation}\label{eq:bridge-local}
 e_t\le g_t e_{t-1}+b_t R_j\quad\hbox{on non-holds entering }D_j,
 \qquad e_t=e_{t-1}\quad\hbox{on holds}.
\end{equation}
Here $e_t=d(s_t,\widehat s_t)$ and $R_j=C_g r_j(k_j)$ for the exact model. A finite-precision implementation instead may use $R_j=C_g r_j(k_j)+\xi_j$, with specified deterministic local error allowances $\xi_j\ge0$. In both models the retained machine must implement the declared rule even when different discarded components share the anchor. The analytic index $j$ is not an extra free state register. Section~\ref{sec:moran} gives this implementation explicitly.

Define two positive edge kernels. The first, $A_{t,j}$, runs from component $j$ to itself, with mark weight one on holds, $(1+\eta)g_t^2$ on non-hold non-overwrites, and zero on overwrites. The second, $F_{t,j}$, runs from the entire factor space into component $j$, with weight $(1+\eta^{-1})b_t^2$ on non-holds and zero on holds. Both are obtained by integrating these weights against $P_t$. Write $\lambda A(C)=\int A(z,C)\lambda(dz)$.

\smallskip\noindent\textbf{(TB) Terminal-mass recharge.}
There are known measurable functions, independent of the allocation and local precision allowances, $1\le \upsilon_{t,j}\le V<\infty$ such that, as measures on the destination factor space,
\begin{equation}\label{eq:recharge}
 (\upsilon_{t-1,j}\mu_{t-1}^j)A_{t,j}+\mu_{t-1}F_{t,j}
       \le \upsilon_{t,j}\mu_t^j,\qquad 1\le t\le N.
\end{equation}
The initialized error satisfies
\begin{equation}\label{eq:bridge-seed}
 \E[e_0^2;Z_0\in C\cap D_j]\le R_j^2\int_{C\cap D_j}\upsilon_{0,j}\,d\mu_0.
\end{equation}
The initialization, including any reveal, is paid for in the raw-call ledger. Inequality~\eqref{eq:recharge} is a one-step inequality between known forward measures and marked raw kernels. It contains neither an optimal risk nor an unrolled forced-error moment. Proposition~\ref{prop:recharge-flow} derives it from a scalar balance of raw incoming refresh mass and expansion mass, uniformly through arbitrarily rare components. Conversely, a mere unconditional mean contraction does not imply it.

\begin{theorem}[Terminal-mass acquired-geometry transfer]\label{thm:bridge}
Fix a prepared raw experiment, a legal exploration independent of the encoder, a deterministic horizon, and a common finite squared-prediction task. Suppose the predictive quotient is realized as in Proposition~\ref{prop:quotient}, or by a separately verified Borel realization with declared update versions. Assume \textup{(MG)}, the stated Markov factor and executable local rule, and \textup{(TB)}. Then, for every integer $M\ge2$,
\begin{equation}\label{eq:bridge-main}
 B_N+c\mathcal Q_N(M)\le R_\cp(N,M)\le R_\on(N,M)
       \le B_N+C V\mathcal Q_N(M).
\end{equation}
The last inequality is an inequality of infima. For each $A>1$ it is realized with $C$ replaced by $AC$ by a machine using one anchor and an allocation of at most $M-1$ other representatives whose cost is within factor $A$ of $\mathcal Q_N(M-1)$. If that infimum is zero, use arbitrarily small additive allocation errors instead. The constants $c,C$ depend only on $c_g,c_s,C_g,D_r,L$ and not on the number of components, their masses, the horizon, or the positive scales. The factor $V$ is displayed separately.

For an effective finite-precision rule satisfying \eqref{eq:bridge-local} with $R_j=C_g r_j(k_j)+\xi_j$, and terminal readout error at most $\zeta_N$ in score norm,
\begin{equation}\label{eq:bridge-finite}
 R_{\rm alg}\le B_N+
 \left[C\sqrt{A V\mathcal Q_N(M)}
       +L\sqrt{V\sum_jw_{N,j}\xi_j^2}+\zeta_N\right]^2.
\end{equation}
Here effective allocation, input access, workspace, program, and physical-time bounds are additional data, not consequences of a Borel cover. The precise full-ledger and common-loss composition is that of Theorem~\ref{thm:morphism}. In particular, root-mean-square state errors add before squaring, complete causal total-variation defects add in the common bounded-loss metric, and persistent cardinalities multiply only for the stated product implementation.
\end{theorem}

\begin{proof}
We first prove the terminal-mass assertion without using geometry. Define a finite error measure
\[
 E_{t,j}(C)=\E[e_t^2;Z_t\in C\cap D_j].
\]
On non-holds, Young's inequality and \eqref{eq:bridge-local} give
$e_t^2\le(1+\eta)g_t^2e_{t-1}^2+(1+\eta^{-1})b_t^2R_j^2$.
On holds the old error is copied exactly. On overwrites the old-error term vanishes. Non-overwrite transitions do not import error from another component. Conditional independence in the Markov-factor hypothesis therefore gives the \emph{measure} inequality
\begin{equation}\label{eq:error-measure-step}
 E_{t,j}\le E_{t-1,j}A_{t,j}+R_j^2\mu_{t-1}F_{t,j}.
\end{equation}
All conditional expectations are applied to nonnegative functions, so no integrability interchange is tacit. Starting from \eqref{eq:bridge-seed}, positivity and \eqref{eq:recharge} prove by induction that
\begin{equation}\label{eq:terminal-conditional}
 E_{t,j}\le R_j^2 \upsilon_{t,j}\mu_t^j,\qquad
 \E e_N^2\le V\sum_jw_{N,j}R_j^2.
\end{equation}
The sum is justified by monotone convergence for a countable partition. Components with zero terminal mass automatically have zero terminal error measure. There is no division by their masses. Conditional on a positive-mass component, the same bound is $\E[e_N^2\mid D_j]\le V R_j^2$.

We next supply the geometric lower, including arbitrary off-support decoder centers. For any $M$ centers assign a center to $j$ when it is within $c_s a_j$ of $S_{N,j}$. The neighborhoods in \textup{(MG)} are disjoint, so the assigned counts satisfy $\sum k_j\le M$. A component with $k_j=0$ is everywhere at least $c_s a_j$ from every center. For $k_j>0$, set $r=r_j(k_j)$ and $c_0=\min(c_g,c_s/2)$. Unassigned centers are at distance at least $c_s a_j\ge 2c_0r$. The union of the $k_j$ score balls of radius $c_0r$ around assigned centers has conditional mass at most $1/2$ by \eqref{eq:profile-smallball}. Hence the conditional squared distortion is at least $c_0^2r^2/2$. Summing with the actual weights proves a lower $c\mathcal Q_N(M)$. Randomized encoders or decoders cannot improve it: condition on independent shared randomness, replace a randomized decoder by its mean, and select the closest center for each full-history conditional prediction. Conditional projection, proved in Lemma~\ref{lem:projection}, identifies this distortion with excess risk above $B_N$.

For the online upper, take a finite allocation of at most $M-1$ labels and the anchor. The executable local rule and \eqref{eq:terminal-conditional} give excess at most $L^2VC_g^2\sum w_{N,j}r_j(k_j)^2$. Approximate the infimum and apply \eqref{eq:one-label}. Every online encoder is a checkpoint encoder, proving the middle inequality. For finite precision apply Minkowski to the weighted $\ell^2$ norm of $C_g r_j(k_j)+\xi_j$, then to the terminal score and its readout approximation. This gives \eqref{eq:bridge-finite}. The lower concerns the more informative exact-input checkpoint class and therefore remains a lower for the restricted finite-input class. Resource and deficiency composition uses the same task and the explicitly lifted strategy, as proved independently in Section~\ref{sec:transfer}.
\end{proof}

\subsection{A raw-flow condition that supplies the balance}
The preceding theorem is not restricted to the following class, but this class makes its recurrent rare-component assertion directly verifiable. Let $\sigma_j$ be probability laws on component $j$. A paid initial preparation gives $\mu_0=\sum_jw_{0,j}\sigma_j$. At time $t$ the state-independent command/report type probabilities are $h_t,u_t,v_t\ge0$, summing to one. A hold copies the state; a continuation has kernel $T_{t,j}$ preserving $\sigma_j$; a refresh draws a new component with known probabilities $\pi_{t,j}$ and draws its state from $\sigma_j$, independently of the previous state. The report reveals the refresh state. The actual component masses obey
\begin{equation}\label{eq:raw-flow}
 w_{t,j}=(h_t+u_t)w_{t-1,j}+v_t\pi_{t,j}.
\end{equation}
Assume the gain-weighted continuation satisfies
\[
 \int g_t(z,\xi,z')^2\,P_t^{\rm cont}(dz',d\xi\mid z)\,\sigma_j(dz)
       \le G^2\sigma_j(dz'),
\]
where $P_t^{\rm cont}$ is conditional on continuation; and $b_t\le b_*$ on non-holds. This weighted invariance is stronger than an unconditioned bound on $\E g_t^2$; it controls where transported errors arrive. Deterministic measure-preserving maps with Lipschitz bound $G$ satisfy it.

\begin{proposition}[Recurrent refresh balance without inverse rare masses]\label{prop:recharge-flow}
Let $A_G=\max\{(1+\eta)G^2-1,0\}$ and $C_b=(1+\eta^{-1})b_*^2$. Suppose there is $K>A_G$ such that
\begin{equation}\label{eq:flow-balance}
 v_t\pi_{t,j}\ge K u_t w_{t-1,j}\quad\hbox{for all }t,j.
\end{equation}
Then \textup{(TB)} holds with constant functions $\upsilon_{t,j}=D$, where
\begin{equation}\label{eq:flow-constant}
 D\ge\max\left\{1,D_{\rm seed},\frac{C_b(K+1)}{K-A_G}\right\}.
\end{equation}
Here $D_{\rm seed}$ bounds the normalized initial error measure. The conclusion is uniform as any $w_{t,j}$ or total activity $u_t+v_t$ tends to zero. The probabilities and refresh distributions can vary with time. No stationarity of $\mu_t$ is required.
\end{proposition}
\begin{proof}
Direct integration of the raw kernels proves inductively that $\mu_t^j=w_{t,j}\sigma_j$ and gives \eqref{eq:raw-flow}. The left side of \eqref{eq:recharge} is bounded by
\[
 \left[D\{h_tw_{t-1,j}+(1+\eta)G^2u_tw_{t-1,j}\}
       +C_b\{u_tw_{t-1,j}+v_t\pi_{t,j}\}\right]\sigma_j.
\]
Set $x=u_tw_{t-1,j}$ and $y=v_t\pi_{t,j}$. Subtracting this expression from $D w_{t,j}\sigma_j$ leaves at least
$[(D-C_b)y-(D A_G+C_b)x]\sigma_j$.
For $y\ge Kx$ this is nonnegative by \eqref{eq:flow-constant}; when $x=y=0$ it is zero. Holds cancel on the two sides and never enter the denominator. The seed is covered by $D_{\rm seed}$.
\end{proof}

For example, choose any sequence of probability vectors $\omega_t$, set
$\pi_t=(1-\alpha_t)w_{t-1}+\alpha_t\omega_t$ with $0\le\alpha_t\le1/2$, and require $v_t\ge2K u_t$. This gives a nonstationary family allowing creation of new strata and arbitrarily rare old strata. Its control data are known; no algorithm reads an unknown parameter to implement this refresh law. An arbitrary report-dependent stopping rule does not follow from the deterministic-time bound: stopping can select large error. Only an independently verified stopped analogue of \eqref{eq:recharge}, or the bounded common-law comparison of Section~\ref{sec:transfer}, authorizes that extension.

\subsection{Finite-resource and task consequences}
Suppose a certified implementation has the full resource vector of Section~\ref{sec:transfer}; an auxiliary causal simulator has at most $R$ states, a controller has $D_c$ and a retained phase register has $Q$. A constructive total-state budget $M_{\rm tot}$ permits the preceding theorem with $M=\lfloor M_{\rm tot}/(RD_cQ)\rfloor$, provided $M\ge2$ and the declared workspace, program, precision and time bounds are feasible. This is an upper construction, not a product lower. The converse in Section~\ref{sec:soft-state} deals with information that really survives a cut.

For a common ideal baseline $b_*\le B_N$, the exact theorem gives
\[
 R_\on(N,M)-b_*\asymp (B_N-b_*)+\mathcal Q_N(M)
\]
when $V$ is uniformly bounded. The acquisition term $B_N-b_*$ is evaluated from the same raw kernel, not assigned a universal exponent. If the report continuity bound of \eqref{eq:report} holds for the compared controller, a prefix-wise version of \eqref{eq:terminal-conditional} gives the complete report-law allowance \eqref{eq:pathallowance}. Additional causal deficiencies enter the common bounded-loss risk additively. A sharper interaction between actual acquisition and finite input occurs in Theorem~\ref{thm:single-probe}, rather than being inferred from this additive upper bound.

\section{Finite anisotropic acquired geometry}\label{sec:main}
We retain a complementary finite-chart transfer theorem, with an explicit stationary kernel-derived route to its causal hypothesis. Its intersecting-stratum extension does not require the scale-separated countable partition used in Theorem~\ref{thm:bridge}. Its assumptions consist of experimental certificates, not an assumed optimal quantization curve. Throughout, $J$ and $d_*$ are finite and $D_0>0$. The geometric constants $c,C$ depend only on $J,d_*,b,B,C_g,L,C_0,D_0$ and, in the separated case, the dimensionless separation ratios $D_0/\Delta$ and $LD_0/\Delta$. They do not depend on $N,m$, the positive widths or acquired weights. Transport enters separately through the displayed $\Gamma_N$. For a parameterized family, the following certificates, programs and local error bounds must hold uniformly over the declared family; parameters unknown to the machine cannot occur in its instructions.

\subsection{Geometric certificates extracted from histories}
Let $(D,d)$ be a compact valid predictive-state domain, with $\diam D\le D_0$, containing the states under consideration and all representatives used by the algorithm. This is a \emph{valid envelope}, not a claim about the support of the acquired law. Updates on $D$ must have the declared physical meaning or a specified conditional version. Observation-dependent interface components are included and charged. For a fixed exogenous protocol clock, the state fibres at successive times may be identified with $D$; the clock remains a declared public argument of the updates, not an observation-dependent memory channel. Suppose the task map satisfies
\begin{equation}\label{eq:task-lip}
 \norm{p_N(s)-p_N(s')}\le Ld(s,s').
\end{equation}
For each $N$, disjoint measurable events $E_{N,j}$ partition the actual history law up to a null set. Write $w_{N,j}=P^\beta(E_{N,j})$. Zero-weight components can be omitted. The following conditions are checked by integrating the raw preparation and kernels, not by choosing a convenient auxiliary geometric measure.

\smallskip\noindent\textbf{(G) Actual mass and global geometry.}
There are valid compact component domains $D_j\subset D$ covering $D$, with dimensions $0\le d_j\le d_*$, and maps
\[
 \phi_j:Q_j\longrightarrow D_j,
 \qquad Q_j=\prod_{k=1}^{d_j}[-a_{j,k}/2,a_{j,k}/2],
 \quad a_{j,1}\ge\cdots\ge a_{j,d_j}>0.
\]
We normalize the coordinates so that $a_{j,1}\le D_0$. Zero widths are removed, not inverted. Dimension zero means a singleton. The maps are onto, and, for constants $b,B>0$,
\begin{equation}\label{eq:chart}
 b\norm{x-x'}\le\norm{p_N(\phi_jx)-p_N(\phi_jx')}\le Ld(\phi_jx,\phi_jx')
 \le LB\norm{x-x'}.
\end{equation}
Conditioned on $E_{N,j}$, the state law is $\phi_{j\#}(f_{N,j}\,dx)$, where
\begin{equation}\label{eq:mass}
 0\le f_{N,j}(x)\le C_g/|Q_j|,\qquad \int_{Q_j}f_{N,j}(x)\,dx=1.
\end{equation}
For dimension zero this means unit mass. In particular these are actual acquired masses; $f$ need not be positive everywhere in the envelope. For $J>1$, the compact score images $p_N(D_j)$ have pairwise distance at least $\Delta>0$. Constants below may depend on $D_0/\Delta$ and $LD_0/\Delta$. The components can be curved; no differentiable structure, covariance, or Fisher matrix is prescribed. The same component envelope and grids must serve all intermediate updates of the chosen implementation. A certificate depending on $N$ is permissible, but a uniform-in-$N$ conclusion requires uniform constants and one stated implementation for each budget.

Set $\psi_j(0)=D_0^2$ and, for integers $k\ge1$, define
\begin{equation}\label{eq:psi}
 \psi_j(k)=
 \begin{cases}
 \displaystyle\max_{1\le l\le d_j}
 \left(\frac{\prod_{i=1}^l a_{j,i}}{k}\right)^{2/l},&d_j>0,\\
 0,&d_j=0.
 \end{cases}
 \qquad
 \Psi_N(m)=\min_{\substack{k_j\in\N\cup\{0\}\\1\le\sum_j k_j\le m}}
 \sum_j w_{N,j}\psi_j(k_j).
\end{equation}
There are finitely many allocations, so the minimum is attained. The zero-allocation penalty retains the cost of discarding an entire acquired component. In the case $J=1$ it is never needed. Formula~\eqref{eq:psi} is explicit in widths and probabilities obtained from the experiment; it is not a variational definition of optimal distortion.

All objects in \textup{(G)} may depend on the selected horizon: more fully they are \[D^{(N)},\quad D_j^{(N)},\quad\phi_j^{(N)},\quad Q_j^{(N)},\quad a_{j,k}^{(N)},\quad p_N,\quad d_j^{(N)}.\] Within a theorem at fixed $N$ we suppress these superscripts. A uniform assertion ranges over exactly those horizons with the displayed constants uniform. For that horizon, the codebook covers every representative state needed at intermediate times, not just the terminal support. Actual intermediate laws may be different. Components of zero acquired weight may be removed only when their transport contributions also vanish. Distinct component labels are counted before duplicate physical representatives are merged; merging never increases the budget.

\subsection{Causal certificates}
For every allocation $(k_j)$, finite grids in the domains with $k_j>0$ and a fixed Borel tie rule define a codebook of at most $m$ representatives. A discarded component is encoded into an existing representative. Write $r_j=\sqrt{\psi_j(k_j)}$. The geometric construction in Lemma~\ref{lem:geometry} gives component covering errors at most $C_0r_j$. The following additional certificate is the causal content.

\smallskip\noindent\textbf{(I) Innovation-localized implementation.}
There is a specified parameter-independent update of the representative label, executable from its preceding label, the current command/report, and the declared finite-precision primitives. Along every admitted common command/report path, the exact and representative states satisfy
\begin{align}
 e_0&\le C_0\sum_j B_{0,j}r_j+u_0,\notag\\
 e_{t+1}&\le \gamma_{t+1}e_t+C_0\sum_j B_{t+1,j}r_j+u_{t+1},
 \qquad e_t=d(s_t,\widehat s_t).\label{eq:innovation}
\end{align}
The nonnegative measurable coefficients $\gamma_t,B_{t,j}$ are determined by the raw experiment and specified update bounds, independently of the grid allocation and precision budget. They may use latent marks for analysis; the implemented machine must not read these marks. The nonnegative $u_t$ bound calibration and numerical errors in the same metric. One sufficient verification is a same-input Lipschitz bound $\gamma_t$ for the exact update and a rounding error bounded by $C_0\sum_jB_{t,j}r_j$. An exact hold or grid-preserving transformation is executed without rounding, and then has $B_{t,j}=0$. An informative overwrite has $\gamma_t=0$. These checks concern the actual one-step kernel and the given algorithm, not its optimality.

Define empty products as one and set
\begin{equation}\label{eq:transport}
 Z_{N,j}=\sum_{t=0}^N B_{t,j}\prod_{i=t+1}^N \gamma_i,
 \qquad U_N=\sum_{t=0}^N u_t\prod_{i=t+1}^N \gamma_i.
\end{equation}

\smallskip\noindent\textbf{(T) Acquisition-weighted transport.}
For all nonnegative real vectors $z=(z_j)$,
\begin{equation}\label{eq:energy}
 \E_{P^\beta}\left(\sum_j Z_{N,j}z_j\right)^2
 \le \Gamma_N\sum_j w_{N,j}z_j^2.
\end{equation}
This is a moment certificate on a positive transport process specified before choosing the memory budget. For example, the positive semidefinite inequality $\E[Z_NZ_N^{\mathsf T}]\preceq\Gamma_N\operatorname{diag}(w_N)$ suffices. The weaker inequality on the nonnegative cone in \eqref{eq:energy} is all that is used. When a component has zero acquired mass, its transport contribution must vanish almost surely. Verifying this point prevents unacquired geometry from entering an average risk law.

\smallskip\noindent\textbf{(C) Report continuity and effective primitives.}
On common continuation interfaces the conditional report kernels obey
\begin{equation}\label{eq:report}
 \norm{g_t(\cdot\mid s,a)-g_t(\cdot\mid s',a)}_\TV\le\kappa_t d(s,s').
\end{equation}
A calibrated numerical implementation may additionally incur a same-state report error $\zeta_t$. The grid descriptions, update routines, readouts, and finite allocation search in \eqref{eq:psi} are specified by finite programs with declared calibration/numerical input access. Exact algebraic width and weight data suffice for the finite allocation comparisons. A certified constant-factor allocation is also sufficient, changing only the upper constant; arbitrary comparisons of unspecified real oracles are not assumed effective. A metric covering argument alone does not prove bit computability; this effective certificate is required for the word ``implementable.'' The finite algorithms in Section~\ref{sec:realizations} supply it explicitly. A terminal readout approximation contributes at most $v_N^{\rm num}$ in Euclidean score norm.

\smallskip\noindent\textbf{(K) A structural alternative to (T).}
The local innovation coefficients have the marked raw-kernel factor of Section~\ref{sec:kernel}, with bounded forcing, an invariant environment law and a dominating preparation, and satisfy the one-step Lyapunov inequality \eqref{eq:kernel-drift}. The seed satisfies the initial moment bounds of Theorem~\ref{thm:kernel}. Construct $H$ by the positive resolvents \eqref{eq:resolvent-profiles}, and $\Gamma_N^{\rm ker}$ by \eqref{eq:resolvent-gamma}, with the stated zero-weight proviso. These are kernel computations, not hypotheses on an optimal risk or an unrolled forced process.

\begin{theorem}[Causal acquired geometry and kernel-derived resource transfer]\label{thm:main}
Fix the raw experiment, preparation, exploration, task, and horizon $N$. Under the predictive realization certificate and \textup{(G), (I), (C)}, and either \textup{(T)} or \textup{(K)}, there are $c,C>0$ with the dependence specified above such that, for every integer $m\ge1$, taking $\Gamma_N=\Gamma_N^{\rm ker}$ under \textup{(K)},
\begin{equation}\label{eq:main}
 B_N+c\Psi_N(m)\le R_\cp(N,m)\le R_\on(N,m)
 \le B_N+\left(C\sqrt{\Gamma_N\Psi_N(m)}
             +L\norm{U_N}_{L^2}+v_N^{\rm num}\right)^2.
\end{equation}
For exact comparable allocation data, the online upper is attained as an inequality by the specified finite representative machine with a minimizing allocation. If only an $A_{\rm alloc}$-approximate allocation is computable, replace $C$ by $C\sqrt{A_{\rm alloc}}$ in the upper bound. The exact Borel model and the finite-precision program are distinct: when exact Borel primitives and exact calibration are available, use $u_t=v_N^{\rm num}=0$; the latter has the displayed nonzero error allowances. If $\Psi_N(m)=0$, matching means both exact excess risks are zero, not a ratio of zero quantities. In particular, with exact arithmetic and calibration, and $\sup_N\Gamma_N<\infty$,
\begin{equation}\label{eq:matching}
 R_\cp(N,m)-B_N\asymp R_\on(N,m)-B_N\asymp\Psi_N(m)
\end{equation}
uniformly over the certified horizons, positive widths, and acquisition weights. Widths may vanish and the acquired rank may change, provided the surviving-coordinate certificates have the same constants. There is no inverse of a vanished width in the conclusion.

For a common bounded decision loss of range at most one, the same implementation yields a finite-horizon report-law allowance
\begin{equation}\label{eq:pathallowance}
 \min\left\{1,\sum_{t<N}\E[\kappa_t e_t+\zeta_t]\right\},
\end{equation}
whenever the local implementation and report certificates hold for the controller actually compared. An estimate using a moment certificate for a different controller is not implied. Uniform comparison over a controller class requires all prefix certificates used in the bound to be uniform over that class. A common visible stopping time $\sigma\le N$ inserts $\mathbf1_{\{t<\sigma\}}$ in this sum. The assertion is for the given horizon or the stated bounded stopping time, not an infinite-horizon path law. Additional causal simulators compose this allowance with their certified defects, and compose their resource maps as in Theorem~\ref{thm:morphism}. None of these assertions compares two different full-history oracle baselines.
\end{theorem}

The acquisition term can be made explicit without inventing a universal power of $N$. For a specified common ideal baseline $b_*$ with $B_N\ge b_*$, put $I_N=B_N-b_*$. Then the exact implementation has
\begin{equation}\label{eq:phi}
 R_\on(N,m)-b_*\asymp I_N+\Psi_N(m).
\end{equation}
The quantity $I_N$ must itself be computed or bounded from the actual acquisition kernel. Section~\ref{sec:sensor} computes it exactly as a failed-acquisition probability times a target variance. Equation~\eqref{eq:phi} is not a claim that every experiment has the same statistical exponent.

\section{Proofs of the finite-chart estimates}\label{sec:proofs}\label{sec:proof}
\begin{lemma}[Checkpoint projection]\label{lem:projection}
For the task \eqref{eq:score},
\begin{equation}\label{eq:projection}
 R_\cp(N,m)=B_N+
 \inf_{c_1,\ldots,c_m\in\mathcal C}\int\min_l\norm{x-c_l}^2\,\nu_N(dx),
 \qquad \mathcal C=\prod_i[0,\sqrt{v_i}].
\end{equation}
Randomized encoders and decoders do not improve this infimum.
\end{lemma}
\begin{proof}
For each executed probe, condition on $H_N$ and on any independent encoder randomization. The cross term in
\[
 (T_i-a_i)^2=(T_i-\E[T_i\mid H_N])^2+
 (\E[T_i\mid H_N]-a_i)^2+
 2(T_i-\E[T_i\mid H_N])(\E[T_i\mid H_N]-a_i)
\]
has conditional expectation zero. Sum with weights $v_i$. For a retained label, replace a randomized decoder by its conditional mean; convexity cannot increase the risk. For fixed centers, replace a randomized label by a nearest center with the first-index tie rule. This is measurable. Conversely, this rule is a legal checkpoint encoder and attains its codebook distortion. Independent shared random seeds only average the costs of such deterministic codebooks, so cannot lower their infimum. Projection of a center to the box $\mathcal C$ is nonexpansive relative to all points of $\mathcal C$. Compactness gives attainment of the codebook minimum, although only its infimum is needed.
\end{proof}

\begin{lemma}[Actual anisotropic mass and allocation]\label{lem:geometry}
Under \textup{(G)}, the distortion in \eqref{eq:projection} lies between $c\Psi_N(m)$ and $C\Psi_N(m)$ for every $m\ge1$. Each component with $k\ge1$ has a $k$-point valid-state cover of radius at most $B\sqrt{d_*\psi_j(k)}$. The conclusions hold with the same constants when any collection of widths is removed at zero.
\end{lemma}
\begin{proof}
Fix a component of positive dimension $d$ and write
$h=\sqrt{\psi_j(k)}$. Put $R=2h$. For coordinates with $a_i>R$, divide the interval into $n_i=\lceil a_i/R\rceil$ equal pieces; for the other coordinates take one piece. If there are $l$ active coordinates, they are the first $l$, and
\[
 \prod_i n_i\le\prod_{i=1}^l\frac{2a_i}{R}
 =\frac{\prod_{i=1}^la_i}{h^l}\le k.
\]
The assertion also holds when $l=0$. Every coordinate interval has length at most $R$. Its midpoint product grid has Euclidean covering radius at most $\sqrt d\,h$, and \eqref{eq:chart} gives the valid-state radius. A point component takes one label and has zero radius. An omitted component can use any retained representative with error at most $D_0$. Minimizing the allocation proves the upper bound, including the cases $m<J$ and atoms of arbitrary weight.

For the lower bound, a score ball of radius $r$ has preimage of Euclidean diameter at most $2r/b$. If it meets the component, its projection on any one coordinate is contained in an interval of length at most $4r/b$. Integrating the actual density in \eqref{eq:mass} over the remaining coordinates gives, for $1\le l\le d$,
\begin{equation}\label{eq:smallball}
 P^\beta\{p_N(s_N)\in B(c,r)\mid E_{N,j}\}
 \le C_g\frac{(4r/b)^l}{\prod_{i=1}^la_i}.
\end{equation}
The center $c$ need not lie on the image. This is a bound on the actual experiment, not on a fictitious uniform distribution.

If $J=1$, all $m$ centers can be used in the union bound. For a fixed $l$, choose $r=c_l(\prod_{i=1}^la_i/m)^{1/l}$ so that the union has probability at most $1/2$. The expected squared error is at least $r^2/2$. Maximize over $l$.

If $J>1$, assign each decoder center to a component whose score image is at distance less than $\Delta/3$, if such a component exists. Separation makes this assignment unique. Let $k_j$ be the number assigned. For $r<\Delta/3$, unassigned or differently assigned centers do not cover any point of component $j$. When $k_j\ge1$, use \eqref{eq:smallball} and the preceding radius, capped at $\Delta/4$. Since $\psi_j(k_j)\le D_0^2$, the cap changes the lower-bound constant only by a factor depending on $D_0/\Delta$, $b$, and $C_g$. For $k_j=0$, the component's distortion is at least $\Delta^2/9$, hence at least a fixed multiple of $\psi_j(0)$. At a point with $k_j\ge1$ the required lower bound is zero. Multiply each component bound by its actual weight and sum. We have $\sum_jk_j\le m$. If this sum is zero, give one component a label: its $\psi_j$ cannot increase, since $a_{j,1}\le D_0$. Thus the sum is bounded below by a constant times the minimum in \eqref{eq:psi}. This proves the lower bound for arbitrary decoder centers and all budgets.

Only surviving coordinates were used. No constant in the covering or small-ball argument contains $a_i^{-1}$ except within the displayed physical volume, which is eliminated in the final scale. This proves the rank-uniform assertion.
\end{proof}

\begin{lemma}[Innovation transport]\label{lem:transport}
Under \textup{(I)}, for every admitted path,
\begin{equation}\label{eq:unrolled}
 e_N\le C_0\sum_jZ_{N,j}r_j+U_N.
\end{equation}
Consequently \textup{(T)} implies
\[
 \norm{e_N}_{L^2}\le C_0\sqrt{\Gamma_N\sum_jw_{N,j}r_j^2}
                      +\norm{U_N}_{L^2}.
\]
\end{lemma}
\begin{proof}
Induct in \eqref{eq:innovation}, distributing each nonnegative multiplicative coefficient. This gives precisely \eqref{eq:transport} and \eqref{eq:unrolled}. Minkowski's inequality and \eqref{eq:energy}, evaluated at $z_j=r_j$, give the second assertion. No independence of the coefficients is used.
\end{proof}

\begin{proof}[Proof of Theorem~\ref{thm:main}]
Under \textup{(K)}, Theorem~\ref{thm:kernel}, whose proof uses only the raw marked kernel and positive resolvents, first supplies \textup{(T)}. This invocation is not circular. Lemma~\ref{lem:projection} and Lemma~\ref{lem:geometry} give the lower bound for every checkpoint encoder. The terminal output of any online encoder is a checkpoint encoder, so $R_\cp\le R_\on$. Choose a minimizing allocation in \eqref{eq:psi} and its valid grids. With only a certified $A_{\rm alloc}$-approximate allocation, use it instead; its weighted squared radii increase by at most $A_{\rm alloc}$. Execute the label update in \textup{(I)}. By \eqref{eq:task-lip}, its score error has $L^2$ norm at most $L\norm{e_N}_{L^2}+v_N^{\rm num}$. The conditional projection identity is valid for this online encoder as well. Apply Lemma~\ref{lem:transport}, and absorb $LC_0$ into $C$. This proves \eqref{eq:main}. Its exact-arithmetic specialization gives \eqref{eq:matching}; adding the same nonnegative $I_N$ yields \eqref{eq:phi}.

For path comparison, run the exact and representative report experiments with the same controller at identical visible prefixes. At such a prefix the action kernels coincide, the retained label is the same function of the prefix, and the next-report kernels differ by at most $\kappa_te_t+\zeta_t$. Maximal coupling of the next report is possible on standard Borel spaces. The common-prefix subprobability law is dominated by the exact prefix law. Hence the probability of a first disagreement at step $t+1$ is at most $\E_{P^\beta}[\kappa_te_t+\zeta_t]$, or the corresponding expectation for the controller actually being compared. Summing proves \eqref{eq:pathallowance}. Append the common label update and resource readout; these do not increase the discrepancy. After the common stopping decision, both processes are absorbed, which inserts the active-step indicator. To use the numerical bound from \textup{(T)} for a different controller, that controller must have its own certificate; the fixed-$\beta$ certificate is not automatically uniform over policies. Composition and common-loss transfer are proved in Section~\ref{sec:transfer}.
\end{proof}

\subsection{How to verify the transport certificate}
\begin{proposition}[Three kernel-level mechanisms]\label{prop:mechanisms}
The following are sufficient verifications of \textup{(T)}.

\textup{(a)} With one component of weight one, $\gamma_t\le\rho<1$, $B_{t,1}\le1$, and $B_{0,1}\le1$, one may take $\Gamma_N\le(1-\rho)^{-2}$.

\textup{(b)} Suppose each update is either an exact hold with $\gamma_t=1$, $B_{t,1}=0$, or an active step with $\gamma_t\le\rho<1$, $B_{t,1}\le1$. Then the same constant works, without a lower bound on the probability or frequency of active steps. Here $w_{N,1}=1$; an initial rounding is permitted.

\textup{(c)} Suppose the process starts at one known unresolved point, waits there, and at a first successful acquisition moves into component $j$, thereafter holding. The update on a successful report overwrites the old prediction, including its conditional version at an otherwise impossible report. Let $B_{0,0}=1$, and let $B_{t,j}=1$ only at the acquisition into $j$. Then
\[
 Z_{N,0}=\mathbf1_{\{\text{no acquisition by }N\}},\qquad
 Z_{N,j}=\mathbf1_{\{\text{acquired component }j\text{ by }N\}}.
\]
For the corresponding actual component weights, \eqref{eq:energy} is an equality with $\Gamma_N=1$.
\end{proposition}
\begin{proof}
In (a), sum a geometric series. In (b), delete the exact holds from the error recurrence, not from the raw acquisition or time account. The resulting pathwise series is the same geometric series indexed by active events. In (c), failure and post-acquisition holds have $A=1,B=0$, while the acquisition has $A=0$ and the appropriate $B=1$. Thus the initial contribution is killed by the first acquisition and the new contribution persists. The displayed indicators are disjoint and exhaustive. Squaring their weighted sum and taking expectation gives the assertion.
\end{proof}

Mechanism (b) permits arbitrarily long metastable waiting periods. It is different from replacing the raw horizon by an uncharged number of successful reports. Mechanism (c) additionally preserves the small acquired mass in the upper bound, rather than paying the full geometric cost on paths that never acquire it. Recurrent nonuniform stability is derived from the one-step kernel condition in Theorem~\ref{thm:kernel}; it is not inferred from a negative average logarithmic Lipschitz exponent alone. The relation to average contraction theory \cite{diaconis} is thus precise: here the required object is a forced, acquisition-weighted second moment, not merely convergence of unforced iterates.

\begin{corollary}[Intersecting strata]\label{cor:crossing}
Remove the separation assumption from \textup{(G)}, allow the score images of the $J$ component charts to intersect, and let the raw preparation carry a latent component mark with probabilities $w_{N,j}$. The mark need not be available to the algorithm or retained by the quotient. Suppose the conditional laws and the other certificates remain valid, with chart rounding certified conditionally on the raw mark. Then \eqref{eq:main} and \eqref{eq:matching} hold for every $m\ge2J$, with constants also depending on $J$, and without any inverse distance between strata.
\end{corollary}
\begin{proof}
The global upper cover is still the union of the component grids. Encode into a nearest valid representative in $d$; its error on component $j$ is no greater than that component grid's error. Duplicate representatives may be merged. The algorithm does not need the latent mark.

For the lower bound apply \eqref{eq:smallball} with all $m$ decoder centers to each positive-dimensional conditional law. This gives distortion at least $c\sum_jw_{N,j}\psi_j(m)$; point components contribute zero. Set $k=\lfloor m/J\rfloor\ge1$. Directly from \eqref{eq:psi}, for $m\ge k\ge1$,
\[
 \psi_j(k)\le(m/k)^2\psi_j(m)\le(2J)^2\psi_j(m).
\]
The allocation of $k$ labels to every component is feasible. Thus $\Psi_N(m)\le(2J)^2\sum_jw_{N,j}\psi_j(m)$. Combining the last two bounds proves the claimed lower, and the remaining causal proof is unchanged. In particular neither smoothness at an intersection nor identifiability of the latent mark was used.
\end{proof}

\section{A kernel theorem for recurrent nonuniform transport}\label{sec:kernel}
The moment condition in the transfer theorem is useful only insofar as it can be derived without solving the approximation problem itself. We now derive it from a positive operator attached to a marked factor of the raw experiment. The factor may be used only for analysis. It need not be observed, and it is never supplied to the encoder.

\subsection{One-step data and positive resolvents}
Let $\Xi$ be standard Borel, and suppose that the coefficients of the innovation recurrence are generated by a time-homogeneous marked kernel
\begin{equation}\label{eq:marked-kernel}
 P(d\xi',d\gamma,d\mathbf b\mid\xi),\qquad
 \gamma\ge0,\quad \mathbf b\in[0,\infty)^J,\quad \sum_j b_j\le B_*.
\end{equation}
Thus, conditionally on $\xi_t$, the next environment and marks are independent of all previous marks and of the seed. This assertion is checked from the raw preparation, the chosen exploration and the one-step update bounds. It excludes coefficients chosen after seeing the grid allocation. Controller state may have to be included to make the factor Markov; its implementation cost is not thereby removed.

Assume that the environment marginal has an invariant probability $\mu$. Disintegrate the stationary marked edge measure in the reverse direction:
\[
 \mu(d\xi)P(d\xi',d\gamma,d\mathbf b\mid\xi)
   =\mu(d\xi')\overleftarrow P(d\xi,d\gamma,d\mathbf b\mid\xi').
\]
All the following identities are statements about measurable versions $\mu$-almost everywhere. Define, for nonnegative measurable $f$,
\begin{align}
 (K_p f)(\xi')&=\int\gamma^p f(\xi)\,d\overleftarrow P,
       &&p=1,2,\label{eq:positive-operators}\\
 f_j(\xi')&=\int b_j\,d\overleftarrow P,
 &g_{jk}(\xi')&=\int b_jb_k\,d\overleftarrow P,\notag\\
 (D_jf)(\xi')&=\int\gamma b_j f(\xi)\,d\overleftarrow P.\notag
\end{align}
The structural hypothesis is the one-step inequality
\begin{equation}\label{eq:kernel-drift}
 1\le v\le v_+<\infty,\qquad K_2v\le\lambda v,
 \qquad 0\le\lambda<1.
\end{equation}
Here $v$ is a measurable Lyapunov weight. Neither $\gamma<1$ on each edge nor a bound on a previously unrolled error is assumed. The operators depend only on the experimental factor and the selected local update/rounding bounds.

For a positive weight $u$, write $\|f\|_u=\operatorname*{ess\,sup}|f|/u$. Cauchy--Schwarz gives $K_1\sqrt v\le\sqrt\lambda\sqrt v$. Hence the positive Neumann resolvents $\mathcal S_1=(I-K_1)^{-1}$ on the $\sqrt v$-weighted space and $\mathcal S_2=(I-K_2)^{-1}$ on the $v$-weighted space exist. Set
\begin{align}
 h_j&=\mathcal S_1 f_j,\label{eq:resolvent-profiles}\\
 V_{jk}&=\mathcal S_2\bigl(g_{jk}+D_jh_k+D_kh_j\bigr),
 &H_{jk}&=\int V_{jk}\,d\mu.\notag
\end{align}
The profiles are finite. Indeed $h_j$ is bounded by a constant times $\sqrt v$, and Cauchy--Schwarz bounds $D_jh_k$ by a constant times $\sqrt v$, hence by a constant times $v$. The constants depend only on $B_*,v_+,1-\lambda$ and $J$, not on a memory allocation.

\begin{theorem}[Forced moment resolvent from a raw-kernel factor]\label{thm:kernel}
Suppose \eqref{eq:marked-kernel}--\eqref{eq:kernel-drift} hold. Let
\[
 Z_{0,j}=0,\qquad Z_{t+1,j}=\gamma_{t+1}Z_{t,j}+b_{t+1,j}.
\]
Under the stationary environment preparation, for every $N\ge0$ and every $z\in[0,\infty)^J$,
\begin{equation}\label{eq:kernel-energy}
 \E\left(\sum_j z_jZ_{N,j}\right)^2\le z^{\mathsf T}Hz.
\end{equation}
The matrix $H$ is symmetric positive semidefinite and has nonnegative entries. For the seeded version define separately $\widetilde Z_{0,j}=S_j\ge0$ and $\widetilde Z_{t+1,j}=\gamma_{t+1}\widetilde Z_{t,j}+b_{t+1,j}$. The same inequality holds for $\widetilde Z$ if, conditionally on $\xi_0$, $\E S_j\le h_j$ and $\E(S_jS_k)\le V_{jk}$, entry by entry. In the innovation notation of \eqref{eq:transport}, $S_j$ is precisely the time-zero forcing $B_{0,j}$; subsequent $b_{t,j}$ represent $B_{t,j}$. These indices identify the seed rather than treating it as an uncharged extra call. If the actual initial environment law has density at most $D_*$ with respect to $\mu$, and uses the same conditional seed rule, the right side is multiplied by $D_*$.

Consequently, for the actual acquired weights $w_{N,j}$, condition \textup{(T)} follows with
\begin{equation}\label{eq:resolvent-gamma}
 \Gamma_N^{\rm ker}
   =D_*\left\|\operatorname{diag}(w_N)^{-1/2}
                    H\operatorname{diag}(w_N)^{-1/2}\right\|_{\rm op},
\end{equation}
provided every zero-weight row of $H$ is zero; the inverse is taken only on positive-weight coordinates. If this proviso fails, this stationary majorant does not certify that horizon. A finite-horizon profile or another mechanism must then be used. In particular, for $J=1$, $w_{N,1}=1$, the one-step drift alone gives a uniform finite transport constant, with no assumption of uniform contraction of the active updates.
\end{theorem}
\begin{proof}
First start the environment at $\mu$ and the forcing process at zero. Define conditional moment profiles
\[
 m_{t,j}(\xi)=\E[Z_{t,j}\mid\xi_t=\xi],\qquad
 C_{t,jk}(\xi)=\E[Z_{t,j}Z_{t,k}\mid\xi_t=\xi].
\]
The Markov-factor hypothesis ensures conditional independence of future edge marks and the previous forcing history given the previous environment. Conditioning on the stationary incoming edge therefore gives the exact recursions
\begin{align}
 m_{t+1,j}&=K_1m_{t,j}+f_j,\label{eq:profile-recursion}\\
 C_{t+1,jk}&=K_2C_{t,jk}+D_jm_{t,k}+D_km_{t,j}+g_{jk}.\notag
\end{align}
Truncation of the marks justifies these identities before integrability is known; monotone convergence then removes the truncation. Alternatively, induction using the finite bounds below proves integrability at each finite time.

Positivity and the identities $h_j=K_1h_j+f_j$ and
$V_{jk}=K_2V_{jk}+D_jh_k+D_kh_j+g_{jk}$ imply
$0\le m_{t,j}\le h_j$ and $0\le C_{t,jk}\le V_{jk}$ by induction. In fact $m_t$ and $C_t$ increase from zero. The first recursion converges in the $\sqrt v$ norm to $h$. Subtracting the second recursion from its fixed-point identity, or summing its geometric convolution, proves convergence in the $v$ norm to $V$. To see the latter directly, split the convolution at a fixed lag: its long-lag part is bounded by a geometric $\lambda$-tail, and the finitely many remaining terms converge because $m_t\to h$.

For each $t$ the conditional matrix $C_t(\xi)$ is positive semidefinite. Entrywise convergence in a finite matrix preserves this property outside one common null set. Thus $V(\xi)$ and its integral $H$ are positive semidefinite. All entries are nonnegative. Integrating the entrywise upper bounds and multiplying by $z_jz_k\ge0$ proves \eqref{eq:kernel-energy}. This last nonnegative-cone bound is the one needed for covering radii; it is not an entrywise-to-Loewner inference.

For a seed satisfying the stated conditional inequalities, the same induction in \eqref{eq:profile-recursion} applies, without the monotone-convergence assertion for that seed. An initial density bounded by $D_*$ dominates the entire path-and-seed law by $D_*$, since all subsequent kernels agree. Integrating any nonnegative path functional proves the domination claim. Finally apply the operator norm bound to $\operatorname{diag}(w_N)^{1/2}z$. A zero diagonal of a positive semidefinite matrix forces its row to be zero; the explicit proviso prevents assigning positive transport energy to a component of zero acquired probability.
\end{proof}

This theorem replaces a path-length family of forced-moment inequalities by one-step kernel integration, a Lyapunov inequality, and two positive resolvents. The comparison with actual acquired weights remains visible in \eqref{eq:resolvent-gamma}: small mass is not made free by a stability theorem. Uniform matching across a family means that this explicitly computed normalized matrix is uniformly bounded. For one acquired chart no further mass-normalization condition is needed. For a finite environmental chain, all these objects are finite matrices and vectors, independent of $N$ and $m$.

\begin{corollary}[Finite-state check and recurrent expansion]\label{cor:finite-kernel}
For a finite environment, $\rho(K_2)<1$ is equivalent to the existence of a strictly positive vector $v$ and $\lambda<1$ satisfying $K_2v\le\lambda v$. In particular consider a two-state chain, with states $C,E$, transition matrix
\[
 P=\begin{pmatrix}1-p&p\\1&0\end{pmatrix},\qquad 0<p\le1,
\]
and gain $1/4$ on arrival at $C$, $2$ on arrival at $E$. Under its invariant law, the reversed chain has the same transition matrix and
\[
 K_2=\begin{pmatrix}(1-p)/16&p/16\\4&0\end{pmatrix},
 \qquad K_2\binom{1}{8}\le\tfrac12\binom{1}{8}.
\]
Bounded forcing therefore satisfies Theorem~\ref{thm:kernel}, although expansions recur and the conditional squared gain on arrival at $E$ is four.
\end{corollary}
\begin{proof}
A positive vector inequality makes the weighted sup norm of $K_2$ at most $\lambda$, so its spectral radius is less than one. Conversely, when $\rho(K_2)<1$, let $v=(I-K_2)^{-1}\mathbf1\ge\mathbf1$. Then $K_2v=v-\mathbf1\le(1-1/\max_i v_i)v$; the zero matrix case is included. The displayed chain has invariant probabilities $(1,p)/(1+p)$ and satisfies detailed balance. Multiplication by the squared incoming gain gives the stated matrix. Its first row applied to $(1,8)$ is $(1+7p)/16\le1/2$ and its second is four. These calculations prove the example. It is an operator witness; a complete raw experiment and its acquired geometry are verified separately below.
\end{proof}

\begin{proposition}[Exact suspension preserves the forced profiles]\label{prop:suspension}
Replace each transition of \eqref{eq:marked-kernel} independently by an exact hold with probability $1-\theta$, $0<\theta\le1$. A hold has unchanged environment, gain one and zero forcing. Keep every hold in the raw-call and physical-time accounts. The invariant environment law remains $\mu$, and the profiles $h,V,H$ in \eqref{eq:resolvent-profiles} are unchanged. Thus the bound \eqref{eq:resolvent-gamma} has no inverse power of $\theta$ apart from any change in the independently computed acquired weights. At $\theta=0$ the forcing process remains its seed, so the same upper bound holds for the seeds admitted in Theorem~\ref{thm:kernel}.
\end{proposition}
\begin{proof}
Reversing the stationary held kernel gives the same mixture. Therefore
\[
 K_p^\theta=(1-\theta)I+\theta K_p,
 \quad f_j^\theta=\theta f_j,\quad
 g_{jk}^\theta=\theta g_{jk},\quad D_j^\theta=\theta D_j.
\]
The drift constant becomes $1-\theta(1-\lambda)<1$. Cancelling $\theta$ in the two fixed-point equations proves $h^\theta=h$ and $V^\theta=V$. This cancellation is justified by the positive resolvent's uniqueness, not by replacing physical time with a random stopping time. At zero activity there are no new innovations and direct retention proves the assertion. The independence and state-independent value of $\theta$ are part of this proposition; arbitrary adaptive suspensions require their own kernel calculation.
\end{proof}

For a scalar independent active gain $G$, with forcing one at every active step, write $a_1=\E G$ and $a_2=\E G^2<1$. The equations give
\begin{equation}\label{eq:scalar-resolvent}
 h=\frac1{1-a_1},\qquad H=\frac{1+a_1}{(1-a_1)(1-a_2)}.
\end{equation}
Indeed $a_1\le\sqrt{a_2}<1$, and substitution into $H=a_2H+2a_1h+1$ proves the formula. It also bounds any seed with first moment at most $h$ and second moment at most $H$. This explicit second-moment formula is a consequence of the general kernel theorem, not a replacement for it. Section~\ref{sec:recurrent} verifies a recurrent continuous-state experiment in this class, including its actual density and finite program.

A negative average logarithmic gain would not replace \eqref{eq:kernel-drift}. For example, $G=1/4$ with probability $9/10$ and $G=4$ with probability $1/10$ has negative mean logarithm but $\E G^2=53/32>1$. Already the squared product multiplying one earlier nonzero innovation has expectation growing geometrically. The theorem controls forced second moments, not only almost-sure contraction of unforced iterates.

\section{Mandatory states and a joint minimax experiment}\label{sec:minimax}
The preceding resource map is constructive. Its converse depends on what future continuations must still be executable. We first establish a general state-separation principle, then exhibit a single experiment in which all the proposed risk terms occur simultaneously.

\subsection{A continuation-separation cut}
A \emph{latched continuation tag} is a variable $C\in\{1,\ldots,L\}$, reported before a memory cut and not reissued afterward. A fresh legal challenge after the cut asks for $C$ with zero error. Equivalently, finitely many challenges ask for its coordinates. This is an exact continuation requirement in the algorithm class, not an extra squared-loss penalty with an unstated weight. The prediction loss is unchanged. We call this the \emph{audited continuation interface}. An external clock cannot supply the tag.

\begin{theorem}[Mandatory-state separation and total-memory resolution]\label{thm:tag}
Let $C$ be uniform on $L$ tags and independent of an acquired score $X$ with law $\nu$. A checkpoint machine has at most $M$ total persistent states and must satisfy the audited continuation interface. Let $Q_\nu(k)$ be its ordinary $k$-center squared distortion, with arbitrary centers. For $M\ge L$, its exact checkpoint excess risk is
\begin{equation}\label{eq:tag-exact}
 \min_{\substack{k_c\in\N,\ k_c\ge1\\\sum_c k_c\le M}}
       \frac1L\sum_{c=1}^L Q_\nu(k_c).
\end{equation}
If $M<L$, the admissible class is empty. These statements include randomized fused implementations with independent shared randomness.

Suppose $\nu$ has one anisotropic acquired chart of dimension at most $d_*$ and the constants of \textup{(G)}. Write $\psi$ for its positive-budget scale in \eqref{eq:psi}. Then, for $M\ge2L$, every such machine has excess at least
\begin{equation}\label{eq:tag-lower}
 c\,\psi(\lfloor M/L\rfloor).
\end{equation}
If the geometric experiment also satisfies the bounded-transport online theorem and the tag remains unchanged, an online product construction with $L\lfloor M/L\rfloor$ states attains the corresponding upper bound, with its full workspace/program/precision/time costs charged. Constants are independent of $M,L$ and positive widths.

In particular, if a simulator key, a controller continuation token and an internal phase token are independent uniform tags of sizes $R_0,D_0',Q_0$, and the future interface separately challenges them, then $L=R_0D_0'Q_0$. The growing product overhead is mandatory even for an implementation that fuses all registers. No assertion of necessity is made for registers lacking these continuation distinctions.
\end{theorem}
\begin{proof}
Fix an independent shared seed. A persistent state reached with positive probability under two different tags cannot support a zero-error answer to the same fresh challenge under both tags. Hence the positive-probability state supports of different tags are disjoint. There are only finitely many tags and states, so one may choose a common full-measure set of seeds on which this holds. Let $k_c$ be the size of the support for tag $c$. Then $k_c\ge1$ and $\sum_c k_c\le M$.

For each tag, the terminal prediction vectors have at most $k_c$ centers. Randomized decoding is replaced by its mean, and randomized encoding cannot improve nearest-center distortion. Independence of $C$ and $X$ therefore gives a conditional excess of at least $Q_\nu(k_c)$. Averaging proves the lower bound in \eqref{eq:tag-exact}. Conversely, use disjoint sets of $k_c$ labels for optimal codebooks under each tag and decode the tag from its set. Compactness gives optimal codebooks in our geometric setting; approximating codebooks suffice in the infimum formulation. This proves equality and infeasibility when $M<L$. Random shared seeds only average costs bounded below by the same minimum.

To prove the uniform scale bound, put $x=M/L\ge2$. At least half the tags have $k_c\le2x$. Lemma~\ref{lem:geometry} and monotonicity give an average distortion at least a fixed multiple of $\psi(\lceil2x\rceil)$. If $k\ge n\ge1$, directly from \eqref{eq:psi},
\[
 \psi(k)\ge(n/k)^2\psi(n).
\]
Since $\lceil2x\rceil/\lfloor x\rfloor\le5$, this proves \eqref{eq:tag-lower}. For the online upper, retain the exact tag and run the certified geometric encoder with $\lfloor M/L\rfloor$ labels. The tag affects neither its acquired law nor its error recurrence. Its product state capacity is at most $M$. Independent separately challenged tags are distinguished exactly by their joint value, proving the last assertion.
\end{proof}

This is a converse for a semantic continuation requirement, not a proof that each factor in every chosen implementation must be present. With an allowed tag-error probability the information argument of Theorem~\ref{thm:soft-state} supplies a quantitative lower instead. Likewise a fixed public phase known independently of observations is not one of these tags.

\subsection{One common robust experiment}
Fix an integer $d\ge1$, widths $0\le a_i\le1/4$, an acquisition probability $s\in[0,1]$, a tag size $L\ge1$, $0\le\delta\le1/4$, $0\le\varepsilon\le1/2$, and an integer $b\ge1$. Zero-width coordinates below are deterministic. Prepare independently a tag $C$ uniform on $L$ values and a vector
\[
 X_i\sim\operatorname{Unif}[1/2-a_i/2,1/2+a_i/2].
\]
The physical parameter is $\lambda=(\sigma,\eta,z)\in\{-1,1\}^2\times\{0,1\}$. It is not supplied to the algorithm. Put $h_b=2^{-b-2}$. The common calibration input for the terminal probability $1/2+\sigma\delta$ is $1/2$; the common $b$-precision input for the terminal probability $1/2+\eta h_b$ is also $1/2$. Both points in the latter pair lie in the same width-$2^{-b}$ rounding cell centered at $1/2$. There is no legal preterminal probe revealing either sign. These are specified information interfaces, not optional inaccurate arithmetic.

The first paid call latches $C$ and supplies these identical calibration/numerical packets. The second reports the bit $z$ with probability $1-2\varepsilon$ and an erasure otherwise. The erasure coin is independent of everything else. The next $n$ raw calls are attempts: each unresolved attempt succeeds with independent probability $s$, reveals $X$ and absorbs; failures and post-success calls are exact holds. Including the one selected terminal probe below, all raw calls count, so $N_{\rm raw}=n+3$. The exploration makes these calls, independently of the encoder. At the memory cut the audited continuation interface may challenge $C$. This challenge is a readout of the retained machine, not an additional physical measurement of the hidden apparatus; its output operations are charged in time.

The soft task independently selects one of four blocks with probability $1/4$. In the geometry block it selects a coordinate uniformly and executes a terminal Bernoulli probe of mean $X_i$. The other blocks execute a terminal Bernoulli probe of mean $1/2+\sigma\delta$, a terminal Bernoulli probe of mean $1/2+\eta h_b$, or the deterministic terminal bit $z$. Only the selected probe is executed. A forecast in $[0,1]$ is made before its outcome. Every block uses squared loss, so the common loss lies in $[0,1]$.

Define the exact-bit target by changing only the second call to an exact report of $z$, with the same costs, all other kernels and the same terminal menu. Admitted bit simulators randomize the sole preterminal bit report and parameter-independent auxiliary reports. No terminal readout can be requested before the target forecast: the terminal command follows that forecast and ends the physical experiment. Thus a simulator cannot query the future answer to manufacture an earlier exact bit. The deficiency of this source--target pair is exactly $\varepsilon$, by Proposition~\ref{prop:floors}; independent auxiliary reports do not change the binary testing lower bound.

The ideal oracle knows $X,\sigma,\eta,z$; its risk is the same for every $\lambda$:
\begin{equation}\label{eq:robust-baseline}
 b_* =\frac14\left\{\frac1d\sum_i\E[X_i(1-X_i)]
           +(1/4-\delta^2)+(1/4-h_b^2)\right\}.
\end{equation}
This is the single baseline subtracted below. Write
\[
 H_n=(1-s)^n,\qquad V_a=\frac1{12d}\sum_i a_i^2,
 \qquad \psi_a(k)=\max_{1\le l\le r}
       \left(\frac{\prod_{i=1}^l a_{(i)}}{k}\right)^{2/l},
\]
where $r$ is the number of positive widths, listed in decreasing order; set $\psi_a=0$ if $r=0$. At $n=0$ use $H_0=1$, also when $s=1$. The calibration and precision packets belong to the source experiment's actual reports. They do not manufacture two distinguishable future states before the missing information has been acquired.

Let $\mathcal C(a,s,L,\delta,b,\varepsilon)$ be exactly this eight-parameter-member family with its stated preparation, report interfaces and task. In the exact-report model, current successful geometric reports may be read as Borel inputs. In the finite-bit model they are accessed through a primitive returning a vector within $2^{-b}$ in each physical coordinate; the machine may not bypass this primitive. The primitive depends only on the current geometric report and independent randomization, not on $\sigma,\eta,z$; it is uniform over the family's parameters. Its error guarantee, and finite rational chart data where needed for execution, are part of this second model.

Using the total budget \eqref{eq:total-budget}, define a single minimax functional
\begin{equation}\label{eq:robust-risk}
 \mathfrak R_{\mathcal C}^{\chi}(\mathbf b)
   =\inf_{A\in\mathcal A_{\mathcal C}^{\chi}(\mathbf b)}
       \sup_{\lambda\in\{-1,1\}^2\times\{0,1\}}
          \bigl(\E_\lambda^A\ell-b_*\bigr).
\end{equation}
The infimum is over parameter-independent machines satisfying the same audited continuation interface and input model. The exact-label model places no finite-bit bound on current Borel comparisons; the finite-bit version below does. In both cases total retained state, and not just a geometric subregister, is constrained by $M$.

\begin{theorem}[Joint acquisition--resolution--precision minimax law]\label{thm:robust}
For the family just defined, let $N_{\rm raw}=n+3\ge3$, $M\ge12L$, and $b_{\rm num}=b$. There are positive constants $c_d,C_d$, depending only on $d$, such that
\begin{equation}\label{eq:robust-phi}
 \Phi=H_nV_a+(1-H_n)\psi_a(\lfloor M/L\rfloor)
                  +\delta^2+2^{-2b}+\varepsilon
\end{equation}
and every nonempty finite-bit or exact-report resource class satisfies
\begin{equation}\label{eq:robust-lower}
 \mathfrak R_{\mathcal C}^{\rm pw}(\mathbf b)\ge c_d\Phi.
\end{equation}
An exact Borel finite-state encoder with at most $M$ states attains an upper bound $C_d\Phi$. The same upper holds for a finite-bit program on the following explicitly feasible resource region. Suppose the known box endpoints are rational data of description length $P_0$, with arithmetic bit length at most $B_0$, and the input primitive has the stated error bound. Put
\[
 m=\lfloor M/(3L)\rfloor-1,\qquad
 B=B_0+b+\lceil\log_2(N_{\rm raw}+M+L+2)\rceil+4.
\]
A sequential grid scan and schoolbook fixed-precision arithmetic have bounds
\begin{equation}\label{eq:robust-feasible}
 W\ge C_dB,\quad P\ge P_0+C_d\log_2(N_{\rm raw}+M+L+2)+C_d\log_2(b+2),
 \quad \tau\ge N_{\rm raw}+C_dN_{\rm raw}(m+1)B^2.
\end{equation}
The numerical input precision is $b$; arithmetic guard bits and exact rational data are charged in $W$ and $P$, not supplied as hidden retained state. The calibration input interface must admit the specified common packet. In \eqref{eq:robust-feasible} the time unit charges one raw call and one elementary bit operation; the selected terminal probe is included in $N_{\rm raw}$. The protocol clock is the fixed call index; arithmetic is padded to the stated uniform bound, and no data-dependent timing is available as a later observation. On this feasible region,
\begin{equation}\label{eq:robust-match}
 c_d\Phi\le\mathfrak R_{\mathcal C}^{\rm pw}(\mathbf b)\le C_d\Phi.
\end{equation}
The bounds are uniform in the positive widths, rank collapse, $n,s,L,M,\delta,b,\varepsilon$. Taking $L=R_0D_0'Q_0$ gives a matched law with growing mandatory simulator/controller/phase continuation memory. All five terms occur in the same family, with one common task and oracle, rather than in unformalized unrelated witnesses.
\end{theorem}
\begin{proof}
For the lower bound, put the uniform prior on the eight values of $\lambda$. A supremum risk is at least this Bayes average. Condition on the complete preterminal information and independent algorithm randomness, and apply conditional-mean projection separately to each actually executable terminal probe.

For the geometric block, unsuccessful acquisition leaves $X$ at its prepared law, so its excess over the oracle is at least $H_nV_a$. Conditioned on successful acquisition, its law is still the prepared box law. The tag is independent and must be recoverable exactly. Disjoint tag state supports give integers $k_c$ with sum at most $M$, even if the machine fuses bit, geometry, clock and controller information. The conditional distribution of $X$ given tag and success is uniform on the same box; all other reports can at most randomize its encoder. The argument of Theorem~\ref{thm:tag} and the anisotropic small-ball bound therefore give at least $c_d(1-H_n)\psi_a(\lfloor M/L\rfloor)$. Finite current-report precision only restricts this class further and cannot invalidate the lower bound.

Neither preterminal packet nor any other report depends on $\sigma$ or $\eta$. Their posterior means are consequently $1/2$, and their excesses over the same oracle are $\delta^2$ and $h_b^2=2^{-2b}/16$. On an erasure the bit posterior under the uniform prior is fair; on a nonerasure it is known. Its conditional variance is therefore $2\varepsilon/4=\varepsilon/2$. Multiply these blockwise bounds by the common factor $1/4$. The losses add because the block is chosen independently after encoding; no simultaneous counterfactual probe outcomes are assumed. This proves \eqref{eq:robust-lower}.

For the upper, retain the tag, one of the three bit states $0,1,\bot$, and either the unresolved geometric state or a grid label. Allocate $m$ geometric representatives as displayed. Then $3L(m+1)\le M$, and $m\ge M/(6L)$ for $M\ge12L$. For unresolved geometry predict the exact mean $1/2$; after acquisition round to a valid product grid. The box cover of Lemma~\ref{lem:geometry} gives squared error at most $C_d\psi_a(m)$. Holds do not round again. For the two ambiguous probabilities forecast $1/2$. For the bit use its observed value or $1/2$ on erasure. This rule has the same risk for every $\lambda$. Scaling \eqref{eq:psi} from $m$ to $\lfloor M/L\rfloor$ changes the constant by at most a dimension-independent numerical factor. The exact-label upper follows.

For the finite-bit version, obtain the successful report to coordinate accuracy $2^{-b}$ and approximate the selected readout to that accuracy. Compute candidate squared distances to absolute accuracy $2^{-2b}$ using a constant multiple of $B$ guard bits, and minimize these approximations. If $r$ is the closest distance to the quantized input, the selected distance is at most $\sqrt{r^2+C_d2^{-2b}}\le r+C_d2^{-b}$. Including the input error gives its covering radius plus $C_d2^{-b}$. Its squared error is at most twice the covering error squared plus $C_d2^{-2b}$. The unresolved mean and the three bit forecasts are dyadic and exact. The added acquired-report error is absorbed by the already present $2^{-2b}$ term, without an inverse width.

For completeness, choose grid cell counts by the finite dyadic-scale construction in Appendix~\ref{sec:algorithms}, rather than by comparing arbitrary real algebraic minima. This is a constant-factor cover with at most $m$ points. Iterate through its mixed-radix indices to generate each center on demand. Store a current index, the best index, the tag and bit state, and $O(d)$ current $B$-bit coordinates. Only the chosen index, tag and bit state survive the call. Rounded schoolbook multiplication, division and distance comparison take at most $C_dB^2$ bit operations per candidate, with guard bits ensuring the stated absolute error. Even scanning at every call costs no more than the time in \eqref{eq:robust-feasible}; in fact only the successful report needs a scan. The readonly program holds the rational data and binary descriptions of its budgets, not an observation-dependent table. This proves the resource bounds and the implementable upper. Zero widths are processed as exact point coordinates, never as divisors, so the constants survive their disappearance.
\end{proof}

\begin{proposition}[Exact erasure deficiency and the two elementary floors]\label{prop:floors}
In the declared parameter-independent simulator class, the deficiency of the bit-erasure experiment with erasure probability $\theta$ relative to the exact-bit experiment is exactly $\theta/2$. With a uniform bit preparation its best squared-prediction risk is $\theta/4$. For two identical preterminal calibration reports and terminal probabilities $1/2\pm h$, the minimum excess over the sign-knowing oracle under the equal prior is $h^2$. These statements use the same squared loss as Theorem~\ref{thm:robust}.
\end{proposition}
\begin{proof}
Write $P_0,P_1$ for the source report laws. Their total variation distance is $1-\theta$. Any parameter-independent simulator $K$ contracts total variation. If both simulated target laws are within $e$ of the exact target laws $\delta_0,\delta_1$, the triangle inequality gives
\[
 1=\|\delta_0-\delta_1\|_{\TV}
     \le2e+\|KP_0-KP_1\|_{\TV}\le2e+1-\theta.
\]
Thus $e\ge\theta/2$. Filling an erasure by an independent fair bit attains this bound for each parameter. Including the actual terminal bit in the comparison law changes neither conclusion: the marginal lower persists, and under either fixed parameter the fair-imputation error event has probability exactly $\theta/2$. The posterior variance under the uniform preparation is $1/4$ on an erasure and zero otherwise, proving the risk assertion. For the calibration pair the conditional mean is $1/2$ and conditional-mean projection gives the excess $\E(p-1/2)^2=h^2$, attained by that forecast.
\end{proof}

Theorem~\ref{thm:robust} is a full-ledger risk assertion on its stated feasible region, not a universal lower bound for the shortest program or fastest arithmetic in every kernel class. Its total-memory converse does not require constant overhead, but does require the audited continuation interface. The $b$-indexed information cell is compulsory in this experiment: optional inaccurate arithmetic would not imply its lower floor. These distinctions are part of the theorem's quantifiers, not qualifications inserted after changing the task.

\section{Approximate continuation distinctions and retained information}\label{sec:soft-state}
An exact continuation challenge forces disjoint state supports, as shown in Theorem~\ref{thm:tag}. Approximate challenges need a different argument: state supports may overlap, but their overlap consumes the information available for geometric resolution. The following converse applies to fused, randomized implementations; it does not assume a product register layout.

Let $C$ be uniform on $\{1,\ldots,L\}$, $L\ge2$, and let $U$ be independent of $C$ and uniform on the box $\prod_{i=1}^d[-a_i/2,a_i/2]$, with $a_1\ge\cdots\ge a_d>0$. The prepared experiment reveals $(C,U)$ through a paid input before a memory cut. Thereafter a legal continuation asks for $C$, or asks for a squared-prediction readout of $U$ using independently executable bounded probes. Affine rescaling embeds the coordinates in Bernoulli means in $[0,1]$; below the distortion $D$ is measured in the unscaled Euclidean coordinate metric. The known scaling converts it to the corresponding prediction excess. There is no further observation informative about $(C,U)$ between the cut and these readouts. Let $R$ be an independent shared random seed and $S$ the complete retained state, taking at most $M$ values conditional on $R$. The test selecting the continuation carries no information about $(C,U)$.

The entropy ingredients are classical \cite{cover}; their use here is tied to the explicitly retained-information cut. Use natural logarithms and let $h_2(t)=-t\log t-(1-t)\log(1-t)$. For $0\le\eta\le1-1/L$ write
\[
 f_L(\eta)=\log L-h_2(\eta)-\eta\log(L-1).
\]

\begin{theorem}[Noisy continuation--resolution converse]\label{thm:soft-state}
Suppose a continuation decoder recovers $C$ with error probability at most $\eta$. For any geometric decoder, even one additionally given the true $C$ for free, its mean squared distortion satisfies
\begin{equation}\label{eq:soft-converse}
 D\ge\max_{1\le l\le d}\frac{l}{2\pi e}
       \left(\prod_{i=1}^l a_i\right)^{2/l}
       \exp\left\{\frac{2}{l}\big(f_L(\eta)-\log M\big)\right\}.
\end{equation}
Necessarily $\log M\ge f_L(\eta)$. If a nominal workspace of $W$ bits is permitted to survive the cut, replace $\log M$ in both statements by $\log M+W\log2$. Workspace erased before the cut does not alter the bound. The result holds for every preprocessing algorithm and every program length or runtime compatible with the experiment.
\end{theorem}
\begin{proof}
We give the entropy argument to specify precisely the resource cut. Introduce the error indicator for the tag decoder. Conditional on that indicator and its decoded tag, there is at most one possible correct tag and at most $L-1$ possible incorrect tags. The chain rule and entropy maximization on a finite set give
\[
 H(C\mid S,R)\le h_2(\eta)+\eta\log(L-1),
\]
using monotonicity of the right side on $[0,1-1/L]$. Thus $I(C;S\mid R)\ge f_L(\eta)$. The state alphabet gives $H(S\mid R)\le\log M$, and the chain rule gives
\begin{equation}\label{eq:information-cut}
 I(U;S\mid C,R)\le\log M-f_L(\eta).
\end{equation}
These statements remain true when the encoding is stochastic.

For the first $l$ coordinates, the prior differential entropy conditional on $(C,R)$ is $\log\prod_{i\le l}a_i$. Replace a randomized geometric decoder by its conditional mean. Conditional translation by its output preserves differential entropy. Among distributions on $\R^l$ with mean square at most $d_0$, entropy is at most $(l/2)\log(2\pi e d_0/l)$: compare with the centered spherical Gaussian of variance $d_0/l$ using nonnegativity of relative entropy, and optimize that variance. Apply this conditional bound and then Jensen's inequality to obtain
\[
 I(U_1,\ldots,U_l;S\mid C,R)
 \ge \log\prod_{i\le l}a_i-\frac l2\log\frac{2\pi eD}{l}.
\]
One may first use positive variances and then take a limit when a conditional distribution is degenerate. Finite state entropy bounds the mutual information, so the extended-entropy version gives the same inequality. Data processing and \eqref{eq:information-cut} now give \eqref{eq:soft-converse} for each $l$. Nonnegativity of conditional mutual information gives the feasibility condition. With $W$ surviving bits the retained alphabet has cardinality at most $M2^W$, while erased words convey no information to a later decoder. This proves the last assertion.
\end{proof}

At $\eta=0$ the factor $\exp f_L(0)=L$ recovers the geometric scale associated with $M/L$ effective labels, up to constants. For nonzero error it supplies a continuous information cost, rather than declaring the tag either entirely free or entirely mandatory. The exact support-allocation theorem remains stronger when zero error is required. This converse is a retained-information bound, not an optimal-runtime theorem, and it does not establish a matching noisy-tag upper at every point of the resource vector. It cannot be used to label an arbitrary simulator register as necessary: the distinguishing continuation must be part of the actual experiment.

\section{A coupled acquisition--resolution law for one probe}\label{sec:single-probe}
The robust family in Section~\ref{sec:minimax} separates several lower-bound mechanisms by a common terminal menu. We now give a different experiment: one scalar hidden quantity and one Bernoulli probe. Acquisition, packet precision, retained memory, calibration ambiguity and report thinning act on that same quantity. The interaction through acquired mass is exact, not a sum of independent terminal blocks.

Fix $N\in\N_0$, $0\le s,\rho\le1$, $0<a\le1/4$, $0\le\delta\le1/8$, and an integer $b\ge0$. Put $L=2^b$. Prepare $U$ uniformly on $[-a/2,a/2]$. The unknown apparatus parameter is $\vartheta\in\{-\delta,\delta\}$; its calibration packet is identical at both parameter values. The terminal probe, available only after the forecast, is
\[
 T\mid U,\vartheta\sim\operatorname{Bernoulli}(1/2+U+\vartheta).
\]
Every algorithm is independent of the unknown sign. Calibration ambiguity here is compulsory information loss, not an optional numerical inaccuracy that an optimizer could remove.

Until acquisition, each raw call independently succeeds with probability $s(1-\rho)$. On success it reveals only the index of the equal-width cell of $U$ among $L$ cells, and the apparatus enters an absorbed state. On failure it emits $\perp$. After absorption every call emits $\perp$ and preserves the physical state. Thus $\perp$ never changes a retained label, and success overwrites it. The thinning parameter $\rho$ can be implemented by independently dropping a would-be successful report before absorption. One raw call and one time unit are charged at each attempt or hold. The forecast follows exactly $N$ such calls, with the horizon supplied by an exogenous clock. The terminal probe is a further paid call.

Set
\begin{equation}\label{eq:single-h}
 h=(1-s(1-\rho))^N,\quad
 b_*={1\over4}-{a^2\over12}-\delta^2,
 \quad \bar U=\hbox{the midpoint of the revealed cell}.
\end{equation}
The common ideal oracle knows $U$ and $\vartheta$, so its risk is $b_*$ for both parameter values. The full report history gives posterior mean $0$ on no acquisition and $\bar U$ on acquisition. Conditional on a success, cell indices are uniform and the unresolved within-cell variance is $a^2/(12L^2)$. Let $\mathcal A_M$ denote all algorithms retaining at most $M\ge1$ states, with no surviving report tape; initially no restriction is imposed on temporary computation.

\begin{theorem}[One-probe coupled minimax law]\label{thm:single-probe}
For the preceding two-member experiment family, checkpoint and online minimax excesses over the common oracle are equal. Uniformly over all the displayed parameters,
\begin{equation}\label{eq:single-law}
 \inf_{A\in\mathcal A_M}\sup_{\vartheta=\pm\delta}
       \{R_\vartheta(A)-b_*\}
 \asymp \delta^2+a^2\left[h+(1-h)(M^{-2}+2^{-2b})\right].
\end{equation}
The constants are numerical, independent of acquisition rarity and all resource parameters. An exact equality before comparison is
\begin{equation}\label{eq:single-exact}
 \delta^2+{ha^2\over12}+{(1-h)a^2\over12L^2}
       +q_M\left(h\delta_0+\frac{1-h}{L}\sum_{l=1}^L\delta_{u_l}\right),
\end{equation}
where $u_l=-a/2+(l-1/2)a/L$ and $q_M$ is ordinary squared distortion with $M$ centers. Equation~\eqref{eq:single-law}, not an assumed value of $q_M$, supplies the explicit resource scale.
\end{theorem}
\begin{proof}
Write $\widehat U$ for a forecast minus $1/2$. Conditional variance gives excess over $b_*$ equal to $\E(\widehat U-U-\vartheta)^2$. The preforecast law is independent of $\vartheta$. Averaging its two signs removes the linear cross term and yields
$\delta^2+\E(\widehat U-U)^2$, a lower for the worst sign. Conditional projection onto the full history separates the unresolved variances in \eqref{eq:single-exact} from quantization of the conditional mean. Encoder randomness, decoder randomness and independent shared randomness cannot lower the resulting distortion infimum. This proves the lower in \eqref{eq:single-exact}.

Conversely choose an optimal finite codebook for the finite measure in \eqref{eq:single-exact}. Such a codebook exists by compactness after projecting centers into $[-a/2,a/2]$. Choose the center of each nonempty encoding cell to be its conditional mean. The resulting prediction has $\E(\widehat U-U)=0$, so the two signs have the same excess. Implement it online by starting in the label for $0$, overwriting with the label for $u_l$ on a success of cell $l$, and copying the label on every $\perp$. The hardware absorbs after its first success, so no retained acquisition flag is needed. Even if the labels for $0$ and some $u_l$ coincide, the next $\perp$ has the same valid action, and any success overwrites. This proves equality of the checkpoint and online values and \eqref{eq:single-exact}.

For the scale lower, the unavoidable calibration and no-acquisition terms are explicit. On the acquired event, allowing the encoder to see $U$ exactly can only help. A uniform interval has $M$-point squared distortion $a^2/(12M^2)$: its Voronoi cells are intervals of lengths $d_i$ with total $a$, and their contributions are at least $d_i^3/(12a)$; convexity minimizes their sum at equal lengths. Empty cells do not help. The packet residual also gives $a^2/(12L^2)$. The maximum of these two acquired lower bounds is at least half their sum. This proves a constant multiple of the right side of \eqref{eq:single-law}.

For an upper bound when $M\ge2$, reserve the center $0$ and use $M-1$ equally spaced centers on the interval. A cell midpoint is approximated with squared error at most $a^2/[4(M-1)^2]$. Together with the exact variances in \eqref{eq:single-exact}, and $M-1\ge M/2$, this is the asserted upper. At $M=1$, the constant forecast $1/2$ has excess $\delta^2+a^2/12$ and the displayed scale is comparable to it. This includes $L=1$, $N=0$, and $s(1-\rho)=0$.
\end{proof}

\subsection{What the thinning parameter does and does not mean}
Here $\rho$ is a raw report-erasure probability, not a complete causal deficiency. Regard the thinned protocol as the source and the unthinned protocol with the same cell interface as the target. The identity report transformation uses one source call per target call, copies the declared clock and terminal probe, and retains no additional observation-dependent state. The action class consists only of the prescribed attempts followed by the terminal probe. Couple the two raw protocols with the same hidden $U$, the same unknown $\vartheta$, the same attempt coins until a disagreement, and the same terminal Bernoulli randomness. They agree until a would-be success is dropped. Therefore this particular resource-certified causal simulator has the uniform complete-law allowance
\begin{equation}\label{eq:thinning-path}
 \varepsilon_N\le1-(1-s\rho)^N\le Ns\rho.
\end{equation}
The inequality is an upper certificate, not an assertion of optimality. It remains valid when some calls already follow absorption, because they cannot create a new disagreement before the first dropped success. A simulator in the reverse direction may retain an already obtained packet to manufacture delayed successes, but that is a different resource account. We never replace the complete defect by the single-call coin parameter $\rho$.

For clarity one can compute an exact deficiency for the \emph{fixed-time packet reduction}, where success times are discarded. This is a different, explicitly less informative experiment. Let the parameter be a fixed cell $c\in\{1,\ldots,L\}$; an erasure channel returns $c$ with probability $w$ and $\perp$ otherwise. Compare $w_1=1-h$ to $w_0=1-(1-s)^N$, with $w_1\le w_0$.

\begin{proposition}[Exact terminal packet deficiency]\label{prop:terminal-deficiency}
For $L\ge2$, the deficiency of the $w_1$ packet experiment relative to the $w_0$ packet experiment, with total variation $\sup_A$, is
\begin{equation}\label{eq:packet-def}
 (w_0-w_1)(1-1/L).
\end{equation}
It is a terminal-experiment statement, not an equality for complete path laws, stopping times, or an infinite-horizon simulator.
\end{proposition}
\begin{proof}
Retain a revealed cell. On an erasure, output a uniformly random cell with probability $(w_0-w_1)/(1-w_1)$, and otherwise erase; the case $w_1=1$ has zero defect. The output erasure mass equals the target's, its mass at the true cell is short by $(w_0-w_1)(1-1/L)$, and precisely that mass lies at false cells. This proves the upper.

For the lower use the common decision problem of guessing the cell, with loss zero for a correct guess and one otherwise, under the uniform prior. The best success probability from a $w$ channel is $w+(1-w)/L$. Any parameter-independent randomization of the $w_1$ channel cannot improve this optimum. If every randomized output law were within $\varepsilon$ of its corresponding $w_0$ law, the target optimal decision rule would have averaged success at least $w_0+(1-w_0)/L-\varepsilon$ after that randomization. Comparing these two success bounds gives \eqref{eq:packet-def}. This argument covers all simulators, not only a particular imputation rule.
\end{proof}

Equation~\eqref{eq:single-law} exhibits a nonseparable effect: thinning changes the acquired mass $1-h$, which weights both resolution and packet-information errors, while increasing the unresolved acquisition variance. Calibration ambiguity affects the same Bernoulli mean. This is stronger than assembling these effects by an independently chosen terminal menu. It does not prove that every recurrent calibration problem has this law, nor that \eqref{eq:packet-def} can replace the complete causal defects in Theorem~\ref{thm:morphism}.

For rational $a$, the displayed grid upper has a direct finite-bit realization. The persistent register stores one of $M$ labels; a current $b$-bit cell index is scanned in temporary storage, mapped to a grid center, and erased before the next call. A conservative implementation uses $O(b+\log M+\operatorname{bits}(a)+b_{\rm out})$ temporary bits and a finite program for integer arithmetic, with polynomial bit-operation cost in these quantities per success and constant label-copy cost per hold. Rounded output contributes at most $2^{-b_{\rm out}}$ in the root-mean-square error. These are feasible upper accounts, not optimal lower bounds on workspace or runtime. The physical raw-call count is $N+1$ including the selected probe, regardless of acquisition success.

\section{Recurrent singular acquisition on accumulating components}\label{sec:moran}
We verify the preceding theorem from a raw kernel whose acquired measures need not have a quantization dimension. This verification uses no covariance, differentiable chart, or stationary environment resolvent. The conditional laws are explicit singular product measures; the component masses evolve by the raw refresh flow.

\subsection{A nonhomogeneous binary construction}
Let $(\rho_n)_{n\ge1}$ be any sequence in $[1/5,1/3]$, put $\ell_0=1$ and $\ell_n=\prod_{i=1}^n\rho_i$, and define
\begin{equation}\label{eq:moran-map}
 x(\omega)=\sum_{n\ge1}(1-\rho_n)\ell_{n-1}\omega_n,
 \qquad \omega\in\{0,1\}^{\N}.
\end{equation}
The series is uniformly convergent. Its image $C$ is compact and its digit representation is unique, because the two children at step $n$ are separated by the gap $(1-2\rho_n)\ell_{n-1}>0$. Let $\sigma$ be the image of independent fair digits. A cylinder of length $n$ has mass $2^{-n}$ and diameter $\ell_n$; the endpoints with constant tails belong to $C$.

For $k\ge1$ set
\begin{equation}\label{eq:moran-profile}
 r(k)=\ell_{\lfloor\log_2k\rfloor},\qquad r(0)=1.
\end{equation}
No assumption that $-n^{-1}\log\ell_n$ converges is imposed.

\begin{lemma}[Mass and cover at every binary scale]\label{lem:moran-profile}
There are at most $k$ points of $C$ covering $C$ with radius $r(k)$. For every real center $z$,
\[
 \sigma(B(z,r(k)/100))\le(2k)^{-1}.
\]
Moreover, $r(k-1)^2\le25r(k)^2$ for $k\ge2$. The digit shift
$F(x(\omega))=x(\omega_2,\omega_3,\ldots)$ preserves $\sigma$ and is $15$-Lipschitz on $C$.
\end{lemma}
\begin{proof}
Put $n=\lfloor\log_2k\rfloor$. Selecting one valid point from each length-$n$ cylinder uses $2^n\le k$ points and has covering radius at most $\ell_n$. Distinct length-$(n+2)$ cylinders have mutual distance at least
\[
 \min_{1\le i\le n+2}(1-2\rho_i)\ell_{i-1}
       \ge\ell_{n+1}/3\ge\ell_n/15.
\]
A ball of diameter $\ell_n/50$ therefore meets at most one such cylinder, whose mass is $2^{-(n+2)}\le(2k)^{-1}$. This holds for centers outside $C$ as well. The regularity bound changes only when $k$ is a power of two, when it follows from $\rho_n\ge1/5$.

Shift invariance follows because the shifted digits are again independent and fair, although the coding coefficients are nonhomogeneous. If two sequences first differ at digit $i\ge2$, their original distance is at least $(1-2\rho_i)\ell_{i-1}$, whereas their shifted distance is at most $\ell_{i-2}$. Their ratio is at most $1/[(1-2\rho_i)\rho_{i-1}]\le15$. If $i=1$, use shifted diameter at most one and original distance at least $1-2\rho_1\ge1/3$. This proves the global Lipschitz bound.
\end{proof}

\subsection{The prepared experiment}
Index components by $j\ge0$ and set $a_j=10^{-(j+2)}$. The hidden physical state is a prediction probability
\[
 p=\tfrac12+a_j(4+x),\qquad x\in C,
\]
or the limiting anchor $o=1/2$. Let $D_j=\tfrac12+a_j(4+C)$ and let $\sigma_j$ be the corresponding image of $\sigma$. The union of the $D_j$ and $o$ is compact in the ordinary absolute-value metric. The sets $D_j$ lie in disjoint intervals of length $a_j$ accumulating at $o$; their $a_j/2$ neighborhoods are disjoint. The task is the single executable terminal probe $T\sim\operatorname{Bernoulli}(p)$, with square loss. This probe is generated only after the prediction has been committed.

One paid initialization draws $j$ with known law $w_0$, draws $x$ with law $\sigma$, and reveals the resulting state. At continuation step $t$, a report type $H,E,R$ is chosen with probabilities $h_t,u_t,v_t$, independently of the current state. On $H$ the state is unchanged and only the type is reported. On $E$ the state becomes $1/2+a_j(4+F(x))$ and only the type is reported. On $R$ a new component is drawn with law $\pi_t$, a new independent point with law $\sigma_j$ is drawn, and the new state and type are reported. The anchor is fixed by $H$ and $E$; $R$ has the same state-independent law there. Every call, including a hold, costs one raw acquisition and one unit of physical time. These are nonnegative normalized kernels on compact standard-Borel spaces. A reveal can equivalently report its uniquely determined component and digit sequence; no hidden component is thereby supplied to the terminal decoder after the memory cut.

A finite protocol has $n$ continuation calls after the paid initialization, so its preforecast raw-call count is $N=n+1$. The terminal probe is a further paid call. Nonstationary probabilities are known exogenous protocol data. There are no optional control actions in this realization beyond continue and the declared terminal command; controller optimization is not smuggled into a fixed-protocol statement.

The observer's exact predictive state is $p$. Indeed the initial reveal and every later update determine $p$ from the full history. Equal $p$ values give the same future laws by the explicit kernel; different $p$ values are separated by the terminal Bernoulli probe. Thus the quotient is exactly the displayed compact set with the declared protocol interface. The shift is continuous on each component and at the anchor, because component diameters tend to zero. Reports $H,E$ have state-independent probabilities, and the entire reveal law under $R$ is state independent. Report total variation between any two preterminal states is consequently zero; terminal Bernoulli total variation equals $|p-p'|$. This supplies report continuity without falsely treating Dirac reveal laws as globally TV-continuous in an unrevealed state. The reveal is a freshly drawn state, not a noiseless observation of an arbitrary old state.

\begin{theorem}[Nonstationary recurrent singular realization]\label{thm:moran}
Fix a contraction sequence in $[1/5,1/3]$, any prepared law $w_0$ on the countable components, and any known nonstationary raw-flow protocol satisfying \eqref{eq:flow-balance} with $G=15$, some $\eta>0$, and $K>(1+\eta)225-1$. Let $w_n$ be given by \eqref{eq:raw-flow}, and define
\begin{equation}\label{eq:moran-Q}
 \mathcal Q_n(M)=\inf_{\sum_jk_j\le M}
       \sum_{j\ge0}w_{n,j}a_j^2
       \ell_{\lfloor\log_2(\max\{1,k_j\})\rfloor}^{\,2}.
\end{equation}
For every $n\ge0$ and $M\ge2$,
\begin{equation}\label{eq:moran-risk}
 R_\cp(n+1,M)-B_{n+1}\asymp
 R_\on(n+1,M)-B_{n+1}\asymp\mathcal Q_n(M).
\end{equation}
The constants depend on $K,\eta$ and the fixed scale-separation constants, but not on $n$, the number or rarity of acquired components, or the holding probabilities. The same statement applies to a finite subset of components with arbitrary positive scales obeying the displayed relative separation. No power-law or limiting dimension for $\sigma$ is required.
\end{theorem}
\begin{proof}
By Lemma~\ref{lem:moran-profile}, \textup{(MG)} holds with $r_j(k)=a_jr(k)$, $C_g=5$, $c_g=1/100$, $c_s=1/2$, and $D_r=25$; the factor five allows discarding to the anchor. The raw computation in Proposition~\ref{prop:recharge-flow} gives
$\mu_n=\sum_jw_{n,j}\sigma_j$. These are the acquired laws, not an auxiliary choice of measure on a reachable set.

Here is the local machine for an allocation. On an allocated component with $k_j>0$, retain its component index and a length-$\lfloor\log_2k_j\rfloor$ prefix, represented by the point with that prefix and all-zero tail. The entire pair is one of at most $\sum_j k_j$ labels. A reset into that component reads the prefix and overwrites the label. A reset into an unallocated component writes the common anchor label. A hold copies the label exactly. On $E$, an allocated label shifts its stored digit word and appends zero; the anchor remains the anchor. Thus the machine never needs to recover the index of an unallocated component.

On an allocated component, shifting the representative and rounding within the same component has error at most $15e_{t-1}+a_jr(k_j)$; indeed the shifted all-zero-tail representative is already a valid grid point, so the latter term is a harmless upper allowance. On an unallocated component, the distance to the anchor is at most $5a_j$. Under $E$, the ratio of the new distance to the old is at most $5/4$, so $15e_{t-1}$ still bounds it. On reset the old error vanishes and the new error is at most $5a_jr(k_j)$, including the zero-allocation case. Thus \eqref{eq:bridge-local} holds with $R_j=5a_jr(k_j)$ and $b_t\le1$. It is executable with the stated labels, not with an uncharged mode register. The initialized error measure has normalized bound one. Weighted continuation invariance follows from shift invariance and its gain bound. Proposition~\ref{prop:recharge-flow} now verifies \textup{(TB)} from the raw flow condition, and Theorem~\ref{thm:bridge} proves \eqref{eq:moran-risk}. The initialization/continuation indexing is exactly $N=n+1$.
\end{proof}

The class contains repeated active expansions. For example, with $\eta=1$, take $K=500$, $u_t=\theta_t/2001$, $v_t=2000\theta_t/2001$, $h_t=1-\theta_t$, and $\pi_t=(w_{t-1}+\omega_t)/2$. Then \eqref{eq:flow-balance} holds with room to spare, and arbitrary sequences $0\le\theta_t\le1$ are allowed. Whenever active expansions occur infinitely often they are recurrent, not a single initial transient. Arbitrarily long intervals with $\theta_t=0$ add raw time but no new rounding. Nonstationarity of $\omega_t$ can move acquired mass among old and newly introduced components. No constant divides by $\theta_t$ or by $w_{n,j}$.

\subsection{An oscillatory singular law and finite descriptions}
\begin{corollary}[A sharp causal profile without a limiting dimension]\label{cor:no-dimension}
There is a computably specified choice of $(\rho_n)$ for which the one-component squared quantization error $q_M(\sigma)$ satisfies
\[
 q_M(\sigma)\asymp\ell_{\lfloor\log_2M\rfloor}^{\,2},
\]
and $\log q_M(\sigma)/\log M$ has no limit. The checkpoint and online excess risks in the one-component version of Theorem~\ref{thm:moran} have the same nonconvergent profile, with uniform constants through all scales.
\end{corollary}
\begin{proof}
Use alternating blocks of $\rho_n=1/3$ and $\rho_n=1/5$, whose lengths are $2^{2^k}$ for block index $k\ge1$. Each new block length dominates the sum of all previous lengths. The cover and small-ball proof of Theorem~\ref{thm:bridge}, with a single component, gives the displayed estimate for every $M$ (the anchor is immaterial up to the one-label inequality). At $M=2^n$ the quotient of logarithms is
$2\sum_{i\le n}\log\rho_i/(n\log2)+o(1)$.
Along alternate block endpoints it tends to $-2\log3/\log2$ and $-2\log5/\log2$. Uniform multiplicative constants do not change these limits. Theorem~\ref{thm:moran} transfers the same profile to the online class.
\end{proof}

A finite-bit version has a concrete interpretation. Choose a finite active set of component indices and prefix lengths $d_j=\lfloor\log_2k_j\rfloor$. A reset packet supplies the component index and the first $d_j$ digits on those components; all other indices need only be mapped to the anchor by a declared input interface. This truncation is itself the observation kernel of the finite-input experiment, not free access to an infinite word. The representative states are integer labels. Shifts and holds are exact integer operations; a terminal probability is evaluated from a finite sum in \eqref{eq:moran-map}, followed by a declared output rounding. If $\rho_n$ and $a_j$ are rational with explicitly bounded bit length through the used depth, the operation counts and workspace are finite and computable from those lengths and the maximum depth. Surviving digits belong to the persistent label; temporary evaluation words are erased before the next call.

For a countable law, a certified tail bound $\sum_{j>J}w_{n,j}a_j^2$ and computable prefix weights permit finite enumeration of allocations up to any prescribed additive profile tolerance. A constant-factor allocation additionally needs a certified positive lower scale. The general theorem does not turn an arbitrary real sequence or an unbounded integer report into a finite-cost routine. Even without such effectiveness data, its Borel finite-label assertion remains precise. This separation is why program and input-interface assumptions are retained in the resource ledger.

\section{Initialization, exact encoders and executable approximations}\label{sec:algorithms}
This appendix specifies which operations the finite programs perform. A finite set of mathematical labels does not, by itself, certify finite-bit computation. Conversely, a physical program need not solve an exact-real nearest-neighbor tie in order to attain a covering bound.

\begin{lemma}[Paid initialization and continuation]\label{lem:initialization}
Suppose a raw initialization protocol uses $n_0$ calls, has a legal online encoder within the declared state capacity, and produces an observed record $H_0$, an exact predictive state in the valid continuation domain and an initial representative. Suppose the geometric, local update and transport hypotheses are verified under the \emph{joint actual initialization and continuation law}, including the initial error in \textup{(I)}. Then the proof of Theorem~\ref{thm:main} applies to $T$ subsequent calls, with $n_0+T$ preforecast calls and total raw budget $n_0+T+1$ when one terminal physical probe is executed and the initializer's workspace, program, clock and physical-time costs included. No conditioning on a successful initialization is permitted unless its probability and failure branches are retained.
\end{lemma}
\begin{proof}
The projection identity conditions on the complete record $(H_0,H_T)$ and does not require $H_0$ to be deterministic. The geometric lower integrates its actual terminal law. The online proof begins with the stated initial representative error and unrolls exactly the same recurrence for the $T$ subsequent updates. Concatenating the two finite machines is legal by the initializer hypothesis and the common state interface. Its raw calls and physical time add. This proves the assertion without postselection or an unpaid initial observation.
\end{proof}

\subsection{A width-uniform computable grid}
For a box with rational physical widths and a label allocation $k\ge1$, remove zero coordinates. If none remain, use its one point. Otherwise let $a_{\max}$ be its largest width and try scales $R_s=a_{\max}2^{-s}$, beginning with $s=0$. Set $n_i(s)=\max(1,\lceil a_i/R_s\rceil)$ and stop before the first scale with $\prod_i n_i(s)>k$. Multiplication may be capped at $k+1$. Output the midpoint grid of the last feasible scale. These are comparisons of specified rational data. There are at most $2+\lfloor\log_2 k\rfloor$ trials (including the rejected scale) because the first coordinate has $n_1(s)=2^s$.

If $h=\sqrt{\psi(k)}$, the proof of Lemma~\ref{lem:geometry} shows feasibility at scale $2h$. Some tried scale between $2h$ and $4h$ is therefore feasible, unless the initial one already has smaller width. The chosen grid has radius at most $2\sqrt{d_*}h$ and at most $k$ points. No inverse disappearing physical width enters this error constant. Exact comparisons may require more input bits when a nonzero width is extremely small; those bits belong to the rational description length, not to an allegedly uniform free oracle. A finite mixture may use exact algebraic allocation search, or any certified constant-factor allocation; these are separate clauses in the main theorem.

An approximate nearest point is selected as follows. Read the current physical report to norm accuracy $\eta$. Generate candidate centers on demand to accuracy $O(\eta^2)$ and compute their squared distances to absolute accuracy $O(\eta^2)$. Retain a least computed distance, breaking ties by index. If $r$ is the best true distance from the quantized report, the chosen distance is at most $\sqrt{r^2+C\eta^2}\le r+\sqrt C\eta$. Adding the input error gives covering radius plus $C'\eta$. This algorithm compares rational approximations and never has to certify inequality of arbitrary real oracles.

\subsection{Exact Borel transitions}
The following is pseudocode for mathematical label machines. The current real report is a transient Borel input, not a real number retained between calls.
\begin{verbatim}
FILTER(label, command a, current report y):
    p := valid_grid_point(label)
    r := q0 + (1 - 2*q0)*p
    p_new := r*(1 + a*y)/(1 + a*y*(2*r - 1))
    return first_nearest_grid_label(p_new)

ABSORBING_SENSOR(label, current report):
    if report is the hold/failure symbol: return label
    return first_nearest_valid_union_label(report)

FOLD_REFRESH(label, current report):
    if report is H: return label
    x := valid_grid_point(label)
    if report is E: x_new := l + 2*min(x-l, u-x)
    if report is (C,U): x_new := x/4 + 3*(l+a*U)/4
    return first_nearest_grid_label(x_new)
\end{verbatim}
The sensor starts with its unresolved label; its latent chart mark is never read. The recurrent encoder begins by encoding the paid initial reveal into the same grid. Before that reveal its protocol phase is fixed and requires only one known state. Lemma~\ref{lem:initialization} then applies with $n_0=1$; $N$ in Theorem~\ref{thm:recurrent} is the total $1+T$, not an uncharged continuation horizon.

\subsection{Finite-precision versions and local errors}
For a finite-bit machine replace the exact report by its certified numerical input, evaluate the update with interval or rounded rational arithmetic, project to the declared valid domain, and use the approximate nearest-point rule. A copied hold is an exact integer label copy. No arithmetic is done on a hold, so its innovation and local numerical error are zero. The following local error statements quantify the change from the exact Borel machine.

For the filter, work in a fixed safe calibrated parameter rectangle
\[
 q_0\in[\underline q,\overline q]\subset(0,1/2),\qquad
 a\in[a_-,a_+]\subset(0,1).
\]
On it the denominator in \eqref{eq:hmm-update} is at least $1-a_+$, and the invariant probabilities stay a positive distance from zero and one. The update is a rational function with all first derivatives in $(p,q_0,a,y)$ bounded on this compact domain. The derivative of logit is bounded on its image. By the mean-value theorem, parameter error at most $\delta$, report error at most $\eta$, and arithmetic absolute error at most $\eta$ give
\begin{equation}\label{eq:filter-bit-error}
 u_t\le C_{\underline q,\overline q,a_-,a_+}(\delta+\eta),
 \qquad v_N^{\rm num}\le C\eta.
\end{equation}
Round/clip approximate calibration into the safe rectangle. A common valid representative interval may be enlarged using the rectangle bounds; all derivative and density constants remain controlled by these displayed parameters. Consequently the propagated state error is at most $C(\delta+\eta)/(2\underline q)$. Choosing $\delta+\eta=O(m^{-1})$ recovers the $m^{-2}$ risk order with finite bits. Exact calibration and exact-real comparisons are not described as a finite-precision program.

For the singular sensor use physical target coordinates. A successful numerical report within $\eta$, an approximate nearest valid center and a readout within $\eta$ contribute at most $C\eta$ once. Let $\eta_0$ bound a precomputed unresolved mean. The root-mean-square readout allowance is
\begin{equation}\label{eq:sensor-bit-error}
 \sqrt{h_N}\eta_0+C\sqrt{1-h_N}\eta.
\end{equation}
A finite table or a supplied certified chart-evaluation and mean-computation routine must be included in $P$. Computability of a chart map without a complexity bound on its mean integral does not give a universal complexity bound. For affine rational boxes the mean is explicit. No approximation in \eqref{eq:sensor-bit-error} divides by a collapsing width or identifies the latent chart mark.

For the recurrent sensor the maps $\min$ and affine interpolation are Lipschitz in physical coordinates. Endpoint, report and arithmetic errors bounded by $\eta$ therefore give $u_t\le C\eta$ on the initial or active steps and $u_t=0$ on holds. The same forcing profile used for the grid error propagates this numerical error, yielding the last assertion of Theorem~\ref{thm:recurrent}. Calibration is either exact rational data or an input with the stated uniform absolute-error guarantee; no unknown endpoint can be read without it.

\subsection{Separate storage and time accounts}
With a procedurally generated grid and rational data, a one-dimensional filter or recurrent sensor needs a persistent index of $\lceil\log_2m\rceil$ bits and $O(B+\log m)$ transient bits for current reports and arithmetic, where $B$ includes calibration input length, numerical accuracy and guard bits. A scan costs at most $O(mB^2)$ bit operations per active step; a uniform one-dimensional grid also permits a direct rounded-index calculation. The description stores the rational constants and budget integers, not a table of $m$ arbitrary exact reals. For a generic supplied union chart table of $m$ centers in $d$ coordinates at $B$ bits, readonly data instead cost $O(mdB)$ and the distance scan has its own time cost. These are alternative implementations, not interchangeable accounts.

The audited-tag minimax program additionally retains the tag and one three-valued bit state. Its geometric unresolved state and $m$ grid states produce exactly the state-product upper used in Theorem~\ref{thm:robust}. The fixed protocol phase and horizon are supplied by the exogenous public clock; an observation-dependent clock would require additional persistent capacity. A terminal tag challenge emits the stored tag, without consulting a past report. Its $O(\log L)$ output cost is included in the conservative time bound there. Program bits, transient guard bits and the compulsory coarse numerical input are separate resources throughout.

\clearpage
\begin{thebibliography}{99}
\bibitem{blackwell} D. Blackwell, \emph{Equivalent comparisons of experiments}, Ann. Math. Statist. \textbf{24} (1953), 265--272. \href{https://doi.org/10.1214/aoms/1177729032}{doi:10.1214/aoms/1177729032}.
\bibitem{diaconis} P. Diaconis and D. Freedman, \emph{Iterated random functions}, SIAM Rev. \textbf{41} (1999), 45--76. \href{https://doi.org/10.1137/S0036144598338446}{doi:10.1137/S0036144598338446}.
\bibitem{feuer} A. Feuer and G. C. Goodwin, \emph{On-line quantization in nonlinear filtering}, J. Stat. Comput. Simul. \textbf{83} (2013), 1210--1222. \href{https://doi.org/10.1080/00949655.2012.656641}{doi:10.1080/00949655.2012.656641}.
\bibitem{graf} S. Graf and H. Luschgy, \emph{Foundations of Quantization for Probability Distributions}, Lecture Notes in Mathematics 1730, Springer, 2000. \href{https://doi.org/10.1007/BFb0103945}{doi:10.1007/BFb0103945}.
\bibitem{psr} M. L. Littman, R. S. Sutton, and S. Singh, \emph{Predictive representations of state}, Advances in Neural Information Processing Systems \textbf{14} (2001), 1555--1561.
\bibitem{pages} G. Pag\`es and H. Pham, \emph{Optimal quantization methods for nonlinear filtering with discrete-time observations}, Bernoulli \textbf{11} (2005), 893--932. \href{https://doi.org/10.3150/bj/1130077599}{doi:10.3150/bj/1130077599}.
\bibitem{sellami} A. Sellami, \emph{Quantization based filtering method using first order approximation}, SIAM J. Numer. Anal. \textbf{47} (2010), 4711--4734. \href{https://doi.org/10.1137/060652580}{doi:10.1137/060652580}.
\bibitem{zhu} S. Zhu, \emph{Asymptotic local uniformity of the quantization error for Ahlfors--David probability measures}, preprint, 2017, \href{https://arxiv.org/abs/1708.07657}{arXiv:1708.07657}.
\bibitem{shalizi} C. R. Shalizi and J. P. Crutchfield, \emph{Computational mechanics: pattern and prediction, structure and simplicity}, J. Stat. Phys. \textbf{104} (2001), 817--879. \href{https://doi.org/10.1023/A:1010388907793}{doi:10.1023/A:1010388907793}.
\bibitem{ferns} N. Ferns, P. Panangaden, and D. Precup, \emph{Metrics for finite Markov decision processes}, Proceedings of UAI 2004, 162--169. \href{https://arxiv.org/abs/1207.4114}{arXiv:1207.4114}.
\bibitem{saldi} N. Saldi, S. Y\"uksel, and T. Linder, \emph{Finite model approximations for partially observed Markov decision processes with discounted cost}, preprint, 2017. \href{https://arxiv.org/abs/1710.07009}{arXiv:1710.07009}.
\bibitem{costa} O. L. V. Costa, M. D. Fragoso, and R. P. Marques, \emph{Discrete-Time Markov Jump Linear Systems}, Springer, London, 2005. \href{https://doi.org/10.1007/b138575}{doi:10.1007/b138575}.
\bibitem{lecam} L. Le Cam, \emph{Asymptotic Methods in Statistical Decision Theory}, Springer, New York, 1986. \href{https://doi.org/10.1007/978-1-4612-4946-7}{doi:10.1007/978-1-4612-4946-7}.
\bibitem{torgersen} E. Torgersen, \emph{Comparison of Statistical Experiments}, Cambridge University Press, 1991. \href{https://doi.org/10.1017/CBO9780511666353}{doi:10.1017/CBO9780511666353}.
\bibitem{fritz} T. Fritz, \emph{A synthetic approach to Markov kernels, conditional independence and theorems on sufficient statistics}, Adv. Math. \textbf{370} (2020), 107239. \href{https://doi.org/10.1016/j.aim.2020.107239}{doi:10.1016/j.aim.2020.107239}.
\bibitem{ogura} M. Ogura and C. F. Martin, \emph{Stability analysis of positive semi-Markovian jump linear systems with state resets}, SIAM J. Control Optim. \textbf{52} (2014), 1809--1831. \href{https://doi.org/10.1137/130925177}{doi:10.1137/130925177}.
\bibitem{zhu-moran} S. Zhu, \emph{Asymptotic uniformity of the quantization error for Moran measures on $\mathbb R^1$}, preprint, 2018. \href{https://arxiv.org/abs/1802.03723}{arXiv:1802.03723}.
\bibitem{cover} T. M. Cover and J. A. Thomas, \emph{Elements of Information Theory}, second ed., Wiley, 2006. \href{https://doi.org/10.1002/047174882X}{doi:10.1002/047174882X}.
\bibitem{nonanticipative} C.~D. Charalambous, P.~A. Stavrou and N.~U. Ahmed, \emph{Nonanticipative rate distortion function and relations to filtering theory}, arXiv:1210.1266v3 (2013), related published version DOI:10.1109/TAC.2013.2293403.
\end{thebibliography}

\end{document}
